\documentclass[11pt,reqno]{amsart}

\usepackage{iftex}
\ifPDFTeX
  \usepackage[T1]{fontenc}
  \usepackage[utf8]{inputenc}
  \usepackage{lmodern}
\else
  \usepackage[no-math]{fontspec}
\fi
\usepackage{amsmath}
\usepackage{amssymb}
\usepackage{amsfonts}
\usepackage{mathrsfs}
\usepackage{amsthm}
\usepackage{mathtools}
\usepackage{geometry}
\usepackage{enumitem}
\usepackage{longtable}
\usepackage{booktabs}
\usepackage{array}
\usepackage[dvipsnames]{xcolor}
\usepackage[pagebackref=true]{hyperref}
\usepackage{bookmark}

\hypersetup{
  colorlinks=true,
  linkcolor=BrickRed,
  citecolor=OliveGreen,
  urlcolor=RoyalBlue,
  filecolor=RoyalBlue,
  linktoc=all,
  bookmarksnumbered=true,
  bookmarksopen=true,
  bookmarksopenlevel=1,
  pdfstartview=FitH,
  pdftitle={The area function of the Mertens excursion skeleton and horizon-localized criteria for the Riemann hypothesis},
  pdfauthor={V. Casa\~na-Giner},
  pdfsubject={Primary 11N37; Secondary 11M26, 60G70},
  pdfkeywords={Mertens function, Moebius function, excursions, Riemann hypothesis}
}
% Back-references in the bibliography: every entry lists, as links, the
% pages on which it is cited.
\renewcommand*{\backref}[1]{}
\renewcommand*{\backrefalt}[4]{%
  \ifcase #1 %
    \relax
  \or
    \ (Cited on p.~#2.)%
  \else
    \ (Cited on pp.~#2.)%
  \fi}

\geometry{margin=1in}
\numberwithin{equation}{section}
\setcounter{tocdepth}{3}
\setlength{\emergencystretch}{1.5em}

\theoremstyle{plain}
\newtheorem{theorem}{Theorem}[section]
\newtheorem{proposition}[theorem]{Proposition}
\newtheorem{lemma}[theorem]{Lemma}
\newtheorem{corollary}[theorem]{Corollary}

\theoremstyle{definition}
\newtheorem{definition}[theorem]{Definition}
\newtheorem{hypothesis}[theorem]{Hypothesis}
\newtheorem{heuristic}[theorem]{Heuristic}

\theoremstyle{remark}
\newtheorem{remark}[theorem]{Remark}

\DeclareMathOperator{\sgn}{sgn}
\DeclareMathOperator{\li}{li}
\newcommand{\Ms}{M_{\mathrm{s}}}
\newcommand{\Msk}{M_{\mathrm{sk}}}
\newcommand{\Ap}{\mathcal{A}_{+}}
\newcommand{\Am}{\mathcal{A}_{-}}
\newcommand{\Anet}{\mathcal{A}}
\newcommand{\eps}{\varepsilon}
\newcommand{\Pp}{P^{+}}
\newcommand{\Cc}{\mathcal{C}}
\newcommand{\Sq}{\mathcal{S}}
\newcommand{\En}{\mathcal{E}}
\newcommand{\mt}{\widetilde{m}}
\newcommand{\F}{\mathcal{F}}
\newcommand{\tagN}{\textup{[N]}}
\newcommand{\tagNE}{\textup{[N-el]}}
\newcommand{\tagC}{\textup{[C]}}
\newcommand{\tagL}{\textup{[L]}}

\title[The Mertens excursion skeleton]{The area function of the Mertens excursion skeleton\\
and horizon-localized criteria for the Riemann hypothesis}

\author{V.~Casa\~na-Giner}
\date{}

\subjclass[2020]{Primary 11N37; Secondary 11M26, 60G70}
\keywords{Mertens function, M\"obius function, excursions, Riemann hypothesis,
Kaczorowski--Pintz oscillation, lacunary sequences, Dyck paths, Catalan tails, extreme values}

\begin{document}

\begin{abstract}
We decompose the Mertens function $M(x)=\sum_{n\le x}\mu(n)$ into excursions
between zero runs and attach to each completed excursion a triangular
skeleton.  Geometric counting of disjoint excursions, combined with the
oscillation theorem of Kaczorowski--Pintz, yields horizon-localized forms of
Littlewood's criterion.  For every fixed
$\lambda>\lambda_*:=e^{\pi/\gamma_1}$, RH is equivalent to the usual
$M(t)\ll_\eps t^{1/2+\eps}$ bound restricted to
$x_0<t\le\min\{x_1,\lambda x_0\}$ on each excursion $(x_0,x_1]$;
$\lambda_*$ is the sharp threshold of this counting mechanism.

Using the logarithmic mean and late maximal-order theorem of Pintz, the
absolute area of the regularized skeleton has exponent
$1+\Theta/2$, where $\Theta=\sup\Re\rho$, and the Dirichlet series grouped by
completed excursions has absolute block-convergence abscissa $\Theta$.
At zeros $v$ of $M$ we obtain the exact weighted-correlation identity
$\sum_{n<v}M(n)^2=-\sum_{h<v}hC_h^{[v]}$ and localize its mass beyond every
lag scale $v^\alpha$ with $\alpha<\Theta$.  A fixed-mask Catalan description
separates excursion shape from orientation; divisor equations compress the
orientation fibre to polynomial size, while a fixed-shift quadratic form
cannot be uniformly observed from a sublinear number of divisor rows.  The
paper distinguishes results and formulations believed to be new, elementary
internal deductions, direct combinations of published theorems, and quoted
literature by the tags explained in Section~\ref{subsec:classification}.
\end{abstract}

\maketitle

\tableofcontents

\section{Introduction}
\label{sec:introduction}

This introduction has two parts. Sections~\ref{subsec:intro-mobius}--\ref{subsec:intro-walk}
collect, as a short manual, the facts about the Möbius function, the Mertens
function, $\zeta$, its zeros and the Chowla-type problems that are used or
referred to in the paper. Sections~\ref{subsec:classification}--\ref{subsec:organization}
describe what is proved here and classify it.

\subsection{The Möbius function and the Mertens function}
\label{subsec:intro-mobius}
The Möbius function, introduced by Möbius \cite{Mobius1832} in the
inversion formula that bears his name, is $\mu(1)=1$, $\mu(n)=(-1)^k$ if $n$
is a product of $k$ distinct primes, and $\mu(n)=0$ if $p^2\mid n$ for some
prime $p$. It is multiplicative and characterized by
\begin{equation}
\sum_{d\mid n}\mu(d)=\mathbf 1_{n=1}\qquad(\mu*1=\delta),
\end{equation}
which is the source of every recursion used below: for $n\ge2$,
$\mu(n)=-\sum_{d\mid n,\,d<n}\mu(d)$, and summing,
$\sum_{k\le x}M(x/k)=1$ for $x\ge1$ \cite[\S3.2]{MontgomeryVaughan2007}.
The Mertens function $M(x)=\sum_{n\le x}\mu(n)$ was tabulated by Mertens
\cite{Mertens1897}, who conjectured $|M(x)|\le\sqrt x$; Stieltjes had
claimed $M(x)=O(\sqrt x)$ in 1885 \cite{Stieltjes1905}. Squarefree integers
have density $6/\pi^2$: $\sum_{n\le x}\mu(n)^2=\frac6{\pi^2}x+O(\sqrt x)$.
Computationally, $M$ is evaluated by combinatorial algorithms based on the
recursion above \cite{DelegliseRivat1996,HT23}; Hurst \cite{Hurst18}
computed $M(x)$ for all $x\le10^{16}$. $M(n)$ is also the determinant of the
$n\times n$ Redheffer matrix \cite{Redheffer1977} and, up to sign, the
Euler characteristic of a cell complex constructed by Björner
\cite{Bjorner2011}.

\subsection{\texorpdfstring{Euler product, $\zeta$ and $1/\zeta$}{Euler product, zeta and 1/zeta}}
For $\Re s>1$,
\begin{equation}
\frac1{\zeta(s)}=\prod_p\Bigl(1-\frac1{p^s}\Bigr)=\sum_{n\ge1}\frac{\mu(n)}{n^s}
=s\int_1^\infty M(x)x^{-s-1}\,dx .
\end{equation}
The prime number theorem (Hadamard \cite{Hadamard1896}, de la Vallée
Poussin \cite{dlVP1896}) is equivalent to $M(x)=o(x)$ \cite{Landau1911}
and to $\sum_{n}\mu(n)/n=0$ \cite{vonMangoldt1897}; the Vinogradov--Korobov
zero-free region gives
$M(x)\ll x\exp\bigl(-c(\log x)^{3/5}(\log\log x)^{-1/5}\bigr)$
\cite{Walfisz1963,Ford2002}. We also use
$\sum_n\mu(n)\log n/n=-1$ (the value at $s=1$ of
$-\frac{d}{ds}\zeta(s)^{-1}=\sum_n\mu(n)(\log n)n^{-s}$, the convergence at $s=1$
being justified by the prime number theorem with error term).

\subsection{Zeros, the explicit formula and RH}
The nontrivial zeros $\rho=\beta+i\gamma$ of $\zeta$ lie in $0<\beta<1$,
infinitely many lie on $\beta=1/2$ (Hardy \cite{Hardy1914}), and the lowest
positive ordinate is $\gamma_1=14.134725\ldots$. Put
\begin{equation}
\Theta:=\sup\{\Re\rho:\ \zeta(\rho)=0,\ 0<\Re\rho<1\}.
\end{equation}
The functional equation makes the zero set symmetric under
$\rho\mapsto1-\overline\rho$, so $\Theta\ge1/2$, and $\Theta\le1$; RH is the
statement $\Theta=1/2$ \cite{Riemann1859}. It
has been verified numerically up to height $3\cdot10^{12}$
\cite{PlattTrudgian2021}. If all zeros are simple and RH holds, then for non-integer $x>0$,
along a sequence $T_k\to\infty$ on whose horizontal segments
$1/\zeta(\varsigma+iT_k)\ll T_k^{\eps}$ (such a sequence exists under RH,
\cite[\S14.16]{Titchmarsh1986}), \cite[\S14.27]{Titchmarsh1986}
\begin{equation}
M(x)=\lim_{k\to\infty}\sum_{|\gamma|<T_k}\frac{x^\rho}{\rho\zeta'(\rho)}-2+
\sum_{n\ge1}\frac{(-1)^{n-1}(2\pi/x)^{2n}}{(2n)!\,n\,\zeta(2n+1)} ;
\end{equation}
the last series collects the residues at the trivial zeros $-2n$. The
$\rho$-sum is taken in this symmetric truncated sense; its absolute
convergence is not known.
The size of $M$ is governed by $\Theta$: for every $\eps>0$,
\begin{equation}
M(x)\ll_\eps x^{\Theta+\eps},\qquad M(x)=\Omega(x^{\Theta-\eps}),
\end{equation}
the upper bound by Perron's formula and the bound $1/\zeta(s)\ll|t|^{\eps}$
to the right of $\Re s=\Theta$ (Littlewood's argument
\cite{Littlewood1912}, \cite[\S14.25]{Titchmarsh1986}, \cite[\S4]{Pintz2026}),
the lower bound because $\int_1^\infty M(x)x^{-s-1}dx$ would otherwise
continue $1/(s\zeta(s))$ analytically beyond $\Re s=\Theta$. In particular
Littlewood's criterion: RH holds iff $M(x)\ll_\eps x^{1/2+\eps}$ for every
$\eps>0$ \cite{Littlewood1912}. Under RH the best known bound is
$M(x)\ll\sqrt x\exp\bigl((\log x)^{1/2}(\log\log x)^{14}\bigr)$
(Soundararajan \cite{Soundararajan2009}, improving Maier--Montgomery
\cite{MaierMontgomery2009}); under RH and a hypothesis on
$\sum_{\gamma\le T}|\zeta'(\rho)|^{-2}$, Ng \cite{Ng2004} obtained a limiting
distribution for $e^{-y/2}M(e^y)$. If $|M(x)|\le C\sqrt x$ for all $x$, then
$|1/(s\zeta(s))|\le|\int_1^\infty M(x)x^{-s-1}dx|\le C/(\varsigma-\frac12)$
for $\varsigma>\frac12$; hence RH holds.  To read the multiplicity, approach a
zero $\frac12+i\gamma$ horizontally through $s=\varsigma+i\gamma$ with
$\varsigma\downarrow\frac12$.  Along this segment $s$ remains bounded and
bounded away from $0$, while
$|1/\zeta(\varsigma+i\gamma)|\ll_\gamma(\varsigma-\frac12)^{-1}$; a zero of order
$m\ge2$ would instead force growth of order
$(\varsigma-\frac12)^{-m}$.  Thus every zero on the critical line is simple
\cite[\S14.29]{Titchmarsh1986}.

\subsection{Oscillation}
The Mertens conjecture is false (Odlyzko--te Riele \cite{OdlyzkoTeRiele1985};
effective form by Pintz \cite{Pintz1987}); currently
$\limsup M(x)/\sqrt x>1.826$ and $\liminf M(x)/\sqrt x<-1.837$
\cite{Hurst18}, so in particular $M(x)=\Omega_\pm(\sqrt x)$ unconditionally
and $M$ changes sign, and vanishes, at arbitrarily large integers. Ingham
\cite{Ingham1942} showed that $\limsup M(x)/\sqrt x=+\infty$ follows from RH
and linear independence of the ordinates. Kaczorowski and Pintz \cite{KP86}
proved that the number of sign changes of $M$ in $[1,Y]$ is at least
$(\gamma_1/\pi-o(1))\log Y$, with oscillation of size $x^{\Theta-\eps}$, and
localized sign changes with large values of both signs in every range $[Y^{1-\eps},Y]$
(Theorem~\ref{thm:KP-quoted}). Pintz proved $D_M(Y)=\frac1Y\int_1^Y|M|\to\infty$
with $D_M(Y)\ge Y^{\beta_0}/(6|\gamma_0|^3)$ for every zero
\cite{Pintz1983}, and recently that the mean and the late maximal order of
$|M|$ are both $Y^{\Theta+o(1)}$, with a finer description through the
zeros \cite{Pintz2026} (Theorem~\ref{thm:pintz-2026}).

\subsection{Liouville's function}
\label{subsec:intro-liouville}
In this subsection only, $\lambda(n)$ and $L(x)$ denote Liouville's function
and its summatory function; everywhere else $\lambda$ is a horizon factor and
$L$ a block length. Liouville's function $\lambda(n)=(-1)^{\Omega(n)}$ has
$\sum\lambda(n)n^{-s}=\zeta(2s)/\zeta(s)$; hence $\lambda=\mathbf 1_{\square}*\mu$
and $\mu(n)=\sum_{d^2\mid n}\mu(d)\lambda(n/d^2)$, so with
$L(x)=\sum_{n\le x}\lambda(n)$,
\begin{equation}
L(x)=\sum_{d\le\sqrt x}M(x/d^2),\qquad M(x)=\sum_{d\le\sqrt x}\mu(d)L(x/d^2),
\end{equation}
and RH is equivalent to $L(x)\ll_\eps x^{1/2+\eps}$. Pólya's conjecture
$L(x)\le0$ is false \cite{Haselgrove1958}. Most correlation results below are
proved for $\lambda$ and $\mu$ by the same arguments; we state them for $\mu$.

\subsection{Short intervals, Chowla, Cesàro and logarithmic averages}
Chowla's conjecture \cite{Chowla1965} asserts, for distinct shifts
$h_1,\dots,h_k$, $\sum_{n\le x}\mu(n+h_1)\cdots\mu(n+h_k)=o(x)$ (Cesàro form);
the logarithmic form replaces $\sum_{n\le x}$ by $\sum_{n\le x}\frac1n$ and
$o(x)$ by $o(\log x)$, and is implied by the Cesàro form. Sarnak
\cite{Sarnak2010} observed that Chowla implies Möbius disjointness from
zero-entropy dynamics, and Tao proved the equivalence of the logarithmic
Chowla and logarithmic Sarnak conjectures \cite{Tao2017}. The known cases
are: the logarithmic two-point case (Tao \cite{Tao2016}), the logarithmic
odd-order cases (Tao--Teräväinen \cite{TaoTeravainen2018,TaoTeravainen2019}),
and Cesàro statements averaged over the shifts: Matomäki--Radziwiłł--Tao
\cite{MRT15} prove $\sum_{h\le K}|\sum_{n\le x}\mu(n)\mu(n+h)|=o(Kx)$ for
every $K=K(x)\to\infty$. The two-point Cesàro case for fixed $h$ remains
unproved. These rest on the Matomäki--Radziwiłł theorem \cite{MR16} on
multiplicative functions in almost all short intervals. For \emph{all}
short intervals, Matomäki--Teräväinen \cite{MatomakiTeravainen2023} prove
$\sum_{x<n\le x+x^\theta}\mu(n)\ll x^\theta(\log x)^{-1/3+\eps}$ for
$\theta>0.55$ and Fourier uniformity for $H\ge x^{3/5+\eps}$, extended by
Matomäki--Shao--Tao--Teräväinen \cite{MSTT23}; primes in all short intervals
are due to Huxley \cite{Huxley1972} and, with the Guth--Maynard large-value
estimate, to length $x^{17/30+o(1)}$ \cite{GuthMaynard2026}. The distinction
between results for all intervals and for almost all intervals is essential
here, because excursion endpoints are selected by $M$ itself.

\subsection{The walk viewpoint}
\label{subsec:intro-walk}
Denjoy \cite{Denjoy1931} observed that RH is the statement that $M$ grows
like a random walk; Wintner \cite{Wintner1944} introduced random
multiplicative functions, whose partial sums have better than square-root
cancellation (Harper \cite{Harper2020}). Areas and quadratic functionals of
lattice excursions are classical
\cite{Takacs91,NguyenThe04,Richard05,Janson07,DKW15,CP17}; here they are
applied to the deterministic Möbius walk.

\subsection{Classification of the results}
\label{subsec:classification}
Every numbered statement carries one tag:
\begin{itemize}[leftmargin=2.5em]
\item[\textup{[N]}] result or formulation believed to be new and requiring a
substantive argument; published inputs, when used, are named in the statement
or proof;
\item[\textup{[N-el]}] elementary exact result or short internal deduction
believed to be new in this form;
\item[\textup{[C]}] direct corollary, reformulation, or short combination of
published theorems with the excursion framework; no independent analytic
advance is claimed for the quoted input;
\item[\textup{[L]}] result from the literature, quoted with reference (or a
classical identity recorded for completeness).
\end{itemize}
Novelty tags record that we did not find the statement in the literature;
absence of a match is weak evidence.

\smallskip\noindent\emph{Results and formulations believed to be new.}
The excursion skeleton and its exact formulas; the geometric counting lemmas
for disjoint intervals; the horizon-localization theorem and the descending
horizons down to $\lambda x_0$, with
$\lambda_*=e^{\pi/\gamma_1}$ as the sharp threshold of this counting
mechanism; the exact exponent $1+\Theta/2$ of the skeleton area; the mean
logarithmic-length bound; wide peaks and packing; the prefix-dependency cutoff
$\chi_*=\lambda_*/2$; channel descent and its quantified persistence; the
failure of the relaxed divisor-sum class, of termwise contraction and of
bounded layer depth; the weighted correlation identity at zeros and its
$\Theta$-edge lag localization; the arithmetic selection results on the
Catalan reflection fibre; and the large-prime decomposition of $M$.

\smallskip\noindent\emph{Consequences and reformulations of published inputs.}
The return localization (RL) is a short consequence of the local
Kaczorowski--Pintz theorem; Littlewood's criterion is rewritten in $L^1$,
$L^2$, energy and correlation form; the exact exponents of $B_1,B_2,\En$
use Pintz's logarithmic mean and late maximal order; and averaged Chowla,
Fourier uniformity and short-interval cancellation are transported to
excursion blocks.  The excursion-block series of $1/\zeta$ combines these
inputs with the fixed-sign excursion decomposition.

\smallskip\noindent\emph{Literature.} The results of
Kaczorowski--Pintz, Pintz, Matomäki--Radziwiłł--Tao, Matomäki--Teräväinen,
Heath-Brown's identity, the reflection principle and the classical
identities for $\mu$.

\subsection{Main results}
\emph{Geometric counts} (Lemmas~\ref{lem:power-lacunary}--\ref{lem:linear-lacunary}).
For any ordered family of disjoint integer intervals $(x_0,x_1]$, the members
starting below $X$ with $x_1>x_0^A$ are $O_A(\log\log X)$, with $x_1>cx_0^2$
are $O_c(\log\log X)$, with $x_1>x_0\psi(x_0)$ are $o(\log X)$ when
$\psi\to\infty$, and with $x_1>\lambda x_0$ are at most
$\log X/\log\lambda+O(1)$.

\emph{Horizon localization} (Theorem~\ref{thm:horizon},
Corollaries~\ref{cor:descending}, \ref{cor:two-sided}). If $\Theta>1/2$, at
least $(\gamma_1/\pi-o(1))\log X$ distinct excursions starting below $X$
carry values $\ge\frac12t^{\Theta-\eta}$. Hence any horizon whose crossing
family has at most $(\kappa+o(1))\log X$ members, $\kappa<\gamma_1/\pi$, gives
an RH-equivalent localization; in particular
\begin{equation}
\mathrm{RH}\iff |M(t)|\ll_\eps t^{1/2+\eps}\ \text{ on }\
x_0<t\le\min\{x_1,\lambda x_0\},\qquad \lambda>\lambda_*=e^{\pi/\gamma_1}=1.2489\ldots
\end{equation}
The erosion chain $x_0^2\to x_0^2/2^j\to x_0\psi(x_0)\to\lambda x_0$ is
rigorous for general families; for $M$ itself the power horizons are
inactive by Corollary~\ref{cor:KP-local-returns}.

\emph{Exact exponents} (Corollaries~\ref{cor:pintz-excursion-exponents},
\ref{cor:area-exponent}, Theorem~\ref{thm:excursion-block-abscissa}).
\begin{gather}
\sum_{n\le X}|M(n)|=X^{1+\Theta+o(1)},\qquad \sum_{n\le X}M(n)^2=X^{1+2\Theta+o(1)},\\
\En(X)=X^{1+\Theta+o(1)},\qquad \int_0^X|\Ms|=X^{1+\Theta/2+o(1)},
\end{gather}
and the series $\sum_I D_I(s)$ over completed excursions has abscissa of
absolute block convergence $\Theta$ and equals $1/\zeta(s)$ for
$\Re s>\Theta$.

\emph{Linear window} (Section~\ref{sec:prefix}). For $u<t\le\lambda u<2u$ and
$M(u)=0$,
\begin{equation}
M(t)=-\sum_{d\le\lambda u/2}\mu(d)\bigl(\lfloor t/d\rfloor-\lfloor u/d\rfloor\bigr),
\end{equation}
so the RH-equivalent window is determined by the prefix up to $\chi u$ for
every $\chi>\chi_*=\lambda_*/2$; the channel decomposition, the kernel form
and the multiplicative-layer filtration make the dependency explicit and
identify where cancellation must occur.

\emph{Correlations} (Section~\ref{sec:chowla}). At zeros $v$ of $M$,
$\sum_{n<v}M(n)^2=-\sum_{h<v}hC_h^{[v]}=v^{1+2\Theta+o(1)}$, lags
$h\le\sqrt v$ are negligible at the $v^2$ scale, and lags $h\le v^\alpha$,
$\alpha<\Theta$, carry a vanishing proportion.

\emph{Extreme values} (Section~\ref{sec:extreme}). Under a uniform tail
hypothesis on a normalized excursion statistic in a stochastic model, each
level of lacunarity removes one logarithm from the extreme scale.

\subsection{Organization}
\label{subsec:organization}
Sections~\ref{sec:definitions}--\ref{sec:internal_corr}: excursion geometry.
Sections~\ref{sec:zeta}--\ref{sec:criteria}: the relation with $\zeta$, area
bounds, global criteria, exact exponents and the block series of $1/\zeta$.
Section~\ref{sec:lacunary}: counts, Kaczorowski--Pintz input, horizon
localization, Catalan reference laws, wide peaks.
Section~\ref{sec:prefix}: the localized window and the large-prime
decomposition. Section~\ref{sec:chowla}: correlations.
Section~\ref{sec:extreme}: stochastic models. Then a one-frequency model,
numerical data, and conclusions.

\subsection*{Conventions}
$\rho=\beta+i\gamma$ runs over nontrivial zeros;
$\Theta=\sup\{\Re\rho:\zeta(\rho)=0,\ 0<\Re\rho<1\}\in[1/2,1]$;
$\gamma_1=14.134725\ldots$; RH means $\Theta=1/2$. $M(x)=M(\lfloor x\rfloor)$
for real $x$, $M(x)=0$ for $x<1$. Heuristics are marked as such.

\section{Notation and variables}
\label{sec:notation}

Global symbols keep the meanings below. Symbols marked ``local'' are
declared again where used and never leave that statement or subsection.
Constants $c,C$ are positive and may change from line to line.

\small
\setlength{\LTpre}{4pt}
\setlength{\LTpost}{8pt}
\begin{longtable}{@{}p{0.24\textwidth}p{0.47\textwidth}p{0.23\textwidth}@{}}
\toprule
\textbf{Symbol} & \textbf{Meaning} & \textbf{Scope}\\
\midrule
\endfirsthead
\toprule
\textbf{Symbol} & \textbf{Meaning} & \textbf{Scope}\\
\midrule
\endhead
$\rho=\beta+i\gamma$ & Nontrivial zero of $\zeta$; $\beta$ is used only here & Global\\
$\Theta$ & $\sup\{\Re\rho\}$ & Global\\
$\gamma_1$ & Lowest positive ordinate, $14.134725\ldots$ & Global\\
$\lambda_*,\ \chi_*$ & $e^{\pi/\gamma_1}$, $\lambda_*/2$ & Global\\
$\lambda$ & Horizon factor, $\lambda>1$ (Liouville's $\lambda(n)$ only in Section~\ref{subsec:intro-liouville}) & Global\\
$\psi$ & Function with $\psi(x)\to\infty$ in the superlinear horizon $x\psi(x)$ & Section~\ref{sec:lacunary}\\
$M,\ M_0$ & Mertens function; $M_0=M+2$ (normalization of \cite{KP86}) & Global\\
$I_\nu=(x_{0,\nu},x_{1,\nu}]$ & Completed excursion; also $(u_\nu,v_\nu]$ & Global\\
$\sigma_\nu$ & Orientation $\pm1$ of $I_\nu$ & Global\\
$\ell_\nu,\ H_\nu,\ \tau_\nu$ & Length, height, first apex & Global\\
$a_\nu,\ b_\nu$ & $\tau_\nu-x_{0,\nu}$, $x_{1,\nu}-\tau_\nu$ & Global\\
$\mathcal Z=[p,q]$ & Zero run & Section~\ref{sec:formulas}\\
$u(x),\ L_x$ & Last zero $\le x$; terminal length $\lfloor x\rfloor-u(x)$ & Global\\
$\Ms,\ \Msk$ & Regularized / unregularized skeleton & Global\\
$\Ap,\ \Am,\ \Anet,\ \mathcal A_{\rm abs}$ & Positive, signed negative, net, absolute area & Global\\
$\widetilde A$ & $\int_0^xM$ & Section~\ref{sec:zeta}\\
$\mathrm{ch}_\nu$ & Chord of $\Anet$ over $I_\nu$ & Theorem~\ref{thm:energy}\\
$\mathfrak q_\nu=q_\nu$ & Number of active steps $\#\{n\in I_\nu:\mu(n)\ne0\}$; $q$ in a general block & Global\\
$Q_\nu,\ \mathfrak b_\nu$ & $\sum_{n\in I_\nu}M(n)^2$, $\sum_{n\in I_\nu}|M(n)|$ & Global\\
$Q_{\rm term},\ \mathfrak b_{\rm term}$ & Same for the terminal open prefix & Global\\
$B_1,\ B_2,\ \En$ & $\sum_{n\le X}|M|$, $\sum_{n\le X}M^2$, chord energy & Global\\
$C_h^{(\nu)},\ C_h^{[v]},\ C_h^{(u,X]}$ & Lag-$h$ correlation in $I_\nu$, in $(0,v]$, in $(u,X]$ & Global\\
$C_g^{\langle h\rangle}$ & Lag-$g$ correlation in a truncated prefix of length $h$ & Local\\
$S_t$ & Oriented partial sum inside a declared block & Local\\
$\F,\ E_\F(X)$ & Family of excursions and its counting function & Section~\ref{sec:lacunary}\\
$\mathfrak h,\ \F_{\mathfrak h}$ & Horizon and crossing family & Section~\ref{sec:lacunary}\\
$A$ & Power exponent in $x_1>x_0^A$ & Section~\ref{sec:lacunary}\\
$\eta$ & Loss in the Kaczorowski--Pintz exponent $\Theta-\eta$ & Global\\
$\kappa$ & Coefficient of a crossing count $(\kappa+o(1))\log X$ & Theorem~\ref{thm:horizon}\\
$\Sq,\ L^\#,\ r_I$ & Squarefree clock, compressed length, support time & Catalan subsection\\
$L$ & Length of a declared block (also support length $L^\#$ in Catalan lemmas) & Local\\
$p_{\rm Cat},\ p^{\max}_{\rm Cat},\ p_{\rm open}$ & Reference-law tail weights & Catalan subsection\\
$J_N,\ E_\nu$ & Number of completed excursions up to a zero cutoff $N$; active support of $I_\nu$ & Arithmetic selection\\
$\mathbf z,\ \boldsymbol\sigma$ & $0/1$ reflection vector; $\pm1$ orientation vector & Arithmetic selection\\
$\mathfrak D_N,\ G_N$ & Divisor-defect functional and its orientation Gram matrix & Arithmetic selection\\
$m_\nu,\ \phi_\nu,\ \mathscr C^{\rm div}$ & Support centre, positive oriented increments, central divisor matrix & Arithmetic selection\\
$\mathbf D_N,\ \mathsf A_h,\ \mathfrak S_h$ & Divisor matrix; fixed-shift form; pair density & Arithmetic selection\\
$\theta_{\rm w},\ \theta'$ & Short-interval exponents $>0.55$ & Wide peaks\\
$\Lambda,\ h_u,\ \omega_u$ & Horizon factor, block length, Lipschitz modulus & Wide peaks\\
$G$ & $(H-h_u)_+$ & Proposition~\ref{prop:wide-peaks}\\
$\mu_u$ & $\mu\mathbf 1_{[1,u]}$ & Section~\ref{sec:prefix}\\
$\Delta_k,\ \mathcal R_K$ & Channel increment, deep remainder & Section~\ref{sec:prefix}\\
$\vartheta$ & Normalizing exponent in $|M(y)|/y^\vartheta$ & Section~\ref{sec:prefix}\\
$c_y(t,u),\ K_y$ & Kernel coefficients (argument $y$ of $M(y)$) & Section~\ref{sec:prefix}\\
$\Upsilon,\ \Cc$ & Multiplicity-weighted and coherent coefficients & Section~\ref{sec:prefix}\\
$\mathbf a=(a_d),\ D_{\mathbf a}$ & Relaxed prefix vector and its divisor sum & Section~\ref{sec:prefix}\\
$P$ & Prefix cutoff in layer statements & Section~\ref{sec:prefix}\\
$\Pp(n)$ & Largest prime factor & Section~\ref{sec:prefix}\\
$U,\ W,\ M^{\rm sm}_U,\ M^{\rm lp}_U,\ \mt$ & Large-prime level, $t/U^2$, smooth and large-prime parts, $\sum_{a\le y}\mu(a)/a$ & Section~\ref{sec:prefix}\\
$\Omega(n),\ \omega(n)$ & Prime factors with / without multiplicity & Global\\
$D_I(s)$ & Dirichlet block of a completed excursion & Section~\ref{sec:criteria}\\
$s=\varsigma+i\tau$ & Complex variable & Global\\
$e(x),\ \xi$ & $e^{2\pi ix}$; frequency & Section~\ref{sec:chowla}\\
$K$ & Lag cutoff in averaged-Chowla statements & Section~\ref{sec:chowla}\\
$R,\ F_R,\ V_R$ & Fejér scale and sums & Section~\ref{sec:chowla}\\
$\iota$ & Exponent gap in $L\ge u^{3/5+\iota}$ & Section~\ref{sec:chowla}\\
$(\mathfrak X,\mathcal G,\Pr),\ Z_k$ & Probability space and excursion statistic & Section~\ref{sec:extreme}\\
$\mathcal N(X)$ & Deterministic count bound & Section~\ref{sec:extreme}\\
$\kappa_0,\ \varpi$ & Tail constant and exponent & Section~\ref{sec:extreme}\\
\bottomrule
\end{longtable}
\normalsize

\section{Excursions and the regularized skeleton}
\label{sec:definitions}

Let $\mu$ be the M\"obius function and
\begin{equation}
\label{eq:mertens_def}
M(x)=\sum_{n\le x}\mu(n),\qquad x\ge0,
\end{equation}
piecewise constant and right-continuous, with jumps $\mu(n)\in\{-1,0,1\}$.
Throughout the paper, for real $x\ge0$ we use the convention
\begin{equation}
M(x)=M(\lfloor x\rfloor),
\end{equation}
and we set $M(x)=0$ for $x<1$. Thus expressions such as $M(t/k)$ always
mean $M(\lfloor t/k\rfloor)$. The segmentation of $M$ by zero runs and the record of its
extrema are standard computational objects; Hurst~\cite{Hurst18} computed
$M(x)$ for all $x\le10^{16}$. The area and quadratic functionals of lattice
excursions are classical
\cite{Takacs91,NguyenThe04,Richard05,Janson07,DKW15,CP17}; here they are
applied to the deterministic M\"obius walk.

\begin{definition}[Excursions and zero runs]
\label{def:excursion}
An \emph{excursion} is an interval
\begin{equation}
I_\nu=(x_{0,\nu},x_{1,\nu}]
\end{equation}
with integer endpoints, $M(x_{0,\nu})=M(x_{1,\nu})=0$,
$M(x_{0,\nu}+1)=\sigma_\nu\in\{\pm1\}$, and
$\sigma_\nu M(n)\ge1$ for every integer $x_{0,\nu}<n<x_{1,\nu}$. Thus
$x_{1,\nu}$ is the first return to zero after the departure from
$x_{0,\nu}$. Consecutive excursions may be separated by a \emph{zero run}
$[x_{1,\nu},x_{0,\nu+1}]$, a maximal set of consecutive integers on which
$M=0$; it may consist of a single point. The first excursion is $(0,2]$.
Since $M$ assumes positive and negative values at arbitrarily large
arguments ($M(x)=\Omega_\pm(\sqrt x)$, Remark~\ref{rem:omega}) and its
integer jumps belong to $\{-1,0,1\}$, between two integers where $M$ has
opposite signs it must vanish at an integer.  Hence $M$ has arbitrarily large zeros and every
excursion is finite. We also write $u_\nu=x_{0,\nu}$ and
$v_\nu=x_{1,\nu}$.

Put
\begin{equation}
\ell_\nu=x_{1,\nu}-x_{0,\nu},\qquad
H_\nu=\max_{x_{0,\nu}<n<x_{1,\nu}}|M(n)|,
\end{equation}
\begin{equation}
\tau_\nu=\min\{n\in I_\nu:|M(n)|=H_\nu\},\qquad
a_\nu=\tau_\nu-x_{0,\nu},\qquad b_\nu=x_{1,\nu}-\tau_\nu.
\end{equation}
Every excursion has $H_\nu\ge1$ and $\ell_\nu\ge2$.
\end{definition}

\begin{definition}[Completed excursions and open prefixes]
\label{def:completed-open}
An excursion in the sense of Definition~\ref{def:excursion} is always
\emph{completed}: both $M(x_{0,\nu})=0$ and $M(x_{1,\nu})=0$ are part of
the definition. A segment $(u,x]$ with $M(u)=0$ and $M(n)\ne0$ for
$u<n\le x$, with no return condition at $x$, is an \emph{open prefix}; the
terminal block of Definition~\ref{def:Ms-terminal} is one. Any estimate
whose proof uses the backward bound $|M(u+r)|\le x_{1,\nu}-u-r$ holds for
completed excursions only and is never applied to an open prefix.
\end{definition}

\begin{lemma}[Height and step count]
\label{lem:height_steps}
\tagNE\ For every excursion,
\begin{equation}
\label{eq:height_steps}
H_\nu\le\min(a_\nu,b_\nu)\le\frac{\ell_\nu}{2}.
\end{equation}
\end{lemma}

\begin{proof}
The steps of $M$ are $0$ or $\pm1$, so at least $H_\nu$ steps are needed to
reach height $H_\nu$ from zero and at least $H_\nu$ more to return to zero.
\end{proof}

\begin{definition}[Crossing and touching zero runs]
\label{def:crossing_touching}
The zero run $[x_{1,\nu},x_{0,\nu+1}]$ between consecutive excursions is a
\emph{crossing run} if $\sigma_{\nu+1}=-\sigma_\nu$ and a \emph{touching run}
if $\sigma_{\nu+1}=\sigma_\nu$.
\end{definition}

\begin{definition}[Regularized skeleton]
\label{def:Ms}
The \emph{skeleton} $\Ms$ is the continuous, piecewise linear function
defined on every excursion by
\begin{equation}
\label{eq:Ms_def}
\Ms(t)=
\begin{cases}
\sigma_\nu\sqrt{H_\nu}\dfrac{t-x_{0,\nu}}{a_\nu},
 &x_{0,\nu}\le t\le\tau_\nu,\\[2ex]
\sigma_\nu\sqrt{H_\nu}\dfrac{x_{1,\nu}-t}{b_\nu},
 &\tau_\nu\le t\le x_{1,\nu},
\end{cases}
\end{equation}
and by $\Ms\equiv0$ on zero runs.  Here ``regularized'' refers only to
the normalization of the apex from $H_\nu$ to $H_\nu^{1/2}$; no analytic
smoothing operation is intended.  The \emph{unregularized skeleton} $\Msk$
is defined in the same way with peak $\sigma_\nu H_\nu$.
\end{definition}

\begin{definition}[Terminal block of a truncated computation]
\label{def:Ms-terminal}
Let $N\ge2$ and $u=\max\{n\le N:M(n)=0\}$. If $u=N$ there is no terminal
block. Otherwise $M(n)\ne0$ for $u<n\le N$, and since $M$ has steps
$0,\pm1$ all these values have the sign of $M(N)$. Put
\begin{equation}
\sigma_{\rm term}=\sgn M(N),\qquad
H_{\rm term}=\max_{u<n\le N}|M(n)|,\qquad
\tau_{\rm term}=\min\{n\in(u,N]:|M(n)|=H_{\rm term}\}.
\end{equation}
On $[u,N]$ the regularized skeleton is the segment from $(u,0)$ to
$(\tau_{\rm term},\sigma_{\rm term}\sqrt{H_{\rm term}})$ followed, if
$\tau_{\rm term}<N$, by the segment ending at
$(N,\sigma_{\rm term}\sqrt{|M(N)|})$. Its signed area is
\begin{equation}
\Delta\Anet_{\rm term}(u,N]=\frac{\sigma_{\rm term}}2
\Bigl[\sqrt{H_{\rm term}}(N-u)+
\sqrt{|M(N)|}(N-\tau_{\rm term})\Bigr].
\end{equation}
The terminal skeleton is provisional: after the next return to zero the
completed excursion acquires a different descending side. Independently of
it, for real $x$ with last zero $u(x)\le\lfloor x\rfloor$ define
\begin{equation}
\mathfrak b_{\rm term}(x):=\sum_{u(x)<n\le\lfloor x\rfloor}|M(n)|,\qquad
Q_{\rm term}(x):=\sum_{u(x)<n\le\lfloor x\rfloor}M(n)^2,\qquad
L_x:=\lfloor x\rfloor-u(x).
\end{equation}
These depend only on $M$.
\end{definition}

\begin{definition}[The three curves]
\label{def:areas}
With $\Ms^{\pm}=\max(\pm\Ms,0)$,
\begin{equation}
\label{eq:areas_def}
\Ap(x)=\int_0^x\Ms^+(t)\,dt,\qquad
\Am(x)=-\int_0^x\Ms^-(t)\,dt,\qquad
\Anet(x)=\int_0^x\Ms(t)\,dt.
\end{equation}
$\Ap$ is nondecreasing, $\Am$ is nonincreasing, and both are constant on
zero runs and on excursions of the opposite sign.  The curve $\Am$ records
negative area with its sign; it is not the absolute value of the negative
area.
\end{definition}

\begin{remark}[The family of $g$-skeletons]
\label{rem:g-skeleton}
For any $g:\mathbb N_0\to[0,\infty)$ with $g(0)=0$ and $g(h)>0$ for $h\ge1$,
one may replace the apex height $\sqrt H$ by $g(H)$. The signed area
increment becomes $\tfrac12\sigma_\nu g(H_\nu)\ell_\nu$. The correlation
identities of Section~\ref{sec:internal_corr} concern the M\"obius
increments and do not depend on $g$. The choice $g(H)=\sqrt H$ is a
normalization.
\end{remark}


\section{Exact formulas}
\label{sec:formulas}

\begin{lemma}[Increment per excursion]
\label{lem:jump}
\tagNE\ For every complete excursion,
\begin{equation}
\label{eq:jump}
\Anet(x_{1,\nu})-\Anet(x_{0,\nu})
=\sigma_\nu\frac{\sqrt{H_\nu}\,\ell_\nu}{2},
\qquad
\sum_{n\in I_\nu}\Ms(n)
=\sigma_\nu\sqrt{H_\nu}\frac{\ell_\nu}{2}.
\end{equation}
\end{lemma}

\begin{proof}
The two pieces are triangles of height $\sqrt{H_\nu}$ and bases $a_\nu,b_\nu$.
For the discrete sum, the ascending values at $j=1,\ldots,a_\nu$ sum to
$\sqrt{H_\nu}\sum_{j\le a_\nu}j/a_\nu=\sqrt{H_\nu}(a_\nu+1)/2$; the
descending values at $\tau_\nu+r$, $r=1,\ldots,b_\nu$, sum to
$\sqrt{H_\nu}\sum_{r\le b_\nu}(b_\nu-r)/b_\nu=\sqrt{H_\nu}(b_\nu-1)/2$.
\end{proof}

\begin{proposition}[Explicit form of the three curves]
\label{prop:explicit}
\tagNE\ Let $x\in I_\nu$. Then
\begin{equation}
\label{eq:explicit_A}
\Anet(x)=
\begin{cases}
\Anet(x_{0,\nu})+\sigma_\nu\sqrt{H_\nu}
\dfrac{(x-x_{0,\nu})^2}{2a_\nu},&x\le\tau_\nu,\\[2ex]
\Anet(x_{0,\nu})+\sigma_\nu\sqrt{H_\nu}
\left(\dfrac{a_\nu}{2}+(x-\tau_\nu)-
\dfrac{(x-\tau_\nu)^2}{2b_\nu}\right),&x\ge\tau_\nu.
\end{cases}
\end{equation}
At excursion endpoints,
\begin{equation}
\label{eq:explicit_net}
\Anet(x_{1,k})=\sum_{\nu\le k}\sigma_\nu\frac{\sqrt{H_\nu}\ell_\nu}{2},\qquad
\Ap(x_{1,k})=\sum_{\substack{\nu\le k\\\sigma_\nu=+1}}
\frac{\sqrt{H_\nu}\ell_\nu}{2},\qquad
\Am(x_{1,k})=-\sum_{\substack{\nu\le k\\\sigma_\nu=-1}}
\frac{\sqrt{H_\nu}\ell_\nu}{2}.
\end{equation}
The curves are piecewise quadratic and $C^1$.
\end{proposition}

\begin{proof}
Integrate \eqref{eq:Ms_def}; the endpoint formulas follow from
Lemma~\ref{lem:jump}.
\end{proof}

A point $y$ is a \emph{strict inflection point} of a function if there
exists $\delta>0$ such that the function is strictly convex on one side of
$y$ and strictly concave on the other, within distance $\delta$. The second
derivative need not exist at $y$.

\begin{lemma}[Critical geometry of $\Anet$]
\label{lem:critical}
\tagNE\ Let $\mathcal Z=[p,q]$ be the zero run between an excursion of sign $\sigma_-$ ending
at $p$ and the next excursion of sign $\sigma_+$ starting at $q$.
\begin{enumerate}[label=\textup{(\roman*)}]
\item Every apex $\tau_\nu$ is a strict inflection point.
\item If $\sigma_+=-\sigma_-$, then $\mathcal Z$ is an extremal plateau; it is a
strict local extremum exactly when $p=q$.
\item If $\sigma_+=\sigma_-$, then $\Anet$ has no local extremum on
$\mathcal Z$. Its concavity has sign $-\sigma_-$ immediately to the left
and $+\sigma_+$ immediately to the right. If $p=q$ the point is a strict inflection; if
$p<q$, no point of $\mathcal Z$ is a strict inflection, and the change of concavity
is carried by the whole plateau.
\item Strict local extrema occur exactly at isolated crossing zeros, and
plateaus of extrema exactly at crossing runs of positive length; every point
of such a plateau is a non-strict local extremum.
\end{enumerate}
\end{lemma}

\begin{proof}
On the descending side of the left excursion,
$\sgn\Anet'=\sigma_-$ and $\sgn\Anet''=-\sigma_-$. On the ascending side of
the right excursion, $\sgn\Anet'=\sigma_+$ and $\sgn\Anet''=\sigma_+$. Inside an excursion
$\Anet'=\Ms\ne0$, so no extremum occurs there. On
$\mathcal Z$, $\Anet'=0$. The assertions follow, and at every apex the sign of
$\Anet''$ reverses.
\end{proof}


\section{Internal correlations of an excursion}
\label{sec:internal_corr}

For an integer $L\ge1$ and a real block $x_1,\dots,x_L$ (local notation:
here $x_j$ are the entries of the block, not excursion endpoints) put
\begin{equation}
\label{eq:block_def}
S_t=\sum_{j=1}^t x_j,\qquad
C_h=\sum_{j=1}^{L-h}x_jx_{j+h}.
\end{equation}

\begin{lemma}[Triangular identity]
\label{lem:triangular-corr}
\tagNE\ For every real finite sequence,
\begin{equation}
\label{eq:general-triangular}
\sum_{t=1}^{L-1}S_t^2+\sum_{h=1}^{L-1}hC_h
=S_L\sum_{t=1}^{L-1}S_t.
\end{equation}
\end{lemma}

\begin{proof}
Expanding the partial sums,
\begin{equation}
\sum_{t=1}^{L}S_t^2
=\sum_{i,j=1}^{L}x_ix_j\bigl(L+1-\max(i,j)\bigr).
\end{equation}
Write $\max(i,j)=\frac{i+j+|i-j|}{2}$. Since
$\frac12\sum_{i,j}(i+j)x_ix_j=S_L\sum_i ix_i$ and each unordered pair at
lag $h=|i-j|\ge1$ occurs twice,
\begin{equation}
\sum_{t=1}^{L}S_t^2=(L+1)S_L^2-S_L\sum_{i=1}^{L}ix_i
-\frac12\sum_{i,j=1}^{L}x_ix_j|i-j|
=(L+1)S_L^2-S_L\sum_{i=1}^{L}ix_i-\sum_{h=1}^{L-1}hC_h .
\end{equation}
Moreover $\sum_i ix_i=LS_L-\sum_{t=1}^{L-1}S_t$.
Hence
$\sum_{t=1}^{L}S_t^2=S_L^2+S_L\sum_{t<L}S_t-\sum_hhC_h$; subtract $S_L^2$.
\end{proof}

For a bridge, $S_L=0$, so the preceding identity gives
\begin{equation}
\label{eq:bridge-triangular}
\sum_{t=1}^{L-1}S_t^2=-\sum_{h=1}^{L-1}hC_h.
\end{equation}

\begin{remark}[Hermitian form]
\label{rem:hermitian}
For complex $x_j$ and $C_h=\sum_{j=1}^{L-h}x_j\overline{x_{j+h}}$, the same
computation applies. The ordered pairs $(i,j)$ and $(j,i)$ at lag $h$
contribute $x_i\overline{x_j}+x_j\overline{x_i}$, so each lag contributes
$2\Re C_h$ to $\sum_{i,j}x_i\overline{x_j}|i-j|$, and the cross term
becomes $\Re(\overline{S_L}\sum_i ix_i)$. Hence
\begin{equation}
\sum_{t=1}^{L-1}|S_t|^2+\sum_{h=1}^{L-1}h\Re C_h
=\Re\!\left(\overline{S_L}\sum_{t=1}^{L-1}S_t\right).
\end{equation}
\end{remark}

For a complete excursion orient the increments by
\begin{equation}
\label{eq:orientation}
x_j=\sigma_\nu\mu(x_{0,\nu}+j),\qquad
S_t=|M(x_{0,\nu}+t)|.
\end{equation}
Then $S_t>0$ for $1\le t<\ell_\nu$ and $S_{\ell_\nu}=0$. Put
\begin{equation}
\label{eq:Q_def}
Q_\nu=\sum_{t=1}^{\ell_\nu-1}S_t^2,\qquad
C_h^{(\nu)}=\sum_{j=1}^{\ell_\nu-h}x_jx_{j+h},
\qquad\text{so that}\qquad
Q_\nu=-\sum_{h=1}^{\ell_\nu-1}hC_h^{(\nu)}.
\end{equation}

\begin{lemma}[Zeroth moment and extreme lag]
\label{lem:zeroth-corr}
\tagNE\ With $q_\nu=\sum_jx_j^2$,
\begin{equation}
\label{eq:zeroth-corr}
\sum_{h=1}^{\ell_\nu-1}C_h^{(\nu)}=-\frac{q_\nu}{2},\qquad
C_{\ell_\nu-1}^{(\nu)}=-1,
\end{equation}
and $q_\nu$ is even and at least $2$.
\end{lemma}

\begin{proof}
$0=(\sum_jx_j)^2=q_\nu+2\sum_hC_h^{(\nu)}$. The first oriented step is $+1$;
immediately before the first return the path is at height $1$, so the last
step is $-1$. The bridge condition forces equally many positive and negative
active steps.
\end{proof}

\begin{proposition}[Quadratic decomposition by excursions]
\label{prop:quadratic-decomp}
\tagNE\ For real $x\ge1$ let $u(x)$ and $L_x$ be as in
Definition~\ref{def:Ms-terminal}, and for $1\le j\le L_x$ put
$y_j=\mu(u(x)+j)$ and $C_h^{\rm term}=\sum_{j=1}^{L_x-h}y_jy_{j+h}$. Then
\begin{equation}
\label{eq:quadratic-decomp-general}
\sum_{n\le x}M(n)^2
=\sum_{\nu:x_{1,\nu}\le u(x)}Q_\nu+Q_{\rm term}(x),
\end{equation}
\begin{equation}
\label{eq:terminal-Q-exact}
Q_{\rm term}(x)=M(\lfloor x\rfloor)^2
+M(\lfloor x\rfloor)\sum_{t=1}^{L_x-1}M(u(x)+t)
-\sum_{h=1}^{L_x-1}hC_h^{\rm term},
\end{equation}
with $Q_{\rm term}(x)=0$ if $L_x=0$. At a zero $x$ of $M$,
\begin{equation}
\label{eq:quadratic-decomp-zero}
\sum_{n\le x}M(n)^2
=\sum_{x_{1,\nu}\le x}Q_\nu
=-\sum_{x_{1,\nu}\le x}\sum_hhC_h^{(\nu)}.
\end{equation}
\end{proposition}

\begin{proof}
Every integer up to $u(x)$ lies either in a zero run or in a unique excursion
interior. Apply Lemma~\ref{lem:triangular-corr} to the terminal block and
add its final term $S_{L_x}^2$.
\end{proof}

\begin{lemma}[Elementary bound for an open terminal block]
\label{lem:terminal-Q-elementary}
\tagNE\ Let $u$ be a zero of $M$ and $L\ge0$. Then $|M(u+r)|\le r$ for
$1\le r\le L$, and
\begin{equation}
\label{eq:terminal-Q-elementary}
\sum_{r=1}^{L}M(u+r)^2\le\frac{L(L+1)(2L+1)}6 .
\end{equation}
In particular $Q_{\rm term}(x)\ll L_x^3$. The bridge bound
\eqref{eq:Q_upper}, whose leading coefficient is $1/12$ rather than $1/3$,
applies only to complete excursions.
\end{lemma}

\begin{proof}
$M(u+r)=\sum_{j\le r}\mu(u+j)$ with $|\mu|\le1$; square and sum.
\end{proof}


Thus the quadratic sum sees only lag correlations internal to excursions:
pairs of M\"obius increments in different completed excursions do not
contribute.

Let
\begin{equation}
\label{eq:beta_def}
\mathfrak b_\nu=\sum_{t=1}^{\ell_\nu-1}S_t=\sum_{n\in I_\nu}|M(n)|.
\end{equation}
Cauchy--Schwarz gives
\begin{equation}
\label{eq:beta-Q}
\mathfrak b_\nu^2\le(\ell_\nu-1)Q_\nu,
\end{equation}
with equality exactly for a flat excursion of height $1$. To separate
support from sign geometry, let $P_h^{(\nu)}=\#\{j:x_jx_{j+h}\ne0\}$ and let
$D_h^{(\nu)}$ count the opposite-sign pairs among them. Then
\begin{equation}
\label{eq:C-P-D}
C_h^{(\nu)}=P_h^{(\nu)}-2D_h^{(\nu)},\qquad
Q_\nu=\sum_{h=1}^{\ell_\nu-1}h\bigl(2D_h^{(\nu)}-P_h^{(\nu)}\bigr).
\end{equation}

\begin{proposition}[Sharp bounds at fixed length and height]
\label{prop:sharp-LH}
\tagNE\ Let $L\ge2$ be an integer and let $(S_t)_{0\le t\le L}$ be integers
with $S_0=S_L=0$, $S_t\ge1$ for $0<t<L$, $S_t-S_{t-1}\in\{-1,0,1\}$, and
$H=\max S_t$. Then $L\ge2H$ and
\begin{align}
H^2+L-2H&\le\mathfrak b\le H(L-H),\label{eq:beta-LH}\\
\frac{2H^3+H}{3}+L-2H&\le Q\le
\frac{2H^3+H}{3}+(L-2H)H^2.\label{eq:Q-LH}
\end{align}
All four bounds are attained. In particular
\begin{equation}
\label{eq:height_Q_min}
Q\ge\frac{2H^3+H}{3},\qquad Q\ge L-1.
\end{equation}
\end{proposition}

\begin{proof}
Before and after the first visit to $H$ the path visits each level
$1,\dots,H-1$. Together with one visit to $H$, these $2H-1$ designated
interior times contribute $2(1+\cdots+(H-1))+H=H^2$ to $\mathfrak b$ and
$2\sum_{j<H}j^2+H^2=(2H^3+H)/3$ to $Q$. The remaining $L-2H$ interior times
lie at heights between $1$ and $H$, which gives the four inequalities. The
lower bounds are attained by inserting the $L-2H$ remaining steps as
horizontal steps at height $1$ into the tent path of height $H$; the upper
bounds by inserting all of them as horizontal steps at height $H$. Both
constructions keep $S_t\ge1$ on $0<t<L$ and have increments in
$\{-1,0,1\}$.
\end{proof}

\begin{corollary}[Sharp fixed-length bounds]
\label{cor:fixed-L}
\tagNE\ For $L\ge2$,
\begin{equation}
\mathfrak b\le\left\lfloor\frac{L^2}{4}\right\rfloor,\qquad
Q\le
\begin{cases}
L(L^2+2)/12,&L\text{ even},\\
L(L^2-1)/12,&L\text{ odd},
\end{cases}
\end{equation}
and in particular
\begin{equation}
\label{eq:Q_upper}
Q\le\frac{L^3+2L}{12}.
\end{equation}
Both bounds are attained by the tent path, with one horizontal step at
maximal height when $L$ is odd.
\end{corollary}

\begin{proof}
Use $S_t\le\min(t,L-t)$ and sum this bound and its square.
\end{proof}


\section{Where \texorpdfstring{$\zeta$}{zeta} enters}
\label{sec:zeta}

The net area curve $\Anet$ depends nonlinearly on $M$. The dependence enters through
$\sqrt{H_\nu}$ and through the selection of
$x_{0,\nu},\tau_\nu,x_{1,\nu}$; no closed expression of its Mellin transform
in terms of $\zeta$ is known. The nodes
$x_{0,\nu},x_{1,\nu}$ are zeros of $M$. If RH holds and all zeros are
simple, there is a sequence $T_k\to\infty$ such that, for non-integer $x$,
\begin{equation}
\label{eq:explicitM}
  M(x)=\lim_{k\to\infty}\sum_{|\gamma|<T_k}\frac{x^\rho}{\rho\,\zeta'(\rho)}-2
  +\sum_{n\ge1}\frac{(-1)^{n-1}(2\pi/x)^{2n}}{(2n)!\,n\,\zeta(2n+1)}
\end{equation}
\cite[\S14.27]{Titchmarsh1986}. For $\widetilde{A}(x):=\int_0^xM(t)\,dt$ one has
$\widetilde A(1)=0$ (the value of $M$ at the single point $1$ does not
affect the integral).  Integration by parts therefore gives, for
$\Re s>1$,
\begin{equation}
\label{eq:mellin_A}
  \int_1^\infty \widetilde{A}(x)x^{-s-2}\,dx=\frac{1}{s(s+1)\zeta(s)}.
\end{equation}
The residues of $x^{s+1}/(s(s+1)\zeta(s))$ are
\begin{equation}
\label{eq:explicitA}
  \sum_\rho\frac{x^{\rho+1}}{\rho(\rho+1)\zeta'(\rho)}-2x+12
  +(\text{residues at the trivial zeros}),
\end{equation}
where $x/\zeta(0)=-2x$ is the residue at $s=0$ and
$1/((-1)\zeta(-1))=12$ the residue at $s=-1$. Under RH and simplicity of the zeros, the contour argument of
\cite[\S14.27]{Titchmarsh1986}, with the extra decay $(s+1)^{-1}$, turns
\eqref{eq:explicitA} into an identity for $\widetilde A(x)$ with the
$\rho$-sum taken as $\lim_k\sum_{|\gamma|<T_k}$; absolute convergence of
the $\rho$-sum needs moment bounds for $1/\zeta'(\rho)$ of the type assumed
in \cite{Ng2004}. Formula \eqref{eq:explicitA} is not used in any proof. The area of $\Msk$ differs from
$\widetilde A$ by explicit local terms.

\begin{proposition}[Dipole identity]
\label{prop:dipoles}
\tagNE\ For each excursion let its \emph{legs} be the integer sets
$J=(u_\nu,\tau_\nu]\cap\mathbb{Z}$ and
$(\tau_\nu,v_\nu]\cap\mathbb{Z}$, with midpoints $c_J=(u_\nu+\tau_\nu)/2$
and $(\tau_\nu+v_\nu)/2$. Then for every zero $v_k$,
\begin{equation}
\label{eq:dipoles}
  \int_0^{v_k}\Msk(t)\,dt=\widetilde{A}(v_k)
  +\sum_{J\subset(0,v_k]}\ \sum_{m\in J}\mu(m)\,(m-c_J).
\end{equation}
\end{proposition}

\begin{proof}
Since $M$ is constant on each $[n,n+1)$, $\int_a^bM(t)\,dt=\sum_{n=a}^{b-1}M(n)$
for integers $a<b$. On an ascent leg, $\int_{u}^{\tau}\Msk=\sigma_\nu H a/2$ and
$\int_u^\tau M=\sum_{m=u+1}^{\tau}\mu(m)(\tau-m)$, while
$\sigma_\nu H=\sum_{m=u+1}^{\tau}\mu(m)$; the difference is
$\sum_m\mu(m)(a/2-\tau+m)=\sum_m\mu(m)(m-c_J)$. On a descent leg,
$\int_\tau^v\Msk=\sigma_\nu Hb/2$ and
$\int_\tau^vM=b\,\sigma_\nu H+\sum_{m=\tau+1}^{v}\mu(m)(v-m)$ with
$\sum_{m=\tau+1}^{v}\mu(m)=-\sigma_\nu H$; the difference is
$\sum_m\mu(m)(m-(v-b/2))$, and $v-b/2=(\tau+v)/2=c_J$. On zero runs
both integrands vanish. Sum over legs.
\end{proof}


\section{Bounds for the area curves}
\label{sec:bounds}

\begin{lemma}[Transfer]
\label{lem:transfer}
\tagNE\ Let $x\ge1$ and let $x^*\ge x$ be a zero of $M$ (in a computation truncated
at $N\ge x$, take $x^*=N$ and the terminal skeleton of
Definition~\ref{def:Ms-terminal}). If $|M(n)|\le C_0\,n^{1/2}$ for all
$1\le n\le x^*$ (so necessarily $C_0\ge1$), then
\begin{equation}
\label{eq:transfer_bound}
  |\Ms(t)|\le C_0^{1/2}t^{1/4}\quad(0\le t\le x),\qquad
  \Ap(x)+|\Am(x)|\le\tfrac45\,C_0^{1/2}x^{5/4}.
\end{equation}
\end{lemma}

\begin{proof}
For $t\ge\tau_\nu$, $|\Ms(t)|\le\sqrt{H_\nu}\le C_0^{1/2}\tau_\nu^{1/4}\le
C_0^{1/2}t^{1/4}$. For $u_\nu\le t\le\tau_\nu$ the chord from $(u_\nu,0)$ to
$(\tau_\nu,\sqrt{H_\nu})$ lies below the chord of $C_0^{1/2}t^{1/4}$ over
$[u_\nu,\tau_\nu]$, whose endpoints $C_0^{1/2}u_\nu^{1/4}\ge0$ and
$C_0^{1/2}\tau_\nu^{1/4}\ge\sqrt{H_\nu}$ (as $\tau_\nu\le x^*$) are at
least as high; that chord lies below $C_0^{1/2}t^{1/4}$ by concavity.
Integrate.
\end{proof}

\begin{lemma}[A concave majorant]
\label{lem:concave-majorant}
\tagNE\ For $y>e$ let $\phi(y)=(\log y)^{3/5}(\log\log y)^{-1/5}$ and, for
$\kappa>0$, $F(y)=y^{1/2}e^{-\kappa\phi(y)}$. There is $y_0=y_0(\kappa)$ such that $F$
is increasing and concave on $[y_0,\infty)$.
\end{lemma}

\begin{proof}
With $y=e^s$, $s>1$, and $g(s)=s^{3/5}(\log s)^{-1/5}$, one has
$F=\exp(s/2-\kappa g(s))$. Put $h(s)=\frac12-\kappa g'(s)$. Then
\begin{equation}
yF'(y)=F\,h(s),\qquad
y^2F''(y)=F\bigl(h(s)^2-h(s)-\kappa g''(s)\bigr).
\end{equation}
Since $g'(s)\to0$ and $g''(s)\to0$, $h(s)\to\frac12$ and
$h^2-h-\kappa g''\to-\frac14$. Hence $F'>0$ and $F''<0$ for large $y$.
\end{proof}

\begin{corollary}
\label{cor:uncond}
\tagC\ From the Vinogradov--Korobov bound \cite{Walfisz1963}: there is an
absolute constant $c>0$ with
\begin{equation}
\label{eq:VK_M}
  M(x)\ll x\exp\bigl(-c\,\phi(x)\bigr)\qquad(x>e),
\end{equation}
one obtains
\begin{equation}
\label{eq:uncond_area}
  \Ap(x)+|\Am(x)|\ll x^{3/2}\exp\bigl(-\tfrac c2\,\phi(x)\bigr).
\end{equation}
\end{corollary}

\begin{proof}
Let $K$ be such that $|M(y)|\le Ky\,e^{-c\phi(y)}$ for $y\ge y_1$, and let
$F$ be as in Lemma~\ref{lem:concave-majorant} with $\kappa=c/2$, so that
$\sqrt{|M(y)|}\le K^{1/2}F(y)$ for $y\ge y_1$. Only finitely many excursions
start below $Y_0:=\max(y_0,y_1)$; they are finite, so they contribute
$O(1)$ to the area. Consider an excursion with $u_\nu\ge Y_0$. Its apex
satisfies $\sqrt{H_\nu}=\sqrt{|M(\tau_\nu)|}\le K^{1/2}F(\tau_\nu)$; the
bound is taken at the apex $\tau_\nu$, which lies in the excursion, and not
at its future endpoint $v_\nu$. For $t\ge\tau_\nu$ monotonicity of $F$ gives
$|\Ms(t)|\le\sqrt{H_\nu}\le K^{1/2}F(t)$. For $u_\nu\le t\le\tau_\nu$,
$|\Ms(t)|$ is the chord from $(u_\nu,0)$ to $(\tau_\nu,\sqrt{H_\nu})$. This
chord lies below the chord from $(u_\nu,K^{1/2}F(u_\nu))$ to
$(\tau_\nu,K^{1/2}F(\tau_\nu))$, which has both endpoints at least as high,
and the latter lies below $K^{1/2}F$ on $[u_\nu,\tau_\nu]\subset[y_0,\infty)$
by concavity. Hence $|\Ms(t)|\le K^{1/2}F(t)$ for all $t\ge Y_0$ outside the
finitely many early excursions, and since $F$ is increasing,
\begin{equation}
\Ap(x)+|\Am(x)|\le O(1)+K^{1/2}\int_{Y_0}^xF(t)\,dt
\le O(1)+K^{1/2}xF(x)=O(1)+K^{1/2}x^{3/2}e^{-\frac c2\phi(x)}.
\end{equation}
\end{proof}

\begin{corollary}
\label{cor:rh}
\tagC\ If RH holds, then
\begin{equation}
\label{eq:rh_area}
  \Ap(x)+|\Am(x)|\ll_\eps x^{5/4+\eps}.
\end{equation}
Consequently, if $\Ap(x)+|\Am(x)|\gg x^{5/4+\delta}$ along an unbounded
sequence for some fixed $\delta>0$, then RH is false. By
Corollary~\ref{cor:area-exponent} the converse also holds: \eqref{eq:rh_area}
is equivalent to RH.
\end{corollary}

\begin{proof}
Under RH, $|M(y)|\le C_\eps y^{1/2+\eps}$, so
$\sqrt{H_\nu}\le C_\eps^{1/2}\tau_\nu^{1/4+\eps/2}$. The chord argument of
Lemma~\ref{lem:transfer}, with the concave majorant
$C_\eps^{1/2}t^{1/4+\eps/2}$ in place of $C_0^{1/2}t^{1/4}$, gives
$|\Ms(t)|\le C_\eps^{1/2}t^{1/4+\eps/2}$ and, after integration,
$\Ap(x)+|\Am(x)|\ll_\eps x^{5/4+\eps/2}$.
\end{proof}

\begin{remark}[Endpoint-indexed bounds]
\label{rem:onedirection}
From $\ell_\nu\ge2H_\nu$ and $\sqrt{H_\nu}\ell_\nu\ll v_\nu^{5/4+\eps}$ one
deduces $H_\nu\ll v_\nu^{5/6+O(\eps)}$, indexed by the future endpoint
$v_\nu$. For a general family of disjoint intervals this gives no bound at
interior points. For $M$ itself Corollary~\ref{cor:KP-local-returns} gives
$v_\nu\le u_\nu^{1+\eps}$ for all late excursions, so the bound becomes
$|M(t)|\ll t^{5/6+O(\eps)}$; the sharp statement is the exact area exponent
of Corollary~\ref{cor:area-exponent}.
\end{remark}

\begin{remark}[Oscillation of the heights]
\label{rem:omega}
The Mertens conjecture is false \cite{OdlyzkoTeRiele1985,KotnikTeRiele2006};
the current bounds are due to Hurst~\cite{Hurst18}, who proves
$\limsup M(x)/\sqrt{x}>1.826054$ and
$\liminf M(x)/\sqrt{x}<-1.837625$. The excursion containing a point $n$ with
$M(n)>1.826054\sqrt n$ has $u_\nu<n$, so $H_\nu>1.826054\,u_\nu^{1/2}$ for
infinitely many positive excursions, and similarly with $1.837625$ for
negative ones. A lower bound indexed by $v_\nu$ does not follow, since
$v_\nu$ may be much larger than $n$.
\end{remark}


\section{Global criteria equivalent to RH}
\label{sec:criteria}

The Mellin identity and the global $L^1/L^2$ equivalences below are
consequences of Littlewood's criterion and partial summation; the excursion
decompositions and pruning lemmas rewrite them in terms of excursions.
For $\Re s>1$,
\begin{equation}
\label{eq:mellin_M}
  \int_1^\infty M(x)x^{-s-1}\,dx=\frac{1}{s\,\zeta(s)} .
\end{equation}
If $\int_1^\infty|M(x)|x^{-\varsigma-1}dx<\infty$ for every $\varsigma>1/2$, the
left side continues analytically to $\Re s>1/2$, so $\zeta$ has no zeros
there.

\begin{proposition}[$L^1$ criterion]
\label{prop:L1}
\tagC\ RH holds if and only if, for every $\eps>0$,
\begin{equation}
\label{eq:L1_criterion}
  \sum_{v_\nu\le x}\mathfrak b_\nu+\mathfrak b_{\rm term}(x)=\sum_{n\le x}|M(n)|
  \ll x^{3/2+\eps}.
\end{equation}
\end{proposition}

\begin{proof}
RH gives $M(x)\ll_\eps x^{1/2+\eps}$, hence the bound. Conversely, fix
$\varsigma_0>1/2$ and $0<\eps<\varsigma_0-1/2$. Partial summation with
$\sum_{n\le x}|M(n)|\ll x^{3/2+\eps}$ gives
$\sum_n|M(n)|n^{-\varsigma_0-1}<\infty$; since $M$ is constant on
$[n,n+1)$, $\int_1^\infty|M(x)|x^{-\Re s-1}dx\le\int_1^\infty|M(x)|x^{-\varsigma_0-1}dx<\infty$
uniformly for $\Re s\ge\varsigma_0$. Hence $F(s):=\int_1^\infty M(x)x^{-s-1}dx$
converges absolutely and uniformly there and is holomorphic on
$\Re s>\varsigma_0$. By \eqref{eq:mellin_M}, $F=1/(s\zeta(s))$
on $\Re s>1$; by uniqueness of analytic continuation the meromorphic function
$1/(s\zeta(s))$ coincides with $F$ on $\Re s>\varsigma_0$, so it has no poles
there, i.e.\ $\zeta$ has no zeros there. As $\varsigma_0>1/2$ is arbitrary,
RH follows.
\end{proof}

\begin{lemma}[Pruning]
\label{lem:pruning}
\tagNE\ $\mathfrak b_\nu\le H_\nu(\ell_\nu-1)\le\ell_\nu^2/2$, hence
\begin{equation}
\label{eq:pruning}
  \sum_{v_\nu\le x,\ \ell_\nu\le L_0}\mathfrak b_\nu\le\frac{xL_0}{2}.
\end{equation}
\end{lemma}

\begin{corollary}[Reduction to long excursions]
\label{cor:long}
\tagC\ For fixed $\delta\ge0$ and $L_0=x^{1/2-\delta}$, RH holds if and only if
\begin{equation}
\label{eq:long_reduction}
  \sum_{\substack{v_\nu\le x\\ \ell_\nu>L_0}}\mathfrak b_\nu
  +\mathbf 1_{\{L_x>L_0\}}\mathfrak b_{\rm term}(x)\ll x^{3/2+\eps}
  \quad\text{for every }\eps>0.
\end{equation}
At most $x/L_0$ long excursions are involved.
\end{corollary}

\begin{proof}
If RH holds, the left side is at most $\sum_{n\le x}|M(n)|\ll x^{3/2+\eps}$.
Conversely, the complete excursions with $\ell_\nu\le L_0$ contribute at
most $xL_0/2\le x^{3/2}$ by Lemma~\ref{lem:pruning}, and if $L_x\le L_0$
the terminal block contributes
$\mathfrak b_{\rm term}(x)\le\sum_{r\le L_x}r=L_x(L_x+1)/2\le L_0^2\le x$.
Hence $\sum_{n\le x}|M(n)|\ll x^{3/2+\eps}$, and
Proposition~\ref{prop:L1} applies.
\end{proof}

\begin{remark}[The $\sqrt{x}$ frontier]
\label{rem:sqrt_frontier}
The cut may depend on $\eps$: with $L_0=x^{1/2+\eps/2}$ the discarded
complete excursions contribute at most $xL_0/2=O(x^{3/2+\eps/2})$, while a
short terminal block contributes
$O(L_0^2)=O(x^{1+\eps})$.  Both are admissible in the
$x^{3/2+\eps}$ budget.  The frontier of pruning is therefore
$x^{1/2+o(1)}$;
what requires new information is a cut $x^{\theta}$ with a fixed
$\theta>1/2$. For $\ell_\nu\asymp x$ the trivial bound gives
$\mathfrak b_\nu\ll x^2$ against a budget $x^{3/2+\eps}$, and there are only
$O(1)$ such excursions.
\end{remark}

\begin{heuristic}[Renewal tail]
\label{heu:renewal}
If excursion lengths behave like those of a recurrent random walk,
\begin{equation}
\label{eq:renewal_tail}
  \#\{v_\nu\le x:\ \ell_\nu>L\}\asymp\sqrt{x}\,L^{-1/2},
\end{equation}
and with $\mathfrak b_\nu\asymp\ell_\nu^{3/2}$ the dyadic range $[L,2L]$
contributes $\asymp\sqrt{x}\,L$ to \eqref{eq:L1_criterion}, so the terminal
scales $L\asymp x$ dominate the $L^1$ sum.
\end{heuristic}

\begin{theorem}[Correlation criterion]
\label{thm:corr-criterion}
\tagC\ RH holds if and only if, for every $\eps>0$,
\begin{equation}
\label{eq:RH-corr}
  \sum_{n\le x}M(n)^2=\sum_{v_\nu\le x}Q_\nu+Q_{\rm term}(x)
  =-\sum_{v_\nu\le x}\sum_{h}hC^{(\nu)}_h+Q_{\rm term}(x)\ll x^{2+\eps}.
\end{equation}
\end{theorem}

\begin{proof}
Under RH, $|M(n)|\ll_\delta n^{1/2+\delta}$ for every $\delta>0$; with
$\delta=\eps/2$, $M(n)^2\ll_\eps n^{1+\eps}$ and the sum is $\ll x^{2+\eps}$.
Conversely, by Cauchy--Schwarz,
\begin{equation}
\sum_{n\le x}|M(n)|\le x^{1/2}\Bigl(\sum_{n\le x}M(n)^2\Bigr)^{1/2}\ll x^{3/2+\eps/2};
\end{equation}
apply
Proposition~\ref{prop:L1}.
\end{proof}

\begin{corollary}[Local sufficient conditions]
\label{cor:local-corr-sufficient}
\tagC\ Suppose that for every $\eps>0$ there is $C_\eps$ such that
\begin{equation}
\label{eq:pruned_Q}
  Q_\nu\le C_\eps\,\ell_\nu^{2+\eps}\quad\text{for every complete excursion with }
  \ell_\nu>v_\nu^{1/2}.
\end{equation}
Then RH holds.
\end{corollary}

\begin{proof}
Let $\eps_0\in(0,1)$, put $\eps=\eps_0/4$, and let $C_\eps$ be as in the
hypothesis. Split the completed excursions at $\ell_\nu=x^{1/2}$. If
$\ell_\nu\le x^{1/2}$, then \eqref{eq:Q_upper} gives
$Q_\nu\le\ell_\nu x/12+\ell_\nu/6$, and disjointness bounds the total by
$x^2/12+x/6$. If $\ell_\nu>x^{1/2}$ then $\ell_\nu>v_\nu^{1/2}$ since
$v_\nu\le x$, and the hypothesis bounds the total by
$C_\eps\sum\ell_\nu^{2+\eps}\le C_\eps(\sum\ell_\nu)^{2+\eps}\le C_\eps x^{2+\eps}$.
The terminal block is not a complete excursion, and \eqref{eq:Q_upper} does
not apply to it. By Lemma~\ref{lem:terminal-Q-elementary},
$Q_{\rm term}(x)\le L_x(L_x+1)(2L_x+1)/6$, which is $O(x^{3/2})$ if $L_x\le x^{1/2}$. If $L_x>x^{1/2}$, then $M(\lfloor x\rfloor)\ne0$.  Let
$I_\nu=(u(x),v_\nu]$ be the unique completed excursion whose open prefix
contains $\lfloor x\rfloor$; since all summands are nonnegative,
$Q_{\rm term}(x)\le Q_\nu$. Fix $A=1+\eps$. By
Corollary~\ref{cor:KP-local-returns}, $v_\nu\le u(x)^A\le x^A$ once $x$ is
large, so $\ell_\nu\le x^A$. If $\ell_\nu>v_\nu^{1/2}$ the hypothesis gives
$Q_\nu\le C_\eps x^{A(2+\eps)}$; otherwise \eqref{eq:Q_upper} gives
$Q_\nu\ll\ell_\nu^3\le v_\nu^{3/2}\le x^{3A/2}$. In all cases
$Q_{\rm term}(x)\ll_\eps x^{(1+\eps)(2+\eps)}\le x^{2+4\eps}=x^{2+\eps_0}$. This step uses the arithmetic
input of Corollary~\ref{cor:KP-local-returns}: for a general family of
disjoint intervals, the future return of the terminal block is not
controlled by $x$, and a bound for complete excursions would not control
$Q_{\rm term}(x)$. Altogether
$\sum_{n\le x}M(n)^2\ll_{\eps_0}x^{2+\eps_0}$ for every $\eps_0\in(0,1)$,
and Theorem~\ref{thm:corr-criterion} gives RH.
\end{proof}

\begin{heuristic}[Ballistic excursions]
\label{heu:local_fails}
In a model where, after a zero of $M$, the next values of $\mu$ behave
locally at random with squarefree density $6/\pi^2$, the pattern ``$k$
non-zero steps of one sign followed by $k$ of the other'' follows a given
zero with probability $\ge c_0^{-k}$. If excursions with $H_\nu=k$ and
$\ell_\nu\asymp k$ occur for unbounded $k$, then by \eqref{eq:height_Q_min}
$Q_\nu/\ell_\nu^{2+\eps}\gg k^{1-\eps}\to\infty$, so the unpruned version of
\eqref{eq:pruned_Q} fails for every $\eps\in(0,1)$. Such excursions have
$\ell_\nu\ll v_\nu^{1/2}$ and do not affect the pruned condition
\eqref{eq:pruned_Q}; this is why the pruning in
Corollary~\ref{cor:local-corr-sufficient} is needed.
\end{heuristic}

\begin{theorem}[Energy criterion]
\label{thm:energy}
\tagC\ With $u$ the last zero $\le x$ and $H_{\rm run}(x)=\max_{u<n\le x}|M(n)|$,
define
\begin{equation}
\label{eq:energy_def}
  \En(x):=\frac14\Bigl(\sum_{v_\nu\le x}H_\nu\ell_\nu+(x-u)H_{\rm run}(x)\Bigr).
\end{equation}
For a complete excursion let $\mathrm{ch}_\nu$ be the chord joining
$(u_\nu,\Anet(u_\nu))$ to $(v_\nu,\Anet(v_\nu))$. By \eqref{eq:jump},
$\mathrm{ch}_\nu'=\sigma_\nu\sqrt{H_\nu}/2$, so $\tfrac14H_\nu\ell_\nu=\int_{I_\nu}(\mathrm{ch}_\nu')^2$
is the Dirichlet energy of the affine interpolant $\mathrm{ch}_\nu$ joining
the endpoint values of $\Anet$ on $I_\nu$; it is not the Dirichlet energy
$\int_{I_\nu}\Ms^2$ of $\Anet$ itself. Then
\begin{equation}
\label{eq:energy_sandwich}
  \sum_{n\le x}|M(n)|\le4\En(x)\le x\max_{n\le x}|M(n)|,
\end{equation}
and RH holds if and only if $\En(x)\ll x^{3/2+\eps}$ for every $\eps>0$.
\end{theorem}

\begin{proof}
$\mathfrak b_\nu\le H_\nu\ell_\nu$, $\mathfrak b_{\rm term}(x)\le(x-u)H_{\rm run}(x)$ and
$\sum\ell_\nu+(x-u)\le x$; apply Proposition~\ref{prop:L1}.
\end{proof}


\subsection{Exact logarithmic exponents from Pintz's mean and late-maximal-order theorem}
\label{subsec:pintz-exponents}

The criteria above are upper-bound equivalences.  Theorems~2.1 and~2.2 of
Pintz~\cite[Theorems~2.1--2.2, equations~(2.14)--(2.15)]{Pintz2026}
determine, in logarithmic scale, the mean modulus and late maximal order of
$|M|$ in terms of the rightmost zeros of $\zeta$. In \cite{Pintz2026}, $D_M(x)=x^{-1}\int_0^x|M|$ (equivalently $\int_1^x$),
$S_{M,\delta}(x)=\max_{x^{1-\delta}\le u\le x}|M(u)|$ for each fixed
$\delta>0$, and $Z(x)=\max_{\gamma>0}x^{\beta}/\gamma$; Theorem~2.2 there,
equation (2.15), states $\log D_M(x)\sim\log S_{M,\delta}(x)\sim\log Z(x)$ as
$x\to\infty$ for every fixed $\delta>0$, and (6.1) there gives
$\log Z(x)\sim\Theta\log x$ (both cases $\Theta<1$ and $\Theta=1$). Since
$\Theta>0$, this is the statement below.
For $\Theta<1$ the logarithmic-scale lower bound for $D_M$ already follows
from Pintz's inequality $D_M(Y)\ge Y^{\beta_0}/(6|\gamma_0|^3)$ for
$Y>e^{|\gamma_0|+4}$, valid for every zero $\beta_0+i\gamma_0$
\cite{Pintz1983}, and the upper bound from the classical estimate
$M(x)\ll_\eps x^{\Theta+\eps}$; the case $\Theta=1$ is the new content of
\cite{Pintz2026}.
We record the consequence for the excursion quantities used in this paper.

\begin{theorem}[Pintz, logarithmic mean and late maximal order
{\cite[Theorems~2.1 and~2.2]{Pintz2026}}]
\label{thm:pintz-2026}
\tagL\ Let $\Theta=\sup\{\Re\rho:\zeta(\rho)=0,\ 0<\Re\rho<1\}$.  Then
\begin{equation}
 D_M(X):=\frac1X\int_1^X|M(t)|\,dt=X^{\Theta+o(1)}.
\end{equation}
Moreover, for every fixed $\delta>0$,
\begin{equation}
 S_{M,\delta}(X):=\max_{X^{1-\delta}\le t\le X}|M(t)|
 =X^{\Theta+o_\delta(1)}.
\end{equation}
Here the equalities are in logarithmic scale, equivalently
\begin{equation}
 \frac{\log D_M(X)}{\log X}\to\Theta,
 \qquad
 \frac{\log S_{M,\delta}(X)}{\log X}\to\Theta
 \quad\text{for each fixed }\delta>0.
\end{equation}
\end{theorem}

\begin{corollary}[Exact exponents of absolute Mertens mass, quadratic mass, and chord energy]
\label{cor:pintz-excursion-exponents}
\tagC\ Put
\begin{equation}
 B_1(X):=\sum_{n\le X}|M(n)|,
 \qquad
 B_2(X):=\sum_{n\le X}M(n)^2.
\end{equation}
Then
\begin{equation}
\label{eq:pintz-B1-B2}
 B_1(X)=X^{1+\Theta+o(1)},
 \qquad
 B_2(X)=X^{1+2\Theta+o(1)},
 \qquad
 \En(X)=X^{1+\Theta+o(1)}.
\end{equation}
Consequently
\begin{equation}
\label{eq:Theta-from-excursions}
 \Theta=\lim_{X\to\infty}\left(\frac{\log B_1(X)}{\log X}-1\right)
 =\frac12\lim_{X\to\infty}
 \left(\frac{\log B_2(X)}{\log X}-1\right)
 =\lim_{X\to\infty}\left(\frac{\log\En(X)}{\log X}-1\right).
\end{equation}
In particular RH is equivalent to any one of the three logarithmic exponents
being respectively $3/2$, $2$, and $3/2$.
\end{corollary}

\begin{proof}
Let $N=\lfloor X\rfloor$.  Since $M(t)=M(n)$ on $[n,n+1)$,
\begin{equation}
 \int_1^X|M(t)|\,dt
 =\sum_{n=1}^{N-1}|M(n)|+(X-N)|M(N)|,
\end{equation}
so
\begin{equation}
 B_1(X)=\int_1^X|M(t)|\,dt+O(|M(N)|+1).
\end{equation}
The boundary term is $O(X)$, whereas Theorem~\ref{thm:pintz-2026} gives
$\int_1^X|M(t)|\,dt=X^{1+\Theta+o(1)}$ and $\Theta\ge1/2$; hence
$B_1(X)=X^{1+\Theta+o(1)}$.
We first note that
\begin{equation}
\label{eq:global-M-max-exponent}
 \max_{n\le X}|M(n)|=X^{\Theta+o(1)}.
\end{equation}
The lower bound follows from $S_{M,\delta}$ with any fixed $\delta$.  For
the upper bound, if $\Theta<1$, fix once and for all
$\delta\in(1-\Theta,1)$, independent of $X$, and split the range
$n\le X$ at $X^{1-\delta}$. On $n<X^{1-\delta}$ the trivial bound gives
$|M(n)|\le X^{1-\delta}=o(X^\Theta)$; on $X^{1-\delta}\le n\le X$,
Theorem~\ref{thm:pintz-2026} with this fixed $\delta$ gives
$|M(n)|\le S_{M,\delta}(X)=X^{\Theta+o(1)}$. When $\Theta=1$ the trivial
bound $|M(n)|\le n$ suffices.

Cauchy--Schwarz and \eqref{eq:global-M-max-exponent} now give
\begin{equation}
 \frac{B_1(X)^2}{X}\le B_2(X)
 \le X\max_{n\le X}|M(n)|^2,
\end{equation}
so both sides have exponent $1+2\Theta$.  Finally
\eqref{eq:energy_sandwich} squeezes $4\En(X)$ between $B_1(X)$ and
\begin{equation}
 X\max_{n\le X}|M(n)|,
\end{equation}
both of exponent $1+\Theta$.
\end{proof}

\begin{corollary}[Length localization of the area and quadratic mass]
\label{cor:pintz-length-localization}
\tagC\ Fix $0\le\alpha<\Theta$.  Among the complete excursions ending at most $X$,
the total contribution of those with $\ell_\nu\le X^\alpha$ is negligible:
\begin{equation}
 \sum_{\substack{v_\nu\le X\\ \ell_\nu\le X^\alpha}}\mathfrak b_\nu
 +\mathbf 1_{\{L_X\le X^\alpha\}}\mathfrak b_{\rm term}(X)
 =o(B_1(X)),
\end{equation}
and
\begin{equation}
 \sum_{\substack{v_\nu\le X\\ \ell_\nu\le X^\alpha}}Q_\nu
 +\mathbf 1_{\{L_X\le X^\alpha\}}Q_{\rm term}(X)
 =o(B_2(X)).
\end{equation}
Thus, on the logarithmic scale forced by the zeta zeros, essentially all
$L^1$ and $L^2$ mass is carried by excursion blocks longer than
$X^\alpha$ for every $\alpha<\Theta$.
\end{corollary}

\begin{proof}
Lemma~\ref{lem:pruning} bounds the first complete-excursion sum by
$O(X^{1+\alpha})$, while the short terminal block contributes
$O(X^{2\alpha})$.  Both are $o(X^{1+\Theta+o(1)})$ when $\alpha<\Theta$.
For the quadratic mass use \eqref{eq:Q_upper} and disjointness:
\begin{equation}
 \sum_{\ell_\nu\le X^\alpha}Q_\nu\ll X^{1+2\alpha},
 \qquad
 Q_{\rm term}(X)\ll X^{3\alpha}
 \quad(L_X\le X^\alpha).
\end{equation}
Because $\alpha<\Theta\le1$, both exponents are strictly below
$1+2\Theta$.
\end{proof}


\begin{corollary}[Exact exponent of the skeleton area]
\label{cor:area-exponent}
\tagN\ The proof uses the return localization (RL) of
Corollary~\ref{cor:KP-local-returns}, which holds unconditionally. Put $\mathcal A_{\rm abs}(x):=\Ap(x)+|\Am(x)|=\int_0^x|\Ms(t)|\,dt$. Then
\begin{equation}
\mathcal A_{\rm abs}(x)=x^{1+\Theta/2+o(1)}\qquad(x\to\infty).
\end{equation}
In particular $\mathrm{RH}\iff\mathcal A_{\rm abs}(x)\ll_\eps x^{5/4+\eps}$
for every $\eps>0$.
\end{corollary}

\begin{proof}
Upper bound. Fix $0<\eps<2-\Theta$. By \eqref{eq:global-M-max-exponent}
there is $C_\eps$ with $H_\nu=|M(\tau_\nu)|\le C_\eps\tau_\nu^{\Theta+\eps}$.
Since $0<(\Theta+\eps)/2<1$, the function $C_\eps^{1/2}t^{(\Theta+\eps)/2}$ is
concave and increasing, and the chord argument of Lemma~\ref{lem:transfer}
gives $|\Ms(t)|\le C_\eps^{1/2}t^{(\Theta+\eps)/2}$, hence
$\mathcal A_{\rm abs}(x)\ll_\eps x^{1+(\Theta+\eps)/2}$. As $\eps$ is
arbitrary, $\limsup_x\log\mathcal A_{\rm abs}(x)/\log x\le1+\Theta/2$.

Lower bound. Let $\delta>0$ and fix $\eps\in(0,1)$ so small that
\begin{equation}
e(\eps):=\frac{1+\Theta-\eps-\frac12(1+\eps)(\Theta+\eps)}{1+\eps}
\ge1+\frac\Theta2-\frac\delta2 ;
\end{equation}
this is possible because $e(\eps)\to1+\Theta/2$ as $\eps\downarrow0$. For
$X\ge2$ let $v^*(X)$ be the least zero of $M$ with $v^*(X)\ge\lfloor X\rfloor$.
By (RL) (Corollary~\ref{cor:KP-local-returns} with $A=1+\eps$) there is
$X_1(\eps)$ such that $v^*(X)\le X^{1+\eps}$ for every $X\ge X_1(\eps)$: if
$M(\lfloor X\rfloor)=0$ then $v^*(X)=\lfloor X\rfloor$; otherwise $v^*(X)$
ends the excursion containing $\lfloor X\rfloor$, whose start $u$ satisfies
$u<X$, and $v^*(X)\le u^{1+\eps}$ once $u\ge U_{1+\eps}$, which holds for all
large $X$ because the excursions starting below $U_{1+\eps}$ are finitely many
and finite. Every $n\le X$ with $M(n)\ne0$ lies in a completed excursion
ending at most $v^*(X)$. With $\mathfrak b_\nu\le H_\nu\ell_\nu$,
$m^*:=\max_{n\le v^*(X)}|M(n)|$ and Lemma~\ref{lem:jump},
\begin{equation}
B_1(X)\le\sum_{v_\nu\le v^*(X)}H_\nu\ell_\nu
\le(m^*)^{1/2}\sum_{v_\nu\le v^*(X)}\sqrt{H_\nu}\,\ell_\nu
=2(m^*)^{1/2}\mathcal A_{\rm abs}(v^*(X)).
\end{equation}
By Corollary~\ref{cor:pintz-excursion-exponents} and
\eqref{eq:global-M-max-exponent} there is $X_2(\eps)$ such that for all
$X\ge X_2(\eps)$, $B_1(X)\ge X^{1+\Theta-\eps}$ and
$m^*\le v^*(X)^{\Theta+\eps}\le X^{(1+\eps)(\Theta+\eps)}$. Hence, for
$X\ge\max(X_1,X_2)$, using monotonicity of $\mathcal A_{\rm abs}$,
\begin{equation}
\mathcal A_{\rm abs}(X^{1+\eps})\ge\mathcal A_{\rm abs}(v^*(X))
\ge\tfrac12X^{(1+\eps)e(\eps)} .
\end{equation}
Given $y\ge\max(X_1,X_2)^{1+\eps}$, put $X=y^{1/(1+\eps)}$; then
$\mathcal A_{\rm abs}(y)\ge\frac12y^{e(\eps)}\ge\frac12y^{1+\Theta/2-\delta/2}
\ge y^{1+\Theta/2-\delta}$ for all large $y$. Since $\delta>0$ is arbitrary,
$\liminf_y\log\mathcal A_{\rm abs}(y)/\log y\ge1+\Theta/2$.
\end{proof}

\begin{remark}[Inputs]
The new part is the inequality $B_1(X)\le2(m^*)^{1/2}\mathcal A_{\rm abs}(v^*(X))$
and its use. The exponent inputs are Corollary~\ref{cor:pintz-excursion-exponents}
(from \cite{Pintz2026}, or from \cite{Pintz1983} and the classical upper bound)
and the return localization (RL) of Corollary~\ref{cor:KP-local-returns}; for an arbitrary family of disjoint
intervals the terminal block of $B_1(X)$ is not controlled by the area up to
$X$. For the $\limsup$ form of the lower bound the classical
$B_1(X)=\Omega(X^{1+\Theta-\eps})$ (Mellin transform and
Proposition~\ref{prop:L1}) suffices instead of \cite{Pintz2026}.
\end{remark}


\subsection{\texorpdfstring{Excursion blocks of $1/\zeta$}{Excursion blocks of 1/zeta}}

\begin{proposition}[Excursion blocks of $1/\zeta$]
\label{prop:excursion-mellin}
\tagNE\ For a complete excursion $I=(x_0,x_1]$ put
$D_I(s)=\sum_{x_0<n\le x_1}\mu(n)n^{-s}$. Then
\begin{equation}
\label{eq:excursion-abel}
D_I(s)=\sum_{n=x_0+1}^{x_1-1}M(n)
\bigl(n^{-s}-(n+1)^{-s}\bigr).
\end{equation}
For real $\varsigma>0$, $D_I(\varsigma)$ has the sign of the excursion and
\begin{equation}
\label{eq:excursion-mass-lower}
|D_I(\varsigma)|\ge
\varsigma\mathfrak b_I x_1^{-\varsigma-1}
\ge\varsigma H_I^2x_1^{-\varsigma-1}.
\end{equation}
For $\Re s>1$,
\begin{equation}
\label{eq:1overzeta-excursions}
\frac1{\zeta(s)}=\sum_I D_I(s),
\end{equation}
zero runs contributing nothing.
\end{proposition}

\begin{proof}
For $a_n=\mu(n)$, $A(n)=M(n)$ and $f(n)=n^{-s}$, finite summation by
parts gives
\begin{equation}
 \sum_{n=x_0+1}^{x_1}a_nf(n)
 =A(x_1)f(x_1)-A(x_0)f(x_0)
 +\sum_{n=x_0+1}^{x_1-1}A(n)\bigl(f(n)-f(n+1)\bigr).
\end{equation}
Since $M(x_0)=M(x_1)=0$, this is \eqref{eq:excursion-abel}.  For real
$\varsigma>0$, all interior $M(n)$ have the excursion sign, and
\begin{equation}
 n^{-\varsigma}-(n+1)^{-\varsigma}
 =\varsigma\int_n^{n+1}x^{-\varsigma-1}\,dx
 \ge\varsigma x_1^{-\varsigma-1}.
\end{equation}
The bound $\mathfrak b_I\ge H_I^2$ is \eqref{eq:beta-LH}.  If $M$ vanishes at
two consecutive integers the intervening $\mu(n)$ is zero, so every
$n$ with $\mu(n)\ne0$ lies in a completed excursion.  For $\Re s>1$
the ordinary Dirichlet series is absolutely convergent, so summing the
finite disjoint blocks gives \eqref{eq:1overzeta-excursions}.
\end{proof}

\begin{theorem}[Abscissa of absolute excursion-block convergence]
\label{thm:excursion-block-abscissa}
\tagC\ Let $D_I(s)$ be the completed-excursion blocks of
Proposition~\ref{prop:excursion-mellin}.  The series
\begin{equation}
 \sum_I D_I(s)
\end{equation}
has abscissa of absolute convergence, with respect to the excursion blocks,
exactly
\begin{equation}
 \boxed{\ \varsigma_{a,\mathrm{exc}}=\Theta\ }.
\end{equation}
Equivalently, the infimum of the real numbers $\varsigma_0$ for which the
series converges absolutely throughout the half-plane
$\Re s>\varsigma_0$ is $\Theta$.  More precisely, it converges absolutely
and locally uniformly for $\Re s>\Theta$ and there represents $1/\zeta(s)$,
while for every real $\varsigma<\Theta$ with $\varsigma\ne0$ one has
\begin{equation}
 \sum_I|D_I(\varsigma)|=\infty .
\end{equation}
No assertion is made on the boundary line $\Re s=\Theta$.
Consequently,
\begin{equation}
 \mathrm{RH}\iff
 \sum_I D_I(s)\ \text{converges absolutely for every }\Re s>\tfrac12.
\end{equation}
This is block absolute convergence; the ordinary termwise Dirichlet series
$\sum_n\mu(n)n^{-s}$ has abscissa of absolute convergence $1$ and, by
$M(x)\ll_\eps x^{\Theta+\eps}$, abscissa of (conditional) convergence
$\Theta$. Grouping by completed excursions converts this conditional
convergence into absolute convergence of the blocks.
\end{theorem}

\begin{remark}
The isolated point $s=0$ does not enter the definition of the abscissa.
For completeness, every completed excursion satisfies
$D_I(0)=\sum_{n\in I}\mu(n)=0$.
\end{remark}

\begin{proof}
Let $s=\varsigma+i\tau$.  Since
\begin{equation}
 n^{-s}-(n+1)^{-s}
 =s\int_n^{n+1}x^{-s-1}\,dx,
\end{equation}
with the real positive branch of $\log x$, one has
\begin{equation}
 |n^{-s}-(n+1)^{-s}|\le |s|n^{-\varsigma-1}.
\end{equation}
Hence
\begin{equation}
 \sum_I|D_I(s)|
 \le |s|\sum_{n\ge1}|M(n)|n^{-\varsigma-1}.
\end{equation}
Let $K_0\Subset\{\Re s>\Theta\}$ be compact and set
$\varsigma_0=\min_{s\in K_0}\Re s>\Theta$.  Choose
$0<\eps_0<\varsigma_0-\Theta$.  From
$B_1(X)=X^{1+\Theta+o(1)}$ we have
$B_1(X)\ll_{\eps_0}X^{1+\Theta+\eps_0}$ for large $X$, and partial
summation gives
\begin{equation}
 \sum_{n\ge1}|M(n)|n^{-\varsigma_0-1}<\infty.
\end{equation}
Since $\sup_{s\in K_0}|s|<\infty$, the Weierstrass criterion gives normal
convergence of the block series on $K_0$.  Thus it is holomorphic in
$\Re s>\Theta$.  There are no zeros of $\zeta$ in that half-plane by the
definition of $\Theta$, and $1/\zeta$ extends holomorphically to $s=1$ with
value $0$; hence $1/\zeta$ is holomorphic throughout the connected
half-plane.  The identity already holds for $\Re s>1$, so the
identity theorem extends it to all $\Re s>\Theta$.

For real $\varsigma>0$, order the completed excursions by their right
endpoints.  If $V$ is a zero of $M$, the fixed sign of each excursion and
finite Abel summation give
\begin{equation}
 \sum_{I:\,x_1(I)\le V}|D_I(\varsigma)|
 =\sum_{n<V}|M(n)|
 \bigl(n^{-\varsigma}-(n+1)^{-\varsigma}\bigr).
\end{equation}
Since $M$ has arbitrarily large zeros (Definition~\ref{def:excursion},
using $M(x)=\Omega_\pm(\sqrt x)$), letting $V\to\infty$ through zeros and
using monotone convergence gives
\begin{equation}
\label{eq:block-absolute-real-axis}
 \sum_I|D_I(\varsigma)|
 =\sum_{n\ge1}|M(n)|
 \bigl(n^{-\varsigma}-(n+1)^{-\varsigma}\bigr),
\end{equation}
even when both sides are infinite.
If the right side were finite for some $0<\varsigma<\Theta$, then for
$n\le X$
\begin{equation}
 n^{-\varsigma}-(n+1)^{-\varsigma}
 =\varsigma\int_n^{n+1}t^{-\varsigma-1}\,dt
 \ge \varsigma(X+1)^{-\varsigma-1},
\end{equation}
so
\begin{equation}
 B_1(X)\le
 \frac{(X+1)^{\varsigma+1}}{\varsigma}
 \sum_{n\ge1}|M(n)|
 \bigl(n^{-\varsigma}-(n+1)^{-\varsigma}\bigr)
 =O(X^{1+\varsigma}),
\end{equation}
contradicting $B_1(X)=\Omega(X^{1+\Theta-\eps})$ for every $\eps>0$.
For real $\varsigma<0$ the same Abel identity holds (the factor
$n^{-\varsigma}-(n+1)^{-\varsigma}$ is negative, so each block again has
constant sign), and
$|n^{-\varsigma}-(n+1)^{-\varsigma}|=|\varsigma|\int_n^{n+1}x^{-\varsigma-1}dx
\ge|\varsigma|\min\{1,(V+1)^{-\varsigma-1}\}$ for $n<V$. Hence
$\sum_{I:\,x_1(I)\le V}|D_I(\varsigma)|\ge|\varsigma|\min\{1,(V+1)^{-\varsigma-1}\}
B_1(V-1)$, which tends to infinity along zeros $V$: for $\varsigma\le-1$ because
$B_1\to\infty$, and for $-1<\varsigma<0$ because
$V^{-\varsigma-1}B_1(V-1)\ge V^{\Theta-\varsigma-\eps}$ for large $V$.  Therefore the block abscissa is
exactly $\Theta$.
\end{proof}

\begin{remark}[Status of the result]
\label{rem:block-abscissa-status}
The analytic input is classical: $\sum_n|M(n)|n^{-\varsigma-1}<\infty$
exactly for $\varsigma>\Theta$ (upper bound $M(x)\ll x^{\Theta+\eps}$;
lower bound by the Mellin argument of Proposition~\ref{prop:L1}). The new
point is the excursion grouping: the fixed sign of each completed excursion
turns Abel summation into the positive identity
\eqref{eq:block-absolute-real-axis}. We did not find this block
formulation in the literature (absence of a match is weak evidence).
\end{remark}


\section{Lacunary excursion families and horizon-localized criteria}
\label{sec:lacunary}

\subsection{Geometric counts}
\label{subsec:geometric}

The lemmas of this subsection use only that excursions are disjoint integer
intervals ordered by their starting points, the next one starting at or
after the end of the previous one. They therefore hold verbatim for every
sequence of disjoint integer intervals, deterministic or random, and in
particular pathwise for any stochastic model of excursions
(Section~\ref{subsec:extreme}). For a family $\F$ of excursions and $X\ge2$
put
\begin{equation}
E_\F(X):=\#\{\nu\in\F:x_{0,\nu}\le X\}.
\end{equation}
For a function $\mathfrak h:[0,\infty)\to[0,\infty]$, the family of
excursions \emph{crossing the horizon} $\mathfrak h$ is
\begin{equation}
\F_{\mathfrak h}:=\{\nu:x_{1,\nu}>\mathfrak h(x_{0,\nu})\}.
\end{equation}
For $A>1$ the excursion $I_\nu$ is \emph{$A$-exceptional} if
$x_{1,\nu}>x_{0,\nu}^A$, i.e.\ $\nu\in\F_{x^A}$, and \emph{$A$-ordinary}
otherwise; for $A=2$ we say \emph{square-exceptional} and
\emph{square-ordinary}. For $\lambda>1$ it is \emph{$\lambda$-long} if
$x_{1,\nu}>\lambda x_{0,\nu}$. We write $E_A=E_{\F_{x^A}}$,
$E_\square=E_2$ and $E_\lambda=E_{\F_{\lambda x}}$; the context separates
$E_A$ ($A$ an exponent) from $E_\lambda$ ($\lambda$ a factor).

\begin{lemma}[Power lacunarity]
\label{lem:power-lacunary}
\tagNE\ For every $A>1$ and $X\ge2$,
\begin{equation}
\label{eq:power-lacunary-count}
E_A(X)\le2+\frac1{\log A}\log\frac{\log X}{\log2}.
\end{equation}
In particular
\begin{equation}
\label{eq:square-exception-count}
E_\square(X)\le2+\log_2\frac{\log X}{\log2},\qquad
E_\square(X)\le4\log\log X\quad(X\ge e^e).
\end{equation}
\end{lemma}

\begin{proof}
No excursion starts at $1$, since $M(1)=1$. List the $A$-exceptional starts
$\ge2$ as $2\le a_1<a_2<\cdots$. The next excursion starts at or after the
end of the preceding one, so $a_{j+1}\ge x_{1,j}>a_j^A$ and
$\log a_j\ge A^{j-1}\log2$. If $a_j\le X$ then
$j\le1+\log(\log X/\log2)/\log A$; add $1$ for the excursion $(0,2]$. For the
second bound in \eqref{eq:square-exception-count},
$\log_2(\log X/\log2)=(\log\log X-\log\log2)/\log2\le1.443\log\log X+0.529$.
\end{proof}

The next lemma is the partition of the square horizon
$[x_0,x_0^2]=[x_0,cx_0^2]\cup[cx_0^2,x_0^2]$: the upper part is accessible
only to a doubly lacunary family, for every fixed $c>0$.

\begin{lemma}[Fractional square horizon]
\label{lem:partition}
\tagNE\ For $0<c\le1$ and $X\ge2$,
\begin{equation}
\label{eq:partition-count}
E_{\F_{cx^2}}(X)=\#\{\nu:x_{0,\nu}\le X,\ x_{1,\nu}>cx_{0,\nu}^2\}
\le\frac2c+2+\log_2\frac{\log X}{\log2}.
\end{equation}
Consequently, the range $cx_0^2<t\le x_0^2$ of the excursions starting at
most $X$ is visited by at most $O_c(\log\log X)$ excursions.
\end{lemma}

\begin{proof}
At most $\lceil2/c\rceil\le2/c+1$ members start below $2/c$. List the others
as $a_1<a_2<\cdots$ and put $b_j=ca_j\ge2$. Then
$a_{j+1}\ge x_{1,j}>ca_j^2$, so
\begin{equation}
b_{j+1}=ca_{j+1}>c^2a_j^2=b_j^2,
\end{equation}
and $\log b_j\ge2^{j-1}\log2$. Since $b_j\le cX\le X$, the number of such
$j$ is at most $1+\log_2(\log X/\log2)$. An excursion containing a point
$t>cx_0^2$ has $x_1\ge t>cx_0^2$.
\end{proof}

\begin{lemma}[Superlinear horizons]
\label{lem:superlinear}
\tagNE\ Let $\psi:[0,\infty)\to(0,\infty)$ satisfy $\psi(x)\to\infty$. Then
\begin{equation}
E_{\F_{x\psi(x)}}(X)
=\#\{\nu:x_{0,\nu}\le X,\ x_{1,\nu}>x_{0,\nu}\psi(x_{0,\nu})\}
=o(\log X).
\end{equation}
\end{lemma}

\begin{proof}
Given $K>0$ choose $x_K$ with $\psi(x)\ge e^K$ for $x\ge x_K$. At most
$x_K+1$ members start below $x_K$; we may take $x_K\ge1$. For consecutive members
$a_j<a_{j+1}$ starting at or above $x_K$,
$a_{j+1}\ge x_{1,j}>a_j\psi(a_j)\ge e^Ka_j$, so $\log a_j$ increases by more
than $K$ at each step and there are at most $1+\log X/K$ of them. Hence
$\limsup E(X)/\log X\le1/K$ for every $K$.
\end{proof}

\begin{lemma}[Linear lacunarity]
\label{lem:linear-lacunary}
\tagNE\ For every $\lambda>1$ and $X\ge2$,
\begin{equation}
\label{eq:linear-lacunary-count}
E_\lambda(X)=\#\{\nu:x_{0,\nu}\le X,\ x_{1,\nu}>\lambda x_{0,\nu}\}
\le2+\frac{\log(X/2)}{\log\lambda}.
\end{equation}
\end{lemma}

\begin{proof}
For starts $a_j\ge2$, $a_{j+1}\ge x_{1,j}>\lambda a_j$, hence
$a_j>2\lambda^{j-1}$ for $j\ge2$; add $1$ for $(0,2]$.
\end{proof}

All four bounds are sharp up to constants for general families of disjoint
intervals, as geometric and doubly geometric chains of intervals show.

\subsection{Arithmetic input: Kaczorowski--Pintz}

We separate the published statement, its translation from $M_0=M+2$ to
$M$, and the consequences for excursions.

\begin{theorem}[Kaczorowski--Pintz {\cite{KP86}}]
\label{thm:KP-quoted}
\tagL\ Let $M_0(x)=M(x)-1/\zeta(0)=M(x)+2$, the function $M_a$ of \cite{KP86} at
$a=0$. Let $V(f,Y)$ be defined as in \cite{KP86}, through the largest number of
points of $[1,Y]$ at which $f$ is nonzero with alternating signs; the number
of sign changes of $f$ on $[1,Y]$ is at least $V(f,Y)-1$, which does not
affect the asymptotic statements below. Then:
\begin{enumerate}[label=\textup{(\roman*)}]
\item \cite[Corollary~4, p.~178, (2.13)--(2.15)]{KP86}
\begin{equation}
\label{eq:KP-signchanges}
\liminf_{Y\to\infty}\frac{V(M_0,Y)}{\log Y}\ge\frac{\gamma_1}{\pi},
\end{equation}
with combined oscillation of size $x^{\Theta-\eta}$ for every fixed
$\eta>0$: there exist ordered points $t_1<\cdots<t_N\le Y$, with
$N\ge(\gamma_1/\pi-o(1))\log Y$, whose signs alternate and which satisfy
$|M_0(t_j)|\ge t_j^{\Theta-\eta}$;
\item The local statement is the specialization to $M_0$ of
Theorem~1, (2.2)--(2.3), via Corollary~1 and Corollary~4 of \cite{KP86}.
Theorem~1 states literally that, for every $\varepsilon>0$ and all
sufficiently large $Y$, every interval $[Y^{1-\varepsilon},Y]$ contains at
least one sign change of $M_0$, and that the sign changes in (2.2) are
combined with oscillation of size $x^{\Theta-\varepsilon}$.  By the
definition of ``combined oscillation'' immediately following (1.1), the
alternating witness points satisfy the corresponding magnitude bound.
Consequently there are two points of that interval, of opposite signs, at
which $|M_0(t)|\ge t^{\Theta-\varepsilon}$.  In particular, for arbitrary fixed
$\delta>0$ and $\eta>0$, taking
$\varepsilon=\min\{\delta,\eta\}$ gives a threshold
$Y_0(\delta,\eta)$ such that $[Y^{1-\delta},Y]$ contains points $t_+,t_-$
with
$M_0(t_+)\ge t_+^{\Theta-\eta}$ and
$M_0(t_-)\le-t_-^{\Theta-\eta}$.
\end{enumerate}
\end{theorem}

\begin{lemma}[From $M_0$ to $M$]
\label{lem:KP-translation}
\tagNE\ If $|M_0(t)|>2$, then $\sgn M(t)=\sgn M_0(t)$ and
$|M(t)|\ge|M_0(t)|-2$. In particular, for every fixed $0<\eta<\Theta$, at points with
$|M_0(t)|\ge t^{\Theta-\eta}$ and $t$ large,
$\sgn M(t)=\sgn M_0(t)$ and $|M(t)|\ge\frac12t^{\Theta-\eta}$.
\end{lemma}

\begin{proof}
$M_0>2$ gives $M>0$; $M_0<-2$ gives $M<-4$. The second assertion uses
$t^{\Theta-\eta}\to\infty$ for $\eta<\Theta$.
\end{proof}

\begin{corollary}[Distinct large excursions]
\label{cor:KP-large-excursions}
\tagNE\ If $\Theta>1/2$ and $0<\eta<\Theta-1/2$, then
\begin{equation}
\label{eq:KP-large-excursions}
N_\eta(Y)\ge\left(\frac{\gamma_1}{\pi}-o(1)\right)\log Y,
\end{equation}
where $N_\eta(Y)$ counts distinct excursions starting at most $Y$ that
contain a point $t\le Y$ with $|M(t)|\ge\tfrac12 t^{\Theta-\eta}$.
\end{corollary}

\begin{proof}
Take the alternating points of Theorem~\ref{thm:KP-quoted}(i) in $[1,Y]$,
discarding the boundedly many below the threshold of
Lemma~\ref{lem:KP-translation}. By that lemma, consecutive points $t<t'$
have $M(t)$ and $M(t')$ of opposite signs. Since $M$ is integer valued with
steps $0,\pm1$, it vanishes at some integer in $(t,t')$, so $t$ and $t'$ lie
in different excursions. Each point with $M(t)\ne0$ belongs to exactly one
excursion, and excursions are ordered, so the excursions containing the
ordered points are pairwise distinct. Their starts are below the points,
hence at most $Y$, and $|M(t)|\ge\frac12t^{\Theta-\eta}$ at each point.
\end{proof}

\begin{corollary}[Multiplicative localization of Mertens returns]
\label{cor:KP-local-returns}
\tagC\ \textup{(Return localization, (RL).)} For every fixed $A>1$ there exists $U_A$ such that every excursion starting
at $u\ge U_A$ ends at $v\le u^A$. Equivalently, $M$ has only finitely many
$A$-exceptional excursions; in particular only finitely many
square-exceptional ones.
\end{corollary}

\begin{proof}
Put $\delta=1-1/A$ and $Y=u^A$, so that
$[Y^{1-\delta},Y]=[u,u^A]$.  Fix once and for all
$\eta=\Theta/2$.  After this choice the resulting threshold depends only
on $A$; denote it by $U_A$.  Enlarge $U_A$ if necessary so that
$t^{\Theta-\eta}>2$ throughout $[u,u^A]$ whenever $u\ge U_A$.
Theorem~\ref{thm:KP-quoted}(ii) and Lemma~\ref{lem:KP-translation} then give
$t_+,t_-\in[u,u^A]$ with $M(t_+)>0>M(t_-)$ and
$|M_0(t_\pm)|>2$.  Neither point can lie in $[u,u+1)$: on that interval
$M=M(u)=0$ and hence $M_0=2$.  Therefore $u<\lfloor t_\pm\rfloor$. Since $M$ changes by
$0,\pm1$ at each integer, it vanishes at some integer $n$ with
$\min\{\lfloor t_+\rfloor,\lfloor t_-\rfloor\}<n<\max\{\lfloor t_+\rfloor,\lfloor t_-\rfloor\}$, and $u<n<u^A$. If
$u=x_{0,\nu}$, the first return satisfies $x_{1,\nu}\le n<u^A$.
\end{proof}

\begin{remark}
A bare sign change of $M_0=M+2$ would not force a zero of $M$. The
decisive point is that the local sign change is combined with oscillation of
size $t^{\Theta-\eta}\to\infty$, so that Lemma~\ref{lem:KP-translation}
applies to both witnesses.
\end{remark}

\begin{remark}[Where (RL) is used]
\label{rem:RL-dependence}
Property (RL) rests on the local statement (2.2)--(2.3) of
Theorem~1 of \cite{KP86}, inherited by Corollary~1 and applied to $M_0$ in
Corollary~4: every sufficiently late interval $[Y^{1-\varepsilon},Y]$
contains a sign change combined with oscillation of size
$x^{\Theta-\varepsilon}$.  A density of sign changes alone would not bound
individual return times. It is used exactly in:
the terminal-block step of Corollary~\ref{cor:local-corr-sufficient}; the
lower bound of Corollary~\ref{cor:area-exponent}; the converse direction of
Theorem~\ref{thm:weighted-corr-RH}; Corollary~\ref{cor:KP-mean-log-length};
Proposition~\ref{prop:positional-surplus}; Remark~\ref{rem:onedirection};
and the statements that power and fractional-square horizons are inactive
for $M$ (Remark~\ref{rem:geom-vs-arith}). All other results are independent
of (RL); in particular the horizon theorem uses only
Theorem~\ref{thm:KP-quoted}(i).
\end{remark}

\begin{remark}[Geometric counts and arithmetic localization]
\label{rem:geom-vs-arith}
For the actual Mertens excursions Corollary~\ref{cor:KP-local-returns} is
stronger than Lemma~\ref{lem:power-lacunary}: each $A$-exceptional family is
finite. Consequently, for $M$ itself the power horizons $x_0^A$ and the
fractional-square horizons $cx_0^2$ are inactive for all late excursions,
and the corresponding RH-equivalences (Corollary~\ref{cor:power-ordinary-RH},
Corollary~\ref{cor:descending}(i)) follow already from
Corollary~\ref{cor:KP-local-returns}; the horizons with genuine content for
$M$ are $x_0\psi(x_0)$ and $\lambda x_0$. The geometric lemmas remain separate results because they are
elementary, independent of $\zeta$, and valid for every ordered family of
disjoint intervals. This is exactly what the localization theorem below and
the extreme-value theorems of Section~\ref{subsec:extreme} use: the former
needs only a count below $(\gamma_1/\pi)\log X$, the latter needs a count
valid pathwise in a model where no analogue of Kaczorowski--Pintz is
available.
\end{remark}

\subsection{Horizon localization}
\label{subsec:horizon}

\begin{theorem}[Horizon-localized RH criterion]
\label{thm:horizon}
\tagN\ Let $\mathfrak h:[0,\infty)\to[0,\infty]$ be a function, let $\F$ be a family
of excursions, and suppose that for some $0\le\kappa<\gamma_1/\pi$
\begin{equation}
\label{eq:horizon-hyp}
E_{\F\cup\F_{\mathfrak h}}(X)\le(\kappa+o(1))\log X\qquad(X\to\infty).
\end{equation}
Then RH is equivalent to the following statement: for every $\eps>0$ there
is $C_\eps$ such that
\begin{equation}
\label{eq:restricted-RH}
|M(t)|\le C_\eps t^{1/2+\eps}
\qquad\text{for all }t\in I_\nu,\ \nu\notin\F,\ t\le \mathfrak h(x_{0,\nu}).
\end{equation}
No condition is imposed on the excursions in $\F$, nor on the part of any
excursion lying beyond its horizon $\mathfrak h(x_0)$.
\end{theorem}

\begin{proof}
RH implies \eqref{eq:restricted-RH} by Littlewood's criterion. Conversely,
suppose $\Theta>1/2$ and fix $\eta,\eps>0$ with
$1/2+\eps<\Theta-\eta$. By \eqref{eq:KP-large-excursions} there are at least
$(\gamma_1/\pi-o(1))\log X$ distinct excursions starting at most $X$, each
containing a witness $t$ with $|M(t)|\ge\frac12t^{\Theta-\eta}$. Consider
such an excursion $\nu$ with witness $t$. If $\nu\notin\F$ and
$t\le \mathfrak h(x_{0,\nu})$, the witness lies in the range of
\eqref{eq:restricted-RH}. Otherwise $\nu\in\F$, or $t>\mathfrak h(x_{0,\nu})$, in
which case $x_{1,\nu}>t>\mathfrak h(x_{0,\nu})$ (as $M(t)\ne0$) and $\nu\in\F_{\mathfrak h}$. By
\eqref{eq:horizon-hyp} at least
\begin{equation}
\Bigl(\frac{\gamma_1}{\pi}-\kappa-o(1)\Bigr)\log X\longrightarrow\infty
\end{equation}
witnesses lie in the range of \eqref{eq:restricted-RH}. They belong to
distinct excursions, so they are unbounded, and at them
$t^{\Theta-\eta}/t^{1/2+\eps}\to\infty$, contradicting
\eqref{eq:restricted-RH}.
\end{proof}

The case $\mathfrak h(x)=\infty$ excludes a sparse family $\F$ of whole excursions;
the case $\F=\emptyset$ truncates every excursion at its horizon.

\begin{corollary}[Square horizon]
\label{cor:power-ordinary-RH}
\tagN\ For every fixed $A>1$, RH is equivalent to \eqref{eq:restricted-RH}
restricted to $A$-ordinary excursions. In particular
\begin{equation}
\label{eq:square-horizon-RH}
\boxed{\ \mathrm{RH}\iff
|M(t)|\ll_\eps t^{1/2+\eps}
\ \text{for all $t$ in excursions with }x_1\le x_0^2.\ }
\end{equation}
A separate application of Theorem~\ref{thm:horizon}, with no excursions
discarded and horizon $\mathfrak h(x)=x^2$, gives the second RH-equivalent
statement
\begin{equation}
 |M(t)|\ll_\eps t^{1/2+\eps}
 \quad(x_{0,\nu}<t\le\min\{x_{1,\nu},x_{0,\nu}^2\})
\end{equation}
over all excursions.  The two restricted statements are therefore two
distinct consequences of the horizon theorem, not an elementary equivalence
between the restrictions themselves.
\end{corollary}

\begin{proof}
Take $\F=\F_{x^A}$, $\mathfrak h=\infty$, respectively $\F=\emptyset$, $\mathfrak h(x)=x^2$, in
Theorem~\ref{thm:horizon}; Lemma~\ref{lem:power-lacunary} gives
$\kappa=0$.
\end{proof}

The count $E_\square(X)=O(\log\log X)$ is thus what makes the square
horizon admissible. The square horizon can be lowered.

\begin{corollary}[Descending horizons]
\label{cor:descending}
\tagN\ RH is equivalent to Littlewood's bound
$|M(t)|\ll_\eps t^{1/2+\eps}$ on the ranges
\begin{equation}
x_{0,\nu}<t\le\min\{x_{1,\nu},\mathfrak h(x_{0,\nu})\}
\end{equation}
over all excursions, for each of the following horizons:
\begin{enumerate}[label=\textup{(\roman*)}]
\item $\mathfrak h(x)=cx^2$, for any fixed $c>0$;
\item $\mathfrak h(x)=x\psi(x)$, for any $\psi:[0,\infty)\to(0,\infty)$ with $\psi(x)\to\infty$;
\item $\mathfrak h(x)=\lambda x$, for any fixed $\lambda>\lambda_*:=e^{\pi/\gamma_1}=1.2489\ldots$
\end{enumerate}
Here $c$, $\psi$ and $\lambda$ do not depend on the excursion.
The same holds with ``truncate at $\mathfrak h(x_0)$'' replaced by ``discard the
excursions with $x_1>\mathfrak h(x_0)$''. In particular it is enough to control the
excursions with $x_1\le\frac54x_0$.
\end{corollary}

\begin{proof}
Apply Theorem~\ref{thm:horizon} with $\F=\emptyset$ (respectively
$\F=\F_{\mathfrak h}$, $\mathfrak h=\infty$). For (i), $c\ge1$ reduces to
Lemma~\ref{lem:power-lacunary} and $0<c<1$ to Lemma~\ref{lem:partition};
$\kappa=0$. For (ii), Lemma~\ref{lem:superlinear} gives $\kappa=0$. For
(iii), Lemma~\ref{lem:linear-lacunary} gives $\kappa=1/\log\lambda$, and
$1/\log\lambda<\gamma_1/\pi$ exactly when $\lambda>\lambda_*$. Finally
$5/4>\lambda_*$.
\end{proof}

At the level of the restricted statements, (i) and (ii) imply (iii)
eventually, because their horizons are at least $\lambda x_0$ for every
fixed admissible $\lambda>\lambda_*$ once $x_0$ is large.  The converse
does not follow from inclusion of ranges; rather, each of (i)--(iii) is
separately RH-equivalent by Theorem~\ref{thm:horizon}.  The finitely many
small excursions affect only the constant $C_\eps$. The horizon has
been lowered from $x_0^2$ to $\lambda x_0$: the possible failure of
Littlewood's bound, if any, is visible within a bounded multiplicative
distance of the start of an excursion. Quantitatively:

\begin{corollary}[Two-sided localization of large values]
\label{cor:two-sided}
\tagN\ Suppose $\Theta>1/2$, and fix $0<\eta<\Theta-\frac12$ and
$\lambda>\lambda_*$ before letting $X\to\infty$.
Among the excursions starting at $x_0\le X$, at least
\begin{equation}
\Bigl(\frac{\gamma_1}{\pi}-\frac1{\log\lambda}-o(1)\Bigr)\log X
\end{equation}
contain a point $t_*$ with $|M(t_*)|\ge\frac12t_*^{\Theta-\eta}$ and
\begin{equation}
x_0+\tfrac12x_0^{\Theta-\eta}\le t_*\le\lambda x_0 .
\end{equation}
\end{corollary}

\begin{proof}
Of the excursions of Corollary~\ref{cor:KP-large-excursions}, those whose
witness exceeds $\lambda x_0$ are $\lambda$-long, hence at most
$\log X/\log\lambda+2$ by Lemma~\ref{lem:linear-lacunary}. For the others,
$M(x_0)=0$ and unit steps give $|M(t_*)|\le t_*-x_0$, so
$t_*-x_0\ge\frac12t_*^{\Theta-\eta}\ge\frac12x_0^{\Theta-\eta}$.
\end{proof}


\begin{remark}[Dyadic erosion and the uniform transition]
\label{rem:dyadic-erosion}
Fix an integer $j\ge0$ before letting $X\to\infty$.  Lemma~\ref{lem:partition} with
$c=2^{-j}$ gives
\begin{equation}
\#\left\{\nu:x_{0,\nu}\le X,\ x_{1,\nu}>
\frac{x_{0,\nu}^2}{2^j}\right\}=O_{2^j}(\log\log X),
\end{equation}
so Theorem~\ref{thm:horizon} successively permits the horizons
\begin{equation}
x_0^2,\quad \frac{x_0^2}{2},\quad \frac{x_0^2}{4},\quad
\frac{x_0^2}{8},\ldots.
\end{equation}
The constant in the fixed-$c$ estimate depends on $c$; hence one may not set
$j=j(x_0)\sim\log_2x_0$ by that lemma alone.  If instead
$2^{j(x)}=x/\psi(x)$ with $\psi(x)\to\infty$, then the corresponding
horizon is $x\psi(x)$ and Lemma~\ref{lem:superlinear} supplies the uniform
$o(\log X)$ crossing count.  Thus the rigorous erosion chain is
\begin{equation}
x_0^2\to\frac{x_0^2}{2}\to\frac{x_0^2}{4}\to\cdots
\to x_0\psi(x_0)\to\lambda x_0,\qquad\lambda>\lambda_*.
\end{equation}
\end{remark}

\subsection{Fixed-mask Catalan compression inside fractional square horizons}
\label{subsec:catalan-square}

The fractional-square localization has a useful interaction with the exact
combinatorics of a completed excursion.  The relevant Catalan time is not
physical time: zero M\"obius increments must first be removed.  This
support-time distinction is essential below.

\begin{definition}[Support-time path and fixed-mask Catalan law]
\label{def:support-catalan}
Let $I=(x_0,x_1]$ be a complete excursion, and list the nonzero increments in
$I$ as
\begin{equation}
 x_0<a_1<a_2<\cdots<a_L\le x_1,
 \qquad \mu(a_j)\ne0.
\end{equation}
Put
\begin{equation}
 r_I(t):=\#\{j:a_j\le t\},\qquad x_0<t\le x_1.
\end{equation}
For the excursion $I_\nu$, let $\sigma_\nu\in\{\pm1\}$ be its orientation and define
\begin{equation}
 S_0=0,\qquad
 S_r:=\sigma_\nu\sum_{j=1}^r\mu(a_j)\qquad(1\le r\le L).
\end{equation}
Then $L=2n$ is even and
\begin{equation}
 S_0=S_L=0,\qquad S_r>0\quad(1\le r<L),\qquad
 S_r-S_{r-1}\in\{-1,+1\}.
\end{equation}
Thus $(S_r)_{0\le r\le L}$ is a primitive Dyck path.  The
\emph{fixed-mask primitive Catalan law} keeps the physical support
$\{a_1,\dots,a_L\}$ fixed and assigns to those support positions a uniformly
chosen primitive Dyck sign sequence of length $L$.  We denote its probability
by $\mathbb P_{\mathrm{pCat},I}$.
\end{definition}

For the observed excursion one has, exactly,
\begin{equation}
\label{eq:support-time-height}
 |M(t)|=S_{r_I(t)},
 \qquad
 r_I(t)\le t-x_0\le t,
\end{equation}
with $r_I(t)<t$ whenever $x_0\ge1$, i.e.\ for every excursion other than $(0,2]$.
The first equality uses the orientation of the excursion and the fact that
$M$ is constant across the zero increments.

We first record the global Catalan estimate, since it is useful in its own
right and clarifies why a different estimate is needed at the square
horizon.

\begin{lemma}[Catalan maximum tail]
\label{lem:catalan-max-tail}
\tagL\ Let $L=2n$ and let $H_{\max}$ be the maximum of a uniformly chosen Dyck path
of length $2n$.  For $1\le h\le n$,
\begin{equation}
\label{eq:catalan-max-tail}
 \mathbb P_{\mathrm{Cat},L}\{H_{\max}\ge h\}
 \le (n+1)\frac{\binom{2n}{n+h}}{\binom{2n}{n}}
 \le (n+1)\exp\!\left(-\frac{h^2}{n+h}\right)
 \le \left(\frac L2+1\right)e^{-h^2/L}.
\end{equation}
For a uniformly chosen primitive Dyck path of length $L=2n$,
\begin{equation}
\label{eq:primitive-catalan-max-tail}
 \mathbb P_{\mathrm{pCat},L}\{H_{\max}\ge h\}
 \le \frac L2\exp\!\left(-\frac{(h-1)^2}{L-2}\right),
 \qquad 2\le h\le n.
\end{equation}
\end{lemma}

\begin{proof}
There are $\mathrm{Cat}_n=(n+1)^{-1}\binom{2n}{n}$ Dyck paths of length $2n$ and
$\binom{2n}{n}$ unrestricted bridges.  Reflecting a bridge after its first
visit to level $h$ gives a bijection between bridges that reach $h$ and
unrestricted walks ending at $2h$, hence their number is
$\binom{2n}{n+h}$.  Therefore
\begin{equation}
 \mathbb P_{\mathrm{Cat},2n}\{H_{\max}\ge h\}
 \le (n+1)\frac{\binom{2n}{n+h}}{\binom{2n}{n}}.
\end{equation}
Moreover
\begin{equation}
 \frac{\binom{2n}{n+h}}{\binom{2n}{n}}
 =\prod_{j=1}^h\frac{n-j+1}{n+j},
\end{equation}
and $\log(1-x)\le-x$ gives
\begin{equation}
 \log\frac{\binom{2n}{n+h}}{\binom{2n}{n}}
 \le-\sum_{j=1}^h\frac{2j-1}{n+j}
 \le-\frac{h^2}{n+h}.
\end{equation}
Since $n+h\le2n=L$, \eqref{eq:catalan-max-tail} follows.  A primitive path
of length $2n$ has the form $+1$, followed by a Dyck path of length $2n-2$,
followed by $-1$; its maximum is one plus the maximum of the interior Dyck
path.  Applying the preceding bound with $n-1$ and $h-1$ proves
\eqref{eq:primitive-catalan-max-tail}.
\end{proof}

\begin{remark}[Why the maximum tail is not the square-horizon estimate]
\label{rem:max-tail-square-limit}
The exponent in \eqref{eq:primitive-catalan-max-tail} is of order
$H^2/L$.  For a general family of disjoint intervals retained by the
square horizon one may have $L\asymp x_0^2$ (for $M$ itself such excursions
are finitely many by Corollary~\ref{cor:KP-local-returns}), while the relevant witness lies much earlier, at physical
time $t\asymp x_0$.  Then a violation of size $t^{1/2+\eps}$ need not make
$H^2/L$ large.  Thus the global maximum tail, although correct, is not the
right Catalan quantity for combining with the fractional-square horizon.
The correct quantity is the pointwise height at the support time
$r_I(t)$.
\end{remark}

\begin{lemma}[Pointwise primitive-Catalan tail]
\label{lem:pointwise-catalan-tail}
\tagNE\ Let $L=2n\ge4$ and let $(S_r)$ be a uniformly chosen primitive Dyck path of
length $L$.  For $1\le r\le L-1$ and real $h\ge1$,
\begin{equation}
\label{eq:pointwise-catalan-tail-left}
 \mathbb P_{\mathrm{pCat},L}\{S_r\ge h\}
 \le 8n^{3/2}\exp\!\left(-\frac{h^2}{2r}\right)
 \le 3L^{3/2}\exp\!\left(-\frac{h^2}{2r}\right).
\end{equation}
By time reversal the same estimate holds with $r$ replaced by $L-r$;
hence
\begin{equation}
\label{eq:pointwise-catalan-tail-two-sided}
 \mathbb P_{\mathrm{pCat},L}\{S_r\ge h\}
 \le 3L^{3/2}
 \exp\!\left[-\frac{h^2}{2\min\{r,L-r\}}\right].
\end{equation}
If $h>\min\{r,L-r\}$ the probability is $0$, so the bound is only
relevant for $h\le\min\{r,L-r\}$.
\end{lemma}

\begin{proof}
Ignore both positivity and the endpoint constraint.  Among all $2^L$ sign
words, the proportion whose first $r$ signs have sum at least $h$ is at most
$\exp(-h^2/(2r))$ by the elementary Chernoff bound for a Rademacher sum.
Therefore the number of primitive Dyck paths satisfying $S_r\ge h$ is at
most
\begin{equation}
 2^L\exp\!\left(-\frac{h^2}{2r}\right).
\end{equation}
The total number of primitive Dyck paths of length $2n$ is $\mathrm{Cat}_{n-1}$.  The
standard central-binomial lower bound
$\binom{2m}{m}\ge4^m/(2\sqrt m)$ gives
\begin{equation}
 \mathrm{Cat}_{n-1}=\frac1n\binom{2n-2}{n-1}
 \ge\frac{4^{n-1}}{2n\sqrt{n-1}}.
\end{equation}
Since $2^L=4^n$, division gives the first inequality in
\eqref{eq:pointwise-catalan-tail-left}; the second follows from
$8n^{3/2}<3(2n)^{3/2}$.  Reversing the path and changing the signs of the
reversed increments preserves the primitive Dyck class and carries height
$S_r$ to height $S_{L-r}$, which proves
\eqref{eq:pointwise-catalan-tail-two-sided}.
\end{proof}

For an observed excursion $I$ define its pointwise Catalan reference-tail weight
by
\begin{equation}
\label{eq:pointwise-catalan-pvalue}
 p_{\mathrm{Cat}}(I,t):=
 \mathbb P_{\mathrm{pCat},I}
 \left\{S_{r_I(t)}\ge |M(t)|\right\}.
\end{equation}
This is a reference-tail probability under the fixed-mask Catalan law; it is
not a $p$-value for a stochastic model of the arithmetic M\"obius path.

\begin{corollary}[Fractional-square Catalan rarity]
\label{cor:fractional-square-catalan-rarity}
\tagNE\ Fix $c>0$ and $\eps>0$.  Let $I=(x_0,x_1]$ be a complete excursion with
$x_1\le cx_0^2$, and keep its zero mask fixed.  For all sufficiently large
$x_0$,
\begin{equation}
\label{eq:fractional-square-catalan-union}
 \mathbb P_{\mathrm{pCat},I}
 \left\{\exists\,t\in I\cap\mathbb Z:\ S_{r_I(t)}\ge \lceil t^{1/2+\eps}\rceil\right\}
 \le 3c^{5/2}x_0^5
 \exp\!\left(-\frac12x_0^{2\eps}\right).
\end{equation}
In particular the right-hand side is smaller than $x_0^{-A}$ for every
fixed $A>0$ once $x_0$ is large enough.
\end{corollary}

\begin{proof}
If $L=L^\#(I)=2$ then $S_r\le1$ and the event is empty for $x_0\ge1$; assume
$L\ge4$. At a support position $t=a_r$ one has $r\le t-x_0<t$.  Hence
Lemma~\ref{lem:pointwise-catalan-tail} with $h=\lceil t^{1/2+\eps}\rceil$ gives
\begin{equation}
 \mathbb P_{\mathrm{pCat},I}\{S_r\ge t^{1/2+\eps}\}
 \le3L^{3/2}\exp\!\left(-\frac12t^{2\eps}\right)
 \le3L^{3/2}\exp\!\left(-\frac12x_0^{2\eps}\right).
\end{equation}
There are at most $L$ support positions.  Since
$L\le x_1-x_0\le x_1\le cx_0^2$, one has
$3L\cdot L^{3/2}=3L^{5/2}\le3(cx_0^2)^{5/2}=3c^{5/2}x_0^5$, and the union
bound gives
\eqref{eq:fractional-square-catalan-union}.  Between support positions the
path is constant, so no additional physical times have to be counted.
\end{proof}

The preceding statement is purely combinatorial: it measures how far a large
value lies in the tail of a reference law, not how the M\"obius constraints
select a path. Combined with the localization theorem it gives the following
quantitative form.

\begin{corollary}[Catalan cost of a Kaczorowski--Pintz witness under $\Theta>1/2$]
\label{cor:catalan-surprise}
\tagN\ Suppose $\Theta>1/2$, fix $0<\eta<\Theta-\frac12$, put
$\delta=\Theta-\eta-\frac12$, and fix $\lambda>\lambda_*$. Then there are
infinitely many excursions $I=(x_0,x_1]$ with $x_1\le\lambda x_0$ containing
a point $t$ with
\begin{equation}
\label{eq:KP-catalan-surprise}
-\log p_{\mathrm{Cat}}(I,t)\ \ge\ \frac18x_0^{2\delta}-O(\log x_0).
\end{equation}
The same holds with $x_1\le\lambda x_0$ replaced by $x_1\le cx_0^2$ for any
fixed $c>0$.
\end{corollary}

\begin{proof}
By Corollary~\ref{cor:KP-large-excursions} and Lemma~\ref{lem:linear-lacunary},
at least $(\gamma_1/\pi-1/\log\lambda-o(1))\log X$ excursions starting at
most $X$ have $x_1\le\lambda x_0$ and contain $t$ with
$h=|M(t)|\ge\frac12t^{1/2+\delta}$; for $cx_0^2$ use
Lemma~\ref{lem:partition} instead. With $r=r_I(t)<t$ and $L\le x_1-x_0$,
Lemma~\ref{lem:pointwise-catalan-tail} gives
$p_{\mathrm{Cat}}(I,t)\le3L^{3/2}\exp(-h^2/(2r))\le3L^{3/2}\exp(-t^{2\delta}/8)$,
and $t>x_0$, $L\le\lambda x_0$ (respectively $L\le cx_0^2$).
\end{proof}

\begin{remark}[Selector bounds are one-sided]
\label{rem:catalan-selector-gap}
Lemma~\ref{lem:pointwise-catalan-tail} gives only the upper-tail estimate
\begin{equation}
 p_{\mathrm{Cat}}(I,t)
 \le 3L^{3/2}\exp\!\left(-\frac{|M(t)|^2}{2r_I(t)}\right),
\end{equation}
and hence the one-sided implication
\begin{equation}
 -\log p_{\mathrm{Cat}}(I,t)
 \ge \frac{|M(t)|^2}{2r_I(t)}-\log(3L^{3/2}).
\end{equation}
No reverse inequality, and therefore no equivalence ``up to polynomial
factors'', is proved here.  Consequently a selector bound
$-\log p_{\mathrm{Cat}}(I,t)\le\Phi(x_0)$ with
$\Phi(x)=x^{o(1)}$ on excursions with $x_1\le\lambda x_0$ implies
$|M(t)|\le t^{1/2}x_0^{o(1)}$ there, so RH follows from
Corollary~\ref{cor:descending}; this is a sufficient condition, not an
independent equivalent formulation.  A polynomial lower bound
$p_{\mathrm{Cat}}\ge x_0^{-A}$ would similarly give
$|M(t)|\ll\sqrt{t\log t}$ on the retained ranges.  To turn the Catalan
weight into a two-sided measure of deviation one would need a matching
lower-tail estimate of the form
$-\log p_{\mathrm{Cat}}(I,t)\ll |M(t)|^2/r_I(t)+\log L$, which is not proved
here.
\end{remark}

\begin{corollary}[Catalan cost along a Littlewood-violating sequence]
\label{cor:catalan-sequence}
\tagNE\ Let $t_j\to\infty$, $A_j\to\infty$ and $\eps_0>0$ with
$|M(t_j)|\ge A_jt_j^{1/2+\eps_0}$, and let $I_j=(x_{0,j},x_{1,j}]$ be the
excursion containing $t_j$. If $x_{1,j}\le x_{0,j}^{B}$ for a fixed $B$,
then
\begin{equation}
-\log p_{\mathrm{Cat}}(I_j,t_j)\ \ge\ \frac{A_j^2}{2}\,t_j^{2\eps_0}
-O_B(\log x_{0,j})\ \ge\ \frac{A_j^2}{2}\,x_{0,j}^{2\eps_0}-O_B(\log x_{0,j}).
\end{equation}
\end{corollary}

\begin{proof}
With $r=r_{I_j}(t_j)<t_j$ and $h=|M(t_j)|$,
$h^2/(2r)\ge A_j^2t_j^{1+2\eps_0}/(2t_j)$; apply
Lemma~\ref{lem:pointwise-catalan-tail} with $L\le x_{1,j}\le x_{0,j}^B$.
\end{proof}

\subsubsection{Selection energies: complete excursions and open prefixes}
\label{subsubsec:selection-energies}

The support-time clock of Definition~\ref{def:support-catalan} is the
squarefree-compressed clock
\begin{equation}
\Sq(x)=\#\{n\le x:\mu(n)\ne0\},\qquad L^\#(I)=\Sq(x_1)-\Sq(x_0)\le x_1-x_0,
\end{equation}
and $r_I(t)=\Sq(t)-\Sq(x_0)$. Two tail quantities are natural: one for the
height of a complete excursion, and one for an open prefix, which does not
require knowledge of the future return.

\begin{definition}[Selection energies]
\label{def:selection-energies}
For a complete excursion $I$ with height $H(I)$ and compressed length
$L^\#=L^\#(I)$, put
\begin{equation}
p^{\max}_{\mathrm{Cat}}(I):=\mathbb P_{\mathrm{pCat},L^\#}\{H_{\max}\ge H(I)\},
\qquad
\mathcal I_{\mathrm{Cat}}(I):=-\log p^{\max}_{\mathrm{Cat}}(I).
\end{equation}
For $t\ge1$ with $M(t)\ne0$ let $a(t)$ be the last zero of $M$ before $t$ and
$r=\Sq(t)-\Sq(a(t))$. After orientation, the compressed segment from $a(t)$ to
$t$ is a positive meander of length $r$: a $\pm1$ path $(S_k)_{k\le r}$ with
$S_0=0$ and $S_k>0$ for $1\le k\le r$. Let $\mathcal M_r$ be the set of such
paths.  Since $S_r\equiv r\pmod 2$, an endpoint event with incompatible
parity has probability zero.  Define
\begin{equation}
p_{\mathrm{open}}(t):=\mathbb P_{\mathcal M_r}\{S_r\ge|M(t)|\},
\end{equation}
the uniform measure on $\mathcal M_r$.
\end{definition}

\begin{lemma}[Height--entropy inequalities]
\label{lem:height-entropy}
\tagNE\ \begin{enumerate}[label=\textup{(\roman*)}]
\item If $L^\#=2$, the compressed path is $0,1,0$ and $H(I)=1$. For every
complete excursion with $L^\#\ge4$,
\begin{equation}
\mathcal I_{\mathrm{Cat}}(I)\ \ge\ \frac{(H(I)-1)^2}{L^\#-2}-\log\frac{L^\#}2 .
\end{equation}
\item $\#\mathcal M_r=\binom{r-1}{\lfloor(r-1)/2\rfloor}\ge2^r/(4\sqrt r)$,
and for $h\ge1$
\begin{equation}
\mathbb P_{\mathcal M_r}\{S_r\ge h\}\le4\sqrt r\,e^{-h^2/(2r)},
\qquad
-\log p_{\mathrm{open}}(t)\ \ge\ \frac{M(t)^2}{2r}-\log\bigl(4\sqrt r\bigr).
\end{equation}
\end{enumerate}
\end{lemma}

\begin{proof}
(i) is \eqref{eq:primitive-catalan-max-tail}. (ii) The count of positive
meanders is classical (ballot theorem; the number ending at $s$ is
$\frac sr\binom{r}{(r+s)/2}$). From
$\binom{2m}{m}\ge4^m/(2\sqrt m)$ and
$\binom{2m+1}{m}=\frac12\binom{2m+2}{m+1}$ one gets
$\#\mathcal M_r\ge2^r/(4\sqrt r)$. Among all $2^r$ words the proportion with
$S_r\ge h$ is at most $e^{-h^2/(2r)}$ by the Chernoff bound; divide.
\end{proof}

\begin{proposition}[Selector criteria]
\label{prop:selector-criteria}
\tagN\ \begin{enumerate}[label=\textup{(\roman*)}]
\item \emph{Complete excursions.} Fix $\lambda>\lambda_*$. If there is $A$
such that $p^{\max}_{\mathrm{Cat}}(I)\ge x_0^{-A}$ for every sufficiently
large complete excursion with $x_1\le\lambda x_0$, then
$H(I)\ll_{A,\lambda}\sqrt{x_0\log x_0}$ on that family, and RH holds.
\item \emph{Open prefixes.} If there is $A$ such that
$p_{\mathrm{open}}(t)\ge t^{-A}$ for all large $t$ with $M(t)\ne0$, then
$|M(t)|\ll_A\sqrt{t\log t}$, and RH holds.
\end{enumerate}
\end{proposition}

\begin{proof}
(i) If $L^\#=2$ then $H(I)=1$. Otherwise $L^\#\ge4$ and, by
Lemma~\ref{lem:height-entropy}(i), $(H-1)^2\le(L^\#-2)(A\log x_0+\log L^\#)$, and $L^\#\le(\lambda-1)x_0$.
Littlewood's bound therefore holds on every excursion with
$x_1\le\lambda x_0$, and Corollary~\ref{cor:descending}(iii), in the form
that discards the $\lambda$-long excursions, gives RH.
(ii) By Lemma~\ref{lem:height-entropy}(ii),
$M(t)^2\le2r(A\log t+\log(4\sqrt r))$ with $r\le t$.
\end{proof}





\begin{remark}[The threshold $\lambda_*$]
\label{rem:lambda-star}
The function $\cos(\gamma_1\log x)$ changes sign $\gamma_1/\pi$ times per
unit of $\log x$, i.e.\ once in every multiplicative interval of ratio
$e^{\pi/\gamma_1}=\lambda_*$; the constant in \eqref{eq:KP-signchanges} is
this density. On the geometric side, a chain of $\lambda$-long intervals
realizes the bound of Lemma~\ref{lem:linear-lacunary}. The method of
Theorem~\ref{thm:horizon} compares these two counts, and $\lambda_*$ is the
exact balance point, hence the sharp threshold of this counting mechanism:
for $\lambda\le\lambda_*$ the $\lambda$-long
family can, as far as counting alone is concerned, absorb every
Kaczorowski--Pintz excursion. Going below $\lambda_*$ requires information
on where inside an excursion the large values occur, not only how many
excursions carry them.
\end{remark}

\begin{proposition}[Logarithmic-length budget]
\label{prop:log-budget}
\tagNE\ Let $\F$ be a family of excursions and let $u_1<\dots<u_N$ be the starts of
its members with $2\le u_j\le X$, with ends $v_j$.  The initial excursion
$(0,2]$ is explicitly excluded, since $\log(v/u)$ is not defined at
$u=0$.  Then
\begin{equation}
\sum_{j=1}^{N-1}\log\frac{v_j}{u_j}\le\log\frac X2 .
\end{equation}
If moreover $v_N\le u_N^A$ for some $A\ge1$, then
$\sum_{j=1}^{N}\log(v_j/u_j)\le A\log X$.
\end{proposition}

\begin{proof}
Disjointness gives $u_{j+1}\ge v_j$, so
$\log(v_j/u_j)\le\log u_{j+1}-\log u_j$; the sum telescopes to at most
$\log u_N-\log u_1\le\log(X/2)$. The last term is at most
$(A-1)\log u_N\le(A-1)\log X$.
\end{proof}

Lemmas~\ref{lem:power-lacunary}--\ref{lem:linear-lacunary} are the special
cases in which every term is bounded below. The budget form shows what the
horizon theorem actually measures.

\begin{corollary}[Mean logarithmic length of Kaczorowski--Pintz excursions]
\label{cor:KP-mean-log-length}
\tagN\ Suppose $\Theta>1/2$, fix $0<\eta<\Theta-\frac12$, and let $\mathcal K(X)$
be any set of distinct excursions starting at most $X$, each containing a
point $t$ with $|M(t)|\ge\frac12t^{\Theta-\eta}$, such that
$\#\mathcal K(X)\ge(\gamma_1/\pi-o(1))\log X$ (such sets exist by
Corollary~\ref{cor:KP-large-excursions}). Then
\begin{equation}
\frac1{\#\mathcal K(X)}\sum_{\nu\in\mathcal K(X)}\log\frac{v_\nu}{u_\nu}
\ \le\ \frac{\pi}{\gamma_1}+o(1)=\log\lambda_*+o(1).
\end{equation}
\end{corollary}

\begin{proof}
Discard the initial excursion $(0,2]$, whose start is $0$ and whose
contribution to the count is $O(1)$. Fix $A>1$ before taking $X\to\infty$.
By Corollary~\ref{cor:KP-local-returns}, all but finitely many excursions
satisfy $v\le u^A$. These exceptions form a fixed global finite set once
$A$ is fixed, so the sum of their individual finite values $\log(v/u)$ is
$O_A(1)$. Proposition~\ref{prop:log-budget} gives a total of at most
$A\log X+O_A(1)$. Since
$\#\mathcal K(X)\ge(\gamma_1/\pi-o(1))\log X$, it follows for every fixed
$A>1$ that
\begin{equation}
\limsup_{X\to\infty}
\frac1{\#\mathcal K(X)}\sum_{\nu\in\mathcal K(X)}
\log\frac{v_\nu}{u_\nu}
\le \frac{A\pi}{\gamma_1}.
\end{equation}
Only after this limit is taken do we let $A\downarrow1$, obtaining
$\pi/\gamma_1$.
\end{proof}

\begin{remark}[What lowering $\lambda_*$ requires]
\label{rem:surplus}
Corollary~\ref{cor:descending}(iii) is the statement that the
excursions carrying Kaczorowski--Pintz witnesses cannot all have
$\log(v/u)>\log\lambda$ with $\lambda>\lambda_*$;
Corollary~\ref{cor:KP-mean-log-length} is its averaged form. Consequently,
any proof that such excursions have mean logarithmic length at least
$\pi/\gamma_1+\delta$ for some fixed $\delta>0$ would contradict
Corollary~\ref{cor:KP-mean-log-length} and prove RH, and a surplus
$\delta>\log(\lambda_*/\lambda)$ among the excursions whose witness lies
beyond $\lambda x_0$ would extend Corollary~\ref{cor:descending}(iii) to
that $\lambda<\lambda_*$. Kaczorowski--Pintz, as stated, gives neither: it
supplies at least one return in each band $[Y^{1-\delta},Y]$ but no upper
separation between consecutive witnesses. Nor do short-interval estimates
for $\mu$ supply it: $|M(x+h)-M(x)|=o(h)$ is compatible with $M$ staying
positive throughout $(x,x+h]$, whereas a return requires a decrement of
size at least $M(x)$ in a prescribed direction. Energy bounds do not help
either, since a flat excursion of height $1$ has $\mathfrak b_I,Q_I\asymp\ell_I$.
\end{remark}

\subsection{A deductive audit below the counting threshold}
\label{subsec:surplus-audit}

The preceding counting argument identifies the exact information that would
be sufficient to push the multiplicative horizon below $\lambda_*$.  We make
that requirement quantitative and record a matching no-go statement showing
that it does not follow from the geometric inputs alone.

Put
\begin{equation}
d_1:=\frac{\gamma_1}{\pi},\qquad
a_*:=\frac1{d_1}=\frac{\pi}{\gamma_1},\qquad
\lambda_*=e^{a_*}.
\end{equation}
Assume $\Theta>1/2$, fix $0<\eta<\Theta-\frac12$, and let
$\mathcal K(X)$ be a Kaczorowski--Pintz witness family as in
Corollary~\ref{cor:KP-mean-log-length}.  For each
$\nu\in\mathcal K(X)$ choose one witness $t_\nu$ in the corresponding
excursion $(u_\nu,v_\nu]$.

Fix $1<\lambda<\lambda_*$ and define the bad subfamily
\begin{equation}
\mathcal B_\lambda(X)
:=\{\nu\in\mathcal K(X):t_\nu>\lambda u_\nu\}.
\end{equation}

\begin{proposition}[Positional surplus criterion]
\label{prop:positional-surplus}
\tagNE\ Suppose there is a fixed $\delta>\log(\lambda_*/\lambda)$ such that
\begin{equation}
\label{eq:positional-surplus}
\frac1{\#\mathcal B_\lambda(X)}
\sum_{\nu\in\mathcal B_\lambda(X)}
\left(\log\frac{v_\nu}{u_\nu}-\log\lambda\right)
\ge \delta-o(1)
\end{equation}
whenever $\mathcal B_\lambda(X)$ is nonempty.  Then
\begin{equation}
\#\mathcal B_\lambda(X)
\le\left(\frac1{\log\lambda+\delta}+o(1)\right)\log X
<(d_1-o(1))\log X.
\end{equation}
Consequently $\gg\log X$ witnesses satisfy $t_\nu\le\lambda u_\nu$, and
the horizon argument gives the same contradiction under $\Theta>1/2$ as in
Theorem~\ref{thm:horizon}.
\end{proposition}

\begin{proof}
Fix first an exponent $A>1$.  Return localization and
Proposition~\ref{prop:log-budget} give
\begin{equation}
\sum_{\nu\in\mathcal K(X)}\log\frac{v_\nu}{u_\nu}
\le A\log X+O_A(1).
\end{equation}
After taking $X\to\infty$ we may let $A\downarrow1$, exactly as in the proof
of Corollary~\ref{cor:KP-mean-log-length}; hence the total logarithmic budget
is $(1+o(1))\log X$.  By \eqref{eq:positional-surplus}, every bad member
consumes on average at least
$\log\lambda+\delta-o(1)$ of this budget.  Therefore
\begin{equation}
(\log\lambda+\delta-o(1))\#\mathcal B_\lambda(X)
\le(1+o(1))\log X.
\end{equation}
Since
$\log\lambda+\delta>\log\lambda_*=a_*=1/d_1$, the resulting coefficient is
strictly smaller than $d_1$.  Subtracting the bad excursions from
$\#\mathcal K(X)\ge(d_1-o(1))\log X$ leaves a positive proportion, on the
logarithmic scale, of witnesses covered by the shorter horizon.
\end{proof}

The surplus has the exact decomposition
\begin{equation}
\label{eq:surplus-decomposition}
\log\frac{v_\nu}{u_\nu}-\log\lambda
=
\log\frac{t_\nu}{\lambda u_\nu}
+
\log\frac{v_\nu}{t_\nu}.
\end{equation}
Thus the existence of a few long excursions is irrelevant: the required
surplus must hold uniformly on average over the bad witness subfamily.

Unit step size alone does not provide such a fixed surplus.  If a witness
beyond the horizon has height $H\asymp t^\alpha$ with $\alpha<1$, and even if
one knows $M(\lambda u)=o(H)$ in the relevant orientation, then the step bound
gives only
\begin{equation}
t-\lambda u\ge H-o(H),\qquad v-t\ge H,
\end{equation}
and consequently
\begin{equation}
\frac vu\ge\lambda+(2+o(1))\frac Hu=\lambda+o(1).
\end{equation}
The gain is additive and sublinear, whereas
Proposition~\ref{prop:positional-surplus} requires a fixed logarithmic gain.

\begin{proposition}[Geometric inputs do not force positional surplus]
\label{prop:surplus-nogo}
\tagNE\ Disjointness of the excursions, increments in $\{-1,0,1\}$, a
localized return, a witness count $(d_1-o(1))\log X$, and a mean logarithmic
length bounded by $a_*+o(1)$ do not imply that a positive proportion of
witnesses lies before $\lambda u$ for any fixed $\lambda<\lambda_*$.  This
remains true even if the witness height is $\asymp t^\alpha$ for a fixed
$1/2<\alpha<1$ and every excursion is eventually shorter than $u^A$ for each
fixed $A>1$.
\end{proposition}

\begin{proof}
Choose
\begin{equation}
1<\lambda<c<r<\lambda_*,\qquad \frac12<\alpha<1,
\end{equation}
and, starting from a sufficiently large integer $u_1$, define recursively
\begin{equation}
u_{j+1}=\lceil r u_j\rceil,\qquad
v_j=u_{j+1},\qquad
t_j=\lfloor c u_j\rfloor.
\end{equation}
For all sufficiently large $j$, one has $u_j<t_j<v_j$.  Since both gaps
$t_j-u_j$ and $v_j-t_j$ are linear in $u_j$, whereas
$H_j:=\lfloor t_j^\alpha/2\rfloor=o(u_j)$, one can construct on
$(u_j,v_j]$ a primitive bridge with increments in $\{-1,0,1\}$ that rises to
height $H_j$ by time $t_j$, stays strictly away from zero while inserting
horizontal steps, and returns to zero at $v_j$.

Moreover,
\begin{equation}
\#\{j:u_j\le X\}=\frac{\log X}{\log r}+O(1)>d_1\log X+O(1),
\end{equation}
while
\begin{equation}
\frac1J\sum_{j\le J}\log\frac{v_j}{u_j}\longrightarrow\log r<a_*,
\qquad
t_j>\lambda u_j.
\end{equation}
Finally $v_j\asymp r u_j$, so $v_j\le u_j^A$ for every fixed $A>1$ and all
sufficiently large $j$.  This is not a construction of the Mertens function;
it is a countermodel to any deduction using only the listed geometric and
counting inputs.
\end{proof}

\subsection{Wide peaks from cancellation in all short intervals}
\label{subsec:wide-peaks}

The only arithmetic input of the following lemma is the quantitative short-interval theorem of
Matom\"aki--Ter\"av\"ainen \cite[Theorem~1.1]{MatomakiTeravainen2023}: for every fixed
$\theta>0.55$ and every fixed $\varepsilon>0$, uniformly for the relevant intervals,
\begin{equation}
 \sum_{x<n\le x+H}\mu(n)
 \ll_{\theta,\varepsilon} \frac{H}{(\log x)^{1/3-\varepsilon}}
 \qquad(H\ge x^\theta).
\end{equation}

\begin{lemma}[Mesoscopic Lipschitz bound]
\label{lem:meso-lipschitz}
\tagC\ Fix $\Lambda>1$ and $\theta_{\rm w}\in(0.55,1)$. For large $u$ put
$h_u=\lceil(\Lambda u)^{\theta_{\rm w}}\rceil$ and
\begin{equation}
\omega_u:=\max\Bigl\{h_u^{-1},\ \max_{u\le a\le\Lambda u}
\frac{|M(a+h_u)-M(a)|}{h_u}\Bigr\}.
\end{equation}
Then $\omega_u\to0$, and for all integers $u\le a<b\le\Lambda u$,
\begin{equation}
\label{eq:meso-lipschitz}
|M(b)-M(a)|\le\omega_u(b-a)+h_u .
\end{equation}
\end{lemma}

\begin{proof}
Fix $\theta'\in(0.55,\theta_{\rm w})$. For all sufficiently large $u$,
$a+h_u\le2\Lambda u$ uniformly for $u\le a\le\Lambda u$. By
\cite[Theorem~1.1]{MatomakiTeravainen2023}, for every fixed
$\varepsilon>0$ there is a constant $C$ such that, uniformly for
$u\le b\le2\Lambda u$,
\begin{equation}
\left|M\!\left(b+\lceil b^{\theta'}\rceil\right)-M(b)\right|
\le C\frac{b^{\theta'}}{(\log u)^{1/3-\varepsilon}}
\end{equation}
for all sufficiently large $u$. Theorem~1.1 of
\cite{MatomakiTeravainen2023} is uniform in the starting point $x$ for
intervals of length $x^{\theta'}$; here $b\asymp_\Lambda u$, so
$\log b=\log u+O_\Lambda(1)$, and replacing $x^{\theta'}$ by
$\lceil b^{\theta'}\rceil$ changes the sum by $O(1)$, absorbed in $C$.

For a given $a$, define recursively
\begin{equation}
 b_0=a,\qquad b_{j+1}=b_j+\lceil b_j^{\theta'}\rceil,
\end{equation}
and stop at the first $J$ for which $b_{J+1}>a+h_u$. Since
$b_j\asymp u$ throughout the construction, each full block has length
$\asymp u^{\theta'}$, whence
\begin{equation}
 J\ll \frac{h_u}{u^{\theta'}}+1,
 \qquad
 \sum_{j<J}b_j^{\theta'}\ll h_u.
\end{equation}
All starting points $b_j$ lie in $[u,2\Lambda u]$. Summing the full-block
estimates and bounding the final incomplete piece trivially gives,
uniformly for $u\le a\le\Lambda u$,
\begin{equation}
\frac{|M(a+h_u)-M(a)|}{h_u}
\ll_{\varepsilon,\theta',\theta_{\rm w},\Lambda}
(\log u)^{-1/3+\varepsilon}+u^{\theta'-\theta_{\rm w}}.
\end{equation}
Consequently
\begin{equation}
\label{eq:eta-uniform}
\omega_u\ll_{\varepsilon,\theta',\theta_{\rm w},\Lambda}
(\log u)^{-1/3+\varepsilon}+u^{\theta'-\theta_{\rm w}}+h_u^{-1}=o(1),
\end{equation}
and the same estimate is uniform for $u\in[U,2U]$ as $U\to\infty$.

For \eqref{eq:meso-lipschitz}, write $b-a=qh_u+r$ with $0\le r<h_u$.
Every complete $h_u$-block starts at a point in $[u,\Lambda u]$ and
contributes at most $\omega_u h_u$, while the terminal remainder contributes
at most $r<h_u$. Thus
$|M(b)-M(a)|\le\omega_u(b-a)+h_u$.
\end{proof}

\begin{proposition}[Wide peaks]
\label{prop:wide-peaks}
\tagN\ With $\Lambda,\theta_{\rm w},h_u,\omega_u$ as in
Lemma~\ref{lem:meso-lipschitz}, let $I_\nu=(u,v]$ be an excursion of
orientation $\sigma_\nu$ with $v\le\Lambda u$ and $u$ large, let $t$ be an
integer with $u<t<v$ and $|M(t)|=H$, and put $G=(H-h_u)_+$ and
$m=\lfloor G/\omega_u\rfloor$. Then
\begin{equation}
t-u\ge\frac G{\omega_u},\qquad v-t\ge\frac G{\omega_u},
\end{equation}
\begin{equation}
\mathfrak b\ge H+2\sum_{j=1}^m(G-\omega_uj),\qquad
Q\ge H^2+2\sum_{j=1}^m(G-\omega_uj)^2 .
\end{equation}
Consequently, uniformly for such excursions with $H/h_u\to\infty$,
\begin{equation}
v-u\ge(2+o(1))\frac H{\omega_u},\qquad
\mathfrak b\ge(1+o(1))\frac{H^2}{\omega_u},\qquad
Q\ge\Bigl(\frac23+o(1)\Bigr)\frac{H^3}{\omega_u},
\end{equation}
so that $(v-u)/H$, $\mathfrak b/H^2$ and $Q/H^3$ tend to infinity, whereas the
unit-step bounds of Lemma~\ref{lem:height_steps} and
Proposition~\ref{prop:sharp-LH} give only $2$, $1$ and $\frac23$.
\end{proposition}

\begin{proof}
By \eqref{eq:meso-lipschitz} applied to $(u,t]$ and $(t,v]$,
$H\le\omega_u(t-u)+h_u$ and $H\le\omega_u(v-t)+h_u$.  Thus, since the two
distances are integers,
\begin{equation}
 t-u,\ v-t\ge\left\lceil\frac{G}{\omega_u}\right\rceil,
\end{equation}
and therefore $m=\lfloor G/\omega_u\rfloor\le t-u,v-t$.  For
$1\le j\le m$ the points $t\pm j$ lie in $[u,v]$, and \eqref{eq:meso-lipschitz} gives
$\sigma_\nu M(t\pm j)\ge G-\omega_uj\ge0$. Summing $t$ and the pairs $t\pm j$
gives the bounds for $\mathfrak b$ and $Q$. Finally
$2\sum_{j\le m}(G-\omega_uj)=G^2/\omega_u+O(G)$ and
$2\sum_{j\le m}(G-\omega_uj)^2=\frac23G^3/\omega_u+O(G^2)$. Under
$H/h_u\to\infty$ we have $G=H-h_u=(1+o(1))H$ and $G\to\infty$.
Equation~\eqref{eq:eta-uniform} gives $\omega_u=o(1)$, hence eventually
$\omega_u\le1$ and the relative errors satisfy $\omega_u/G\le1/G=o(1)$.
This proves the stated uniform asymptotics.
\end{proof}

\begin{corollary}[Packing of high excursions]
\label{cor:packing}
\tagN\ Let $Y=Y(U)$ with $Y/U^{\theta_{\rm w}}\to\infty$. The number of excursions with
$U\le u\le2U$, $v\le\Lambda u$ and height at least $Y$ is
$o(U/Y)$, whereas the unit-step bound gives $O(U/Y)$.
\end{corollary}

\begin{proof}
They are disjoint and contained in $[U,2\Lambda U]$. By
Proposition~\ref{prop:wide-peaks}, each has length at least
\begin{equation}
 (2+o(1))\frac{Y}{\max_{U\le u\le2U}\omega_u}.
\end{equation}
If $N(U,Y)$ denotes their number, disjointness therefore gives
\begin{equation}
 N(U,Y)(2+o(1))\frac{Y}{\max_{U\le u\le2U}\omega_u}
 \le (2\Lambda-1)U.
\end{equation}
By the uniform estimate \eqref{eq:eta-uniform},
\begin{equation}
\omega_U^*:=\max_{U\le u\le2U}\omega_u=o(1).
\end{equation}
Therefore
\begin{equation}
N(U,Y)\ll_\Lambda \omega_U^*\frac{U}{Y}=o(U/Y),
\end{equation}
which completes the packing argument.
\end{proof}

\begin{remark}[Scope]
\label{rem:wide-peaks-scope}
Proposition~\ref{prop:wide-peaks} is informative only for heights above
$u^{0.55+o(1)}$. If $\Theta>0.55$ and $0.55<\theta_{\rm w}<\Theta-\eta$, it
applies to the Kaczorowski--Pintz witness excursions of
Corollary~\ref{cor:two-sided} with $v\le\lambda u$: their peaks are wide in
the sense above. It does not lower $\lambda_*$: the resulting lower bound
for $\log(v/u)$ is of order $H/(\omega_uu)$, which may still tend to $0$,
whereas Remark~\ref{rem:surplus} requires a fixed surplus. The published
quantitative estimate of Matom\"aki--Ter\"av\"ainen
\cite[Theorem~1.1]{MatomakiTeravainen2023} makes the gain explicit. If
$0.55<\theta'<\theta_{\rm w}$, then the blocking argument in
Lemma~\ref{lem:meso-lipschitz} gives, for every fixed $\varepsilon>0$,
\begin{equation}
\omega_u\ll_{\varepsilon,\theta',\theta_{\rm w},\Lambda}
(\log u)^{-1/3+\varepsilon}+u^{\theta'-\theta_{\rm w}}.
\end{equation}
Consequently, in the regime $H/h_u\to\infty$ of
Proposition~\ref{prop:wide-peaks},
\begin{equation}
Q\gg_{\varepsilon}H^3(\log u)^{1/3-\varepsilon}.
\end{equation}
\end{remark}




\section{Arithmetic selection on the Catalan reflection fibre}
\label{sec:arithmetic-selection-fibre}

The fixed-mask Catalan description separates two finite degrees of freedom of a completed Mertens excursion: its positive support-time shape and its orientation.  The Catalan estimates developed above act on the first coordinate.  The present section isolates the second coordinate and shows that it is strongly constrained by divisor inversion.  The resulting arithmetic selector is nevertheless insufficient, by itself, to observe a fixed-shift Chowla form on the ambient squarefree space.  This produces a sharp separation between orientation selection and two-point cancellation.

Throughout this section, let $N$ be an integer with $M(N)=0$, so that the portion of the path up to $N$ consists only of completed excursions and zero runs.  Write
\begin{equation}
I_\nu=(x_{0,\nu},x_{1,\nu}],\qquad 1\le \nu\le J_N,
\end{equation}
for the completed excursions up to $N$, and let
\begin{equation}
E_\nu:=\{d\in I_\nu:\mu(d)\ne0\}
\end{equation}
be the active support of the $\nu$th excursion.  For a word $w\in\{-1,0,1\}^N$ define
\begin{equation}
M_w(n):=\sum_{d\le n}w(d),\qquad
\Pi_N(w):=\bigl(M_w(n)^2\bigr)_{1\le n\le N}.
\end{equation}
We restrict to words having the same zero mask as $\mu$.

\begin{theorem}[Exact reflection fibre]
\label{thm:exact-reflection-fibre}
\tagNE\ Let $w\in\{-1,0,1\}^N$ have the same support as $\mu$.  Then
\begin{equation}
\Pi_N(w)=\Pi_N(\mu)
\end{equation}
if and only if, on each completed excursion $I_\nu$, the word $w$ agrees either with $\mu$ or with $-\mu$.  Consequently the fibre of the complete quadratic profile has exactly
\begin{equation}
2^{J_N}
\end{equation}
elements.
\end{theorem}

\begin{proof}
Both partial-sum paths start from the same value, $M_w(0)=M(0)=0$, and $w,\mu\in\{-1,0,1\}^N$ have the same zero mask.  Independent reflection of a completed excursion multiplies $M$ by $-1$ on the interior of that excursion and leaves it zero at both endpoints; hence it preserves $M(n)^2$ at every integer.  Conversely, equality of the quadratic profiles gives the same zero set and the same absolute path.  On the interior of each completed excursion one has
\begin{equation}
|M_w(n)|=|M(n)|>0,
\end{equation}
so $M_w$ has no interior zero and therefore has constant sign there.  Since $M$ also has constant sign on that excursion, the ratio $M_w/M$ is a constant in $\{\pm1\}$.  Taking first differences successively across that component shows that the same constant relates $w$ and $\mu$ on the active support of the excursion.  The choices are independent because every completed excursion begins and ends at zero.
\end{proof}

Theorem~\ref{thm:exact-reflection-fibre} identifies the precise information lost by every statistic determined by the full quadratic profile: after the absolute excursion shapes are fixed, the remaining ambiguity is a Boolean orientation vector.

For $\mathbf z=(\mathbf z_1,\dots,\mathbf z_{J_N})\in\{0,1\}^{J_N}$, let $w_\mathbf z$ be obtained from $\mu$ by reflecting precisely those excursions for which $\mathbf z_\nu=1$.  Fix arbitrary weights $\varrho_n>0$ and define the divisor-defect functional
\begin{equation}
\mathfrak D_N(w)
:=
\varrho_1|w(1)-1|^2
+
\sum_{n=2}^N\varrho_n
\left|\sum_{d\mid n}w(d)\right|^2.
\label{eq:divisor-defect-functional}
\end{equation}
For every excursion put
\begin{equation}
g_\nu(n):=
\sum_{\substack{d\mid n\\ d\in E_\nu}}\mu(d),
\qquad 1\le n\le N,
\label{eq:excursion-divisor-signal-new}
\end{equation}
and define the $J_N\times J_N$ matrix
\begin{equation}
(G_N)_{\nu\kappa}
:=
\sum_{n=1}^N\varrho_n g_\nu(n)g_\kappa(n).
\label{eq:orientation-gram-new}
\end{equation}

\begin{theorem}[Positive Gram selector for excursion orientations]
\label{thm:positive-gram-selector}
\tagNE\ The matrix $G_N$ is positive semidefinite and
\begin{equation}
\mathfrak D_N(w_\mathbf z)
=
4\,\mathbf z^{\mathsf T}G_N\mathbf z
\qquad
(\mathbf z\in\{0,1\}^{J_N}).
\label{eq:gram-selector-identity}
\end{equation}
Moreover,
\begin{equation}
\mathbf z^{\mathsf T}G_N\mathbf z>0
\qquad\text{for every }\mathbf z\ne0.
\label{eq:boolean-strict-positivity}
\end{equation}
Thus the divisor equations select the arithmetic orientation uniquely inside the complete quadratic-profile fibre.
\end{theorem}

\begin{proof}
For $n\ge2$, M\"obius inversion gives
\begin{equation}
\sum_{d\mid n}\mu(d)=0.
\end{equation}
Reflecting the excursions indexed by $\mathbf z$ therefore gives
\begin{equation}
\sum_{d\mid n}w_\mathbf z(d)
=-2\sum_{\nu=1}^{J_N}\mathbf z_\nu g_\nu(n).
\end{equation}
At $n=1$ one has explicitly
\begin{equation}
w_\mathbf z(1)-1
=-2\sum_{\nu=1}^{J_N}\mathbf z_\nu g_\nu(1),
\end{equation}
with only the excursion containing $1$ able to contribute.  Substitution into \eqref{eq:divisor-defect-functional} yields \eqref{eq:gram-selector-identity}, and \eqref{eq:orientation-gram-new} is a Gram representation, so $G_N\succeq0$.

If a nonzero Boolean vector $\mathbf z$ satisfied $\mathbf z^{\mathsf T}G_N\mathbf z=0$, then $\mathfrak D_N(w_\mathbf z)=0$.  Hence $w_\mathbf z(1)=1$ and
\begin{equation}
\sum_{d\mid n}w_\mathbf z(d)=0\qquad(2\le n\le N).
\end{equation}
Finite M\"obius inversion determines the coefficients recursively and forces $w_\mathbf z(d)=\mu(d)$ for every $d\le N$, contradicting the reflection of a nonempty set of excursions.
\end{proof}

The Gram form is exact but uses all divisor equations.  A much smaller family of specially placed divisor equations already compresses the orientation fibre from exponential to polynomial size.

For the $\nu$th excursion list its active support as
\begin{equation}
a_{\nu,1}<\cdots<a_{\nu,L_\nu^\#},
\qquad L_\nu^\#=2r_\nu,
\end{equation}
and define its support centre
\begin{equation}
m_\nu:=a_{\nu,r_\nu}.
\end{equation}
Set
\begin{equation}
\phi_\nu(d):=\sigma_\nu\mu(d)
\qquad(d\in E_\nu),
\end{equation}
so that $\phi_\nu$ is the oriented increment sequence of the positive support-time excursion.  Define the central divisor matrix
\begin{equation}
\mathscr C^{\mathrm{div}}_{\nu\kappa}
:=
\sum_{\substack{d\mid m_\nu\\ d\in E_\kappa}}
\phi_\kappa(d).
\label{eq:central-divisor-matrix-new}
\end{equation}

\begin{theorem}[Polynomial compression of the orientation fibre]
\label{thm:polynomial-orientation-compression}
\tagN\ The matrix $\mathscr C^{\mathrm{div}}$ is lower triangular.  Its
diagonal entry in row $\nu$ is nonzero whenever
\begin{equation}
x_{0,\nu}\ge \frac{m_\nu}{2}.
\label{eq:central-diagonal-condition}
\end{equation}
The condition \eqref{eq:central-diagonal-condition} can fail for at most
$\lfloor\log_2N\rfloor$ completed excursions.

After the affine normalization $w(1)=1$ is imposed, the solution set of the
central divisor equations is an affine subset of the orientation-coordinate
space whose linear direction has dimension at most
$\lfloor\log_2N\rfloor+1$.  More concretely, once the exceptional orientation
coordinates have been specified, every nonexceptional coordinate not already
fixed by the affine normalization is determined successively by a central
equation with nonzero diagonal coefficient.  Consequently, among the
$2^{J_N}$ Boolean orientations in the quadratic-profile fibre, at most
\begin{equation}
2^{\lfloor\log_2N\rfloor+1}\le2N
\label{eq:polynomial-survivors}
\end{equation}
satisfy the central divisor equations.  Furthermore, at most
$\lfloor\log_2N\rfloor+1$ additional divisor rows suffice to select the
M\"obius orientation uniquely.
\end{theorem}

\begin{proof}
If $\kappa>\nu$, every $d\in E_\kappa$ exceeds $m_\nu$, so no such $d$ can
divide $m_\nu$; hence $\mathscr C^{\mathrm{div}}_{\nu\kappa}=0$.  If
\eqref{eq:central-diagonal-condition} holds, every proper divisor of
$m_\nu$ is at most $m_\nu/2\le x_{0,\nu}$ and therefore lies outside
$E_\nu$.  Thus only $d=m_\nu$ contributes to the diagonal, giving
\begin{equation}
\mathscr C^{\mathrm{div}}_{\nu\nu}=\phi_\nu(m_\nu)\in\{\pm1\}.
\end{equation}

Since $m_\nu\le x_{1,\nu}$, failure of
\eqref{eq:central-diagonal-condition} implies
$x_{1,\nu}>2x_{0,\nu}$.  Count the initial excursion separately.  If
$2\le a_1<\cdots<a_r$ are the positive starts of the remaining failing
completed excursions, then disjointness gives
\begin{equation}
N\ge x_{1,r}>2a_r>2^2a_{r-1}>\cdots>2^r a_1\ge2^{r+1}.
\end{equation}
Thus the total number of failures, including the possible initial one, is
strictly smaller than $\log_2N$ when $N$ is not a power of two and at most
$\log_2N-1$ when it is; in either case it is at most
$\lfloor\log_2N\rfloor$.

Let $\boldsymbol\sigma=(\sigma_1,\dots,\sigma_{J_N})^{\mathsf T}$ be the
arithmetic orientation vector.  For every central row with $m_\nu>1$,
\begin{equation}
(\mathscr C^{\mathrm{div}}\boldsymbol\sigma)_\nu
=
\sum_{d\mid m_\nu}\mu(d)=0.
\end{equation}
The row $n=1$ is imposed separately as the affine normalization $w(1)=1$.
Let $\mathcal E_N$ be the set of indices for which
\eqref{eq:central-diagonal-condition} fails.  Among the central rows with
$m_\nu>1$, every nonexceptional row has a nonzero new diagonal coefficient.
Thus their coefficient matrix has rank at least
\begin{equation}
J_N-|\mathcal E_N|-1,
\end{equation}
the possible loss of one row accounting only for the distinguished index
$m_\nu=1$.  Hence the associated homogeneous linear system has nullity at
most $|\mathcal E_N|+1$.  Intersecting this linear solution space with the
affine normalization $w(1)=1$ gives an affine solution set whose linear
direction has dimension at most $|\mathcal E_N|+1\le
\lfloor\log_2N\rfloor+1$.

For the Boolean count we use the triangular structure directly rather than a
dimension argument.  Choose first the coordinates indexed by $\mathcal E_N$
and, if it is not already among them, the single coordinate attached to the
row $m_\nu=1$.  Then process every remaining nonexceptional row in increasing
order.  All earlier coordinates are already known, while the coefficient of
the new coordinate is $\mathscr C^{\mathrm{div}}_{\nu\nu}\ne0$, so that
coordinate is uniquely determined.  The affine equation $w(1)=1$ can only
remove choices, not create them.  Therefore there are at most
$2^{|\mathcal E_N|+1}\le2^{\lfloor\log_2N\rfloor+1}\le2N$ Boolean
solutions.  In particular the Boolean estimate is not inferred from real
dimension alone.

For the final assertion, for a generic real orientation vector
$\boldsymbol\sigma'=(\sigma'_1,\dots,\sigma'_{J_N})$ define
$w_{\boldsymbol\sigma'}(d)=\sigma'_\nu\phi_\nu(d)$ for $d\in E_\nu$ and
$w_{\boldsymbol\sigma'}(d)=0$ off the active support.  Factor the complete
divisor readout as
\begin{equation}
\boldsymbol\sigma'
\longmapsto
w_{\boldsymbol\sigma'}
\longmapsto
\left(\sum_{d\mid n}w_{\boldsymbol\sigma'}(d)\right)_{1\le n\le N}.
\end{equation}
The first map is injective: the supports $E_\nu$ are disjoint and every
completed excursion contains an active increment, so two distinct
orientation vectors give distinct words.  The second map is injective on
$\mathbb R^N$ because its matrix is
$(\mathbf 1_{d\mid n})_{1\le n,d\le N}$, which is lower unitriangular.
Hence the composition has full
column rank $J_N$ on the orientation coordinates.  The central rows with
$m_\nu>1$ have linear rank at least $J_N-|\mathcal E_N|-1$, independently of
the separate affine normalization $w(1)=1$.  Extending a basis of their row
space to a basis of the full readout therefore needs at most
$|\mathcal E_N|+1\le\lfloor\log_2N\rfloor+1$ additional divisor rows.
The full-rank system has at most one orientation solution, and the arithmetic
orientation $\boldsymbol\sigma$ is such a solution; it is therefore unique.
\end{proof}

The preceding theorem concerns selection inside the highly structured reflection fibre.  It does not imply that a fixed-shift correlation can be recovered from a comparably small family of divisor measurements.  In fact a linear obstruction holds on the ambient squarefree space.

Let $\mathbf D_N$ be the divisor-incidence matrix
\begin{equation}
(\mathbf D_N)_{r,d}:=\mathbf 1_{d\mid r},
\qquad 1\le r,d\le N,
\end{equation}
and, for $\mathcal T\subseteq\{1,\dots,N\}$, let $\mathbf D_\mathcal T$ denote the restriction to rows in $\mathcal T$.  Let $\mathsf A_h$ be the symmetric fixed-shift matrix defined by
\begin{equation}
x^{\mathsf T}\mathsf A_hx
=
\sum_{n\le N-h}x_nx_{n+h}.
\end{equation}
Write
\begin{equation}
\mathfrak S_h
:=
\prod_{p^2\mid h}\left(1-\frac1{p^2}\right)
\prod_{p^2\nmid h}\left(1-\frac2{p^2}\right)>0,
\label{eq:mirsky-density-new}
\end{equation}
so that, by the classical squarefree-pair theorem of Mirsky \cite{Mirsky1949},
\begin{equation}
\#\{n\le X:n,n+h\text{ squarefree}\}
=
\mathfrak S_hX+o_h(X)
\end{equation}
for each fixed nonzero $h$.

\begin{theorem}[Linear obstruction to divisor-row observability]
\label{thm:linear-observability-obstruction}
\tagNE\ Fix $h\ge1$.  Suppose that for some $c_{\rm obs}\in[0,\infty)$, some $\varepsilon\in[0,1/2)$, and some row set $\mathcal T\subseteq\{1,\dots,N\}$ one has
\begin{equation}
|x^{\mathsf T}\mathsf A_hx|
\le
c_{\rm obs}\|\mathbf D_\mathcal T x\|_2^2
+
\varepsilon\|x\|_2^2
\label{eq:uniform-observability-new}
\end{equation}
for every real vector $x$ supported on the squarefree integers in $[1,N]$.  Then
\begin{equation}
|\mathcal T|
\ge
\frac{\mathfrak S_h}{4}N+o_h(N).
\label{eq:linear-row-lower-bound}
\end{equation}
In particular, no $o(N)$ family of divisor rows can uniformly observe the fixed-shift quadratic form with error below $1/2$.
\end{theorem}

\begin{proof}
Let $N/2<b\le N-h$ be such that both $b$ and $b+h$ are squarefree.  Since $b>N/2$, column $b$ of $\mathbf D_N$ is nonzero only in row $b$; the same holds for $b+h$.  If neither row belongs to $\mathcal T$, then for
\begin{equation}
x=e_b\pm e_{b+h}
\end{equation}
one has $\mathbf D_\mathcal T x=0$, while
\begin{equation}
|x^{\mathsf T}\mathsf A_hx|=1,
\qquad
\|x\|_2^2=2.
\end{equation}
Thus \eqref{eq:uniform-observability-new} would force $1\le2\varepsilon$, contrary to $\varepsilon<1/2$.  Hence $\mathcal T$ must meet every edge $\{b,b+h\}$ of the graph of squarefree shifted pairs in $(N/2,N]$.

By \eqref{eq:mirsky-density-new}, this graph has
\begin{equation}
\frac{\mathfrak S_h}{2}N+o_h(N)
\end{equation}
edges.  Its maximum degree is at most $2$, so every vertex cover has cardinality at least half the number of edges.  This gives \eqref{eq:linear-row-lower-bound}.
\end{proof}

\begin{corollary}[Selection--observability separation]
\label{cor:selection-observability-separation}
\tagNE\ For cutoffs $N$ with $M(N)=0$, the complete quadratic profile has an orientation fibre of size $2^{J_N}$.  The central divisor equations reduce that fibre to at most $2N$ candidates, and $O(\log N)$ further divisor equations select the arithmetic orientation uniquely.  In contrast, for every fixed $h\ge1$, uniform observation of the ambient fixed-shift form $x^{\mathsf T}\mathsf A_hx$ with error strictly below $1/2$ requires $\Omega_h(N)$ divisor rows.

Consequently, the unresolved difficulty in fixed-shift M\"obius cancellation is not the finite identification of the orientation bits once the absolute excursion geometry is known.  It is the collective cancellation of the shift interaction across the arithmetic path.  The Catalan selection problem and the Chowla cancellation problem therefore occupy genuinely different complexity scales.
\end{corollary}


\section{Prefix structure of the localized window}
\label{sec:prefix}

\subsection{The linear window and the one-step recursion}
\label{subsec:linear-window}

Once the horizon is below $2x_0$, the square-range bilinear structure of
Proposition~\ref{prop:prefix-bilinear} is no longer needed: on $(u,2u]$ the
M\"obius function is a \emph{linear} expression in the prefix
$\mu\mathbf 1_{[1,u]}$.

\begin{proposition}[Linear window]
\label{prop:linear-window}
\tagL\ Let $u\ge1$. For every integer $u<n\le2u$,
\begin{equation}
\label{eq:linear-window-mu}
\mu(n)=-\sum_{\substack{d\mid n\\ d\le u}}\mu(d).
\end{equation}
Consequently, if $M(u)=0$ and $u<t\le2u$,
\begin{equation}
\label{eq:linear-window-M}
M(t)=-\sum_{d\le u}\mu(d)\left(\left\lfloor\frac td\right\rfloor
-\left\lfloor\frac ud\right\rfloor\right).
\end{equation}
\end{proposition}

\begin{proof}
Since $n\ge2$, $\sum_{d\mid n}\mu(d)=0$. Every proper divisor of $n$ is at
most $n/2\le u$, and every divisor $d\le u$ is proper because $n>u$. Summing
\eqref{eq:linear-window-mu} over $u<n\le t$, the number of multiples of $d$
in $(u,t]$ is $\lfloor t/d\rfloor-\lfloor u/d\rfloor$.
\end{proof}

\begin{corollary}[Localized prefix formula]
\label{cor:localized-prefix}
\tagNE\ Let $M(u)=0$, let $1<\lambda<2$, and let $u<t\le\lambda u$.  Then
\begin{equation}
\label{eq:localized-prefix}
M(t)=-\sum_{d\le\lambda u/2}\mu(d)
\left(\left\lfloor\frac td\right\rfloor-
\left\lfloor\frac ud\right\rfloor\right).
\end{equation}
Indeed, if $\lambda u/2<d\le u$, then $d>u/2$ and the first multiple of
$d$ strictly larger than $u$ is $2d>\lambda u\ge t$.
\end{corollary}

\begin{corollary}[Prefix-dependency cutoff]
\label{cor:prefix-cutoff}
\tagN\ Put
\begin{equation}
\chi_*:=\frac12e^{\pi/\gamma_1}=0.62445\ldots .
\end{equation}
For every fixed $\chi$ with $\chi_*<\chi<1$, RH is equivalent to
Littlewood's bound on every excursion starting at $u$ throughout
\begin{equation}
u<t\le\min\{x_1,2\chi u\},
\end{equation}
where, throughout that window, $M(t)$ is given by the sum
\eqref{eq:localized-prefix} using only the prefix
$\mu(d)$ with $d\le\chi u$.  In particular, the choice
$\chi=5/8$ corresponds to $\lambda=5/4$ and gives the RH-equivalent
window
\begin{equation}
1\le r=t-u\le\min\{\ell-1,u/4\}
\end{equation}
using only $d\le5u/8$. The return point $r=\ell$, when it belongs to the
window, is omitted because $M(u+\ell)=0$.

The number $\chi_*$ is an infimum, not an attained endpoint, because the
horizon theorem requires $\lambda>\lambda_*$.
\end{corollary}

\begin{proof}
Take $\lambda=2\chi>\lambda_*$ in
Corollary~\ref{cor:descending}(iii), then use
Corollary~\ref{cor:localized-prefix}.
\end{proof}


\begin{corollary}[Algebraic dependency gap]
\label{cor:dependency-gap}
\tagNE\ Let $\chi_*<\chi<1$.  If $u$ is a zero of $M$, then every increment
$\mu(n)$ with
\begin{equation}
u<n\le2\chi u
\end{equation}
is determined by the prefix
$\mu(1),\ldots,\mu(\lfloor\chi u\rfloor)$.  Consequently the intermediate
block
\begin{equation}
(\chi u,u]
\end{equation}
does not occur as a direct input in the divisor recursion determining the
localized future window.  As $\chi\downarrow\chi_*$, the omitted block
has relative length tending to
$1-\chi_*=0.3755\ldots$.
\end{corollary}

\begin{proof}
For $u<n\le2\chi u$, every proper divisor $d\mid n$ satisfies
$d\le n/2\le\chi u$.  Hence
\begin{equation}
\mu(n)=-\sum_{\substack{d\mid n\\ d<n}}\mu(d)
\end{equation}
uses only $d\le\chi u$.
\end{proof}

In convolution form, \eqref{eq:linear-window-mu} says that
$(\mu-\mu_u)*1=\mu\mathbf 1_{(u,\infty)}$ on $n\le2u$, the order-one analogue
of \eqref{eq:bilinear-general}, which for $u<n\le u^2$ reads
$\mu(n)=-(\mu_u*\mu_u*1)(n)$; on $(u,2u]$ both give $\mu(n)$, the
bilinear form being the quadratic continuation of the linear one. Equivalently, \eqref{eq:linear-window-M}
follows from $\sum_{d\le t}\mu(d)\lfloor t/d\rfloor=1$, since
$\lfloor t/d\rfloor=1$ for $u<d\le t\le2u$.

\begin{proposition}[One-step recursion]
\label{prop:one-step}
\tagL\ For real $x\ge1$,
\begin{equation}
\label{eq:hyperbola-M}
\sum_{k\le x}M\!\left(\frac xk\right)=1 .
\end{equation}
Hence, if $M(u)=0$ and $t>u$,
\begin{equation}
\label{eq:descent-identity}
M(t)=-\sum_{2\le k\le t}\left[M\!\left(\frac tk\right)-M\!\left(\frac uk\right)\right],
\end{equation}
with $M(y)=0$ for $y<1$. Every argument on the right is at most $t/2$; hence on a localized
window $u<t\le\lambda u$ all arguments are at most $\lambda u/2$.  For
$\lambda=5/4$ this is $5u/8$.
\end{proposition}

\begin{proof}
$\sum_{k\le x}M(x/k)=\sum_{k\le x}\sum_{d\le x/k}\mu(d)
=\sum_{m\le x}\sum_{d\mid m}\mu(d)=1$. Here and below
$M(x/k)=M(\lfloor x/k\rfloor)$ by the global real-argument convention.
Subtract the identity at $u$, where $M(u)=0$, and extend the second sum to
$k\le t$ using $M(u/k)=0$ for $k>u$.
\end{proof}

\begin{proposition}[Scale cascade in the localized excursion window]
\label{prop:scale-cascade}
\tagNE\ Let $(u,v]$ be an excursion, let $\lambda_*<\lambda<2$, and let
$u<t\le\min\{v,\lambda u\}$.  For $k\ge2$ put
\begin{equation}
\Delta_k(u,t):=M(t/k)-M(u/k),
\end{equation}
and for an integer $K\ge2$ put
\begin{equation}
\mathcal R_K(u,t):=-\sum_{K<k\le t}\Delta_k(u,t).
\end{equation}
Then, exactly,
\begin{equation}
\label{eq:scale-cascade}
M(t)=-\sum_{k=2}^{K}\Delta_k(u,t)+\mathcal R_K(u,t),
\end{equation}
and $\mathcal R_K(u,t)$ uses only values $M(y)$ with
\begin{equation}
y\le\frac{\lambda u}{K+1}.
\end{equation}
Equivalently, all dependence on values above $\lambda u/(K+1)$ is
concentrated in the $K-1$ increments
$\Delta_2,\ldots,\Delta_K$.  Moreover
\begin{equation}
\Delta_k(u,t)=\sum_{u/k<n\le t/k}\mu(n),
\end{equation}
so the supporting real interval has length $(t-u)/k$ and contains $\lfloor t/k\rfloor-\lfloor u/k\rfloor$ integers.
\end{proposition}

\begin{proof}
Equation~\eqref{eq:scale-cascade} is Proposition~\ref{prop:one-step} with
the terms $2\le k\le K$ separated.  If $k>K$, then
$t/k\le\lambda u/(K+1)$ and $u/k\le\lambda u/(K+1)$.  The last identity is
the definition of the summatory function on the interval $(u/k,t/k]$.
\end{proof}

\begin{corollary}[First channels]
\label{cor:first-channels}
\tagNE\ For $K=2$,
\begin{equation}
M(t)=-\bigl(M(t/2)-M(u/2)\bigr)+\mathcal R_2(u,t),
\end{equation}
where $\mathcal R_2$ uses only $y\le\lambda u/3$.  With the convenient localized
choice $\lambda=5/4$, the only channel above $5u/12$ is
\begin{equation}
\Delta_2(u,u+r)=\sum_{u/2<n\le u/2+r/2}\mu(n),
\qquad 0<r\le u/4,
\end{equation}
whose supporting interval has length at most $u/8$.  In the limiting
sense $\lambda\downarrow\lambda_*$, the successive remainder cutoffs are
\begin{equation}
\frac{\lambda_*}{2}u=0.62445\ldots u,\quad
\frac{\lambda_*}{3}u=0.4162\ldots u,\quad
\frac{\lambda_*}{4}u=0.3122\ldots u,\quad
\frac{\lambda_*}{5}u=0.2497\ldots u.
\end{equation}
\end{corollary}

\begin{lemma}[Channel persistence extracted from a high excursion]
\label{lem:channel-persistence}
\tagNE\ Let $S_r=\sigma_\nu M(u+r)>0$ on $1\le r\le h$ with $S_0=0$ and
$|S_r-S_{r-1}|\le1$.  Fix $K\ge2$ and write, using
Proposition~\ref{prop:scale-cascade},
\begin{equation}
S_r=\sum_{k=2}^{K}X_k(r)+Y_K(r),
\quad
X_k(r):=-\sigma_\nu\Delta_k(u,u+r),
\quad
Y_K(r):=\sigma_\nu \mathcal R_K(u,u+r).
\end{equation}
If $H=\max_{r\le h}S_r$ and $\tau$ is a point with $S_\tau=H$, then among
the last $\lfloor H/2\rfloor+1$ indices
\begin{equation}
\tau-\lfloor H/2\rfloor,\ldots,\tau
\end{equation}
there is one fixed component among
$X_2,\ldots,X_K,Y_K$ which is at least $H/(2K)$ at no fewer than
$\lfloor H/(2K)\rfloor$ indices.
\end{lemma}

\begin{proof}
Since $S_0=0$ and the path has unit steps, $\tau\ge H$ and
$S_{\tau-j}\ge H-j\ge H/2$ for
$0\le j\le\lfloor H/2\rfloor$.  At every one of these indices the sum of
the $K$ displayed components is at least $H/2$, hence some component is at
least $H/(2K)$.  Pigeonholing the indices among the $K$ components gives the
claim.
\end{proof}

\begin{remark}[What the persistence lemma adds]
\label{rem:persistence-content}
Lemma~\ref{lem:channel-persistence} converts a
large value attained inside one excursion into a quantified persistence
requirement: either one of the finitely many scaled interval sums
$\Delta_2,\ldots,\Delta_K$ repeatedly carries a positive fraction of the
height, or the responsibility lies repeatedly in the deeper remainder,
which is supported entirely below $\lambda u/(K+1)$.  This avoids the
iteration defect of Remark~\ref{rem:iteration-gap}: no new point is claimed
to be the start-window point of another excursion.
\end{remark}
\begin{corollary}[Channel descent with explicit loss]
\label{cor:channel-descent}
\tagNE\ Let $u$ be a zero of $M$, $t>u$, $2\le K\le u$, and $0<\vartheta<1$. Then either
\begin{enumerate}[label=\textup{(\roman*)}]
\item there are $2\le k\le K$ and $y\in\{\lfloor u/k\rfloor,\lfloor t/k\rfloor\}$
with $|M(y)|\ge|M(t)|/(2K)$, hence
\begin{equation}
\frac{|M(y)|}{y^\vartheta}\ge\frac{k^\vartheta}{2K}\cdot\frac{|M(t)|}{t^\vartheta},
\qquad y\le\left\lfloor\frac t2\right\rfloor;
\end{equation}
or
\item $|\mathcal R_K(u,t)|\ge|M(t)|/K$, where $\mathcal R_K$ depends only on
$M(y)$ with $y\le\lfloor t/(K+1)\rfloor$.
\end{enumerate}
The statement holds at every $t>u$, not only in a localized window.
\end{corollary}

\begin{proof}
The proof of Proposition~\ref{prop:scale-cascade} applies verbatim to every
$t>u$, with $\mathcal R_K$ supported on $y\le t/(K+1)$. By \eqref{eq:scale-cascade}, one of the $K$ terms
$\Delta_2,\dots,\Delta_K,\mathcal R_K$ has modulus at least $|M(t)|/K$. If it is
$\Delta_k=M(\lfloor t/k\rfloor)-M(\lfloor u/k\rfloor)$, one of its two values
has modulus at least $|M(t)|/(2K)$. Since $k\le K\le u$, both candidate
arguments are positive integers; hence $y\ge1$. Moreover
$y\le\lfloor t/k\rfloor\le\lfloor t/2\rfloor$, and therefore
\begin{equation}
 \frac{|M(y)|}{y^\vartheta}
 \ge \frac{|M(t)|}{2K}\left(\frac{t}{k}\right)^{-\vartheta}
 =\frac{k^\vartheta}{2K}\frac{|M(t)|}{t^\vartheta}.
\end{equation}
\end{proof}

\begin{remark}[Channels do not iterate to a bound]
\label{rem:channel-loss}
Corollary~\ref{cor:channel-descent} has no iteration gap: it applies at
the new point $y$ relative to the last zero of $M$ before $y$. Its defect
is the constant loss $k^\vartheta/(2K)<1$ per step. If alternative~(i) held at
every step, a chain of length $m\asymp\log t$ would give only
$|M(t)|\le(2K)^m\,O(1)=t^{O(\log K)}$, weaker than the trivial bound.
For fixed $K$ and $0<\vartheta<1$ the guaranteed factor can never reach
$1$, since
\begin{equation}
 \frac{k^\vartheta}{2K}\le\frac{K^\vartheta}{2K}
 =\frac1{2K^{1-\vartheta}}<\frac12.
\end{equation}
Thus no iteration of the same pigeonhole estimate can remove the loss at an
individual step.  A closing argument must instead recover the loss
collectively from the persistence supplied by
Lemma~\ref{lem:channel-persistence}, or prove cancellation in the deep
remainder.  Both require information beyond dependency alone.
\end{remark}



\begin{proposition}[Quantified persistence along channel descent]
\label{prop:quantified-channel-persistence}
\tagN\ Fix $K\ge2$, $0<\vartheta<1$ and $\delta>0$, and put
\begin{equation}
 c_{\vartheta,K}:=\log\!\left(\frac{2K}{2^\vartheta}\right)>0.
\end{equation}
Suppose $t_0\to\infty$ through points satisfying
\begin{equation}
 Z_\vartheta(t_0):=\frac{|M(t_0)|}{t_0^\vartheta}\ge t_0^\delta.
\end{equation}
Since $|M(t_0)|\le t_0$, any nonempty asymptotic regime of this form has
$\delta\le1-\vartheta$.

At stage $j$, let $u_j$ be the last zero of $M$ before $t_j$. Continue the
descent only while $u_j\ge K$ and alternative~\textup{(i)} of
Corollary~\ref{cor:channel-descent} occurs; stop if either the deep-remainder
alternative occurs or $u_j<K$. Whenever a descent step is available there
are $2\le k_j\le K$ and $t_{j+1}\le\lfloor t_j/2\rfloor$ such that
\begin{equation}
 Z_\vartheta(t_{j+1})\ge\frac{k_j^\vartheta}{2K}Z_\vartheta(t_j).
\end{equation}
Consequently, after $m$ successful steps,
\begin{equation}
\label{eq:channel-product-persistence}
 Z_\vartheta(t_m)
 \ge Z_\vartheta(t_0)\prod_{j<m}\frac{k_j^\vartheta}{2K}
 \ge t_0^\delta e^{-m c_{\vartheta,K}}.
\end{equation}
Set
\begin{equation}
 m_*:=\left\lfloor\frac{\delta\log t_0}{2c_{\vartheta,K}}\right\rfloor,
 \qquad
 \eta_{\vartheta,K,\delta}:=
 \frac{\delta\log2}{2c_{\vartheta,K}}.
\end{equation}
Then, for all sufficiently large $t_0$, either a deep-remainder alternative
occurs before step $m_*$, or the descent is valid through step $m_*$ and
there is a descendant satisfying
\begin{equation}
 t_{m_*}\le t_0^{1-\eta_{\vartheta,K,\delta}+o(1)},
 \qquad
 Z_\vartheta(t_{m_*})\ge t_0^{\delta/2}.
\end{equation}
\end{proposition}

\begin{proof}
Each successful application of Corollary~\ref{cor:channel-descent} gives
$t_{j+1}\le\lfloor t_j/2\rfloor$ and multiplies the normalized value by at
least $2^\vartheta/(2K)=e^{-c_{\vartheta,K}}$, proving
\eqref{eq:channel-product-persistence}. Hence for every $m\le m_*$ reached
without a deep remainder,
\begin{equation}
 Z_\vartheta(t_m)\ge t_0^{\delta/2},
 \qquad
 t_m\le t_0/2^m.
\end{equation}
It remains to justify that the stopping condition $u_j<K$ cannot occur
before $m_*$. Since $M$ has infinitely many zeros, choose once and for all
a zero $z_K\ge K$. If the last zero before $t_j$ were $<K$, then necessarily
$t_j<z_K$; otherwise $z_K$ itself would be a preceding zero. Hence
\begin{equation}
 Z_\vartheta(t_j)\le z_K^{1-\vartheta},
\end{equation}
a constant depending only on $K$ and $\vartheta$. This contradicts
$Z_\vartheta(t_j)\ge t_0^{\delta/2}$ for large $t_0$. Therefore the only
possible premature stop is the deep-remainder alternative. If it does not
occur, the iteration reaches $m_*$. Finally
$t_{m_*}\le t_0 2^{-m_*}=t_0^{1-\eta_{\vartheta,K,\delta}+o(1)}$.
\end{proof}

\begin{remark}[What remains after quantified persistence]
\label{rem:quantified-persistence-gap}
Proposition~\ref{prop:quantified-channel-persistence} shows that a
polynomial normalized violation cannot disappear immediately under channel
descent. The pointwise loss accumulates,
and bounding the deep remainder by absolute values loses a factor of the
number of channels.  A closing argument must therefore exploit the
multi-index persistence of Lemma~\ref{lem:channel-persistence} collectively,
or establish cancellation among the channels entering $\mathcal R_K$.
\end{remark}

\begin{corollary}[RH in the localized linear window]
\label{cor:linear-window-RH}
\tagC\ RH holds if and only if, for every $\eps>0$,
\begin{equation}
\left|\sum_{d\le5u/8}\mu(d)\left(\left\lfloor\frac td\right\rfloor
-\left\lfloor\frac ud\right\rfloor\right)\right|
=\left|\sum_{2\le k\le t}\left[M\left(\frac tk\right)-M\left(\frac uk\right)\right]\right|
\ll_\eps t^{1/2+\eps}
\end{equation}
for every excursion starting at $u$ and every
\begin{equation}
u<t\le\min\left\{x_1,\frac54u\right\}.
\end{equation}
Equivalently, in intra-excursion coordinates $t=u+r$ it is enough to
control
\begin{equation}
1\le r\le\min\{\ell-1,u/4\}.
\end{equation}
\end{corollary}

\begin{proof}
Use Corollary~\ref{cor:prefix-cutoff} with $\chi=5/8$, together with
Proposition~\ref{prop:one-step}.
\end{proof}

\begin{proposition}[Kernel form of the recursion]
\label{prop:kernel}
\tagNE\ Let $M(u)=0$ and let $t$ be an integer with $u<t<2u$. For $1\le y\le\lfloor t/2\rfloor$ put
\begin{equation}
c_y(t,u)=\Bigl(\Bigl\lfloor\frac ty\Bigr\rfloor-\Bigl\lfloor\frac t{y+1}\Bigr\rfloor\Bigr)
-\Bigl(\Bigl\lfloor\frac uy\Bigr\rfloor-\Bigl\lfloor\frac u{y+1}\Bigr\rfloor\Bigr).
\end{equation}
Then:
\begin{enumerate}[label=\textup{(\roman*)}]
\item $M(t)=-\sum_{1\le y\le\lfloor t/2\rfloor}c_y(t,u)M(y)$;
\item $c_y(t,u)=\dfrac{t-u}{y(y+1)}+e_y$ with $|e_y|<2$; for $y>\sqrt t$
one has $c_y\in\{-1,0,1\}$, and $c_y\ne0$ only if $y=\lfloor t/k\rfloor$ or
$y=\lfloor u/k\rfloor$ for some $2\le k<\sqrt t$ (the value $k=1$ gives $y\in\{t,u\}>\lfloor t/2\rfloor$);
\item $c_1(t,u)\ge\frac12(t-u-1)$;
\item $M(t)=-(t-u)\,\Upsilon(\sqrt t)+\mathrm{Err}(t,u)$, where $\Upsilon$ is the
coefficient \eqref{eq:Upsilon_def} and
\begin{equation}
|\mathrm{Err}(t,u)|\le2\sum_{y\le\sqrt t}|M(y)|
+\sum_{2\le k<\sqrt t}\Bigl(\Bigl|M\Bigl(\frac tk\Bigr)\Bigr|+\Bigl|M\Bigl(\frac uk\Bigr)\Bigr|\Bigr).
\end{equation}
\end{enumerate}
\end{proposition}

\begin{proof}
(i) In \eqref{eq:descent-identity} group the indices $k\ge2$ by the value
$y=\lfloor t/k\rfloor$, respectively $y=\lfloor u/k\rfloor$. The number of
$k\ge1$ with $\lfloor t/k\rfloor=y$ is
$\lfloor t/y\rfloor-\lfloor t/(y+1)\rfloor$, and for $y\le\lfloor t/2\rfloor$ none of them
is $k=1$. The same holds for $u$, because $k=1$ gives $y=u>\lfloor t/2\rfloor$. (ii) For $x\in\{t,u\}$ the bracket
$\lfloor x/y\rfloor-\lfloor x/(y+1)\rfloor$ equals $x/(y(y+1))$ minus a
difference of two fractional parts, so it differs from $x/(y(y+1))$ by less
than $1$. For $y>\sqrt t\ge\sqrt u$ both
$t/y-t/(y+1)$ and $u/y-u/(y+1)$ are $<1$, so each bracket is $0$ or $1$, and
it is $1$ only if $y=\lfloor x/k\rfloor$ with $k\le x/y<\sqrt t$.
(iii) $c_1=(t-u)-(\lfloor t/2\rfloor-\lfloor u/2\rfloor)$ and
$\lfloor t/2\rfloor-\lfloor u/2\rfloor\le\lceil(t-u)/2\rceil$.
(iv) Insert (ii) into (i) for $y\le\sqrt t$; the smooth part is
$(t-u)\sum_{y\le\sqrt t}M(y)/(y(y+1))=(t-u)\Upsilon(\sqrt t)$. For
$y>\sqrt t$ use the last assertion of (ii).
\end{proof}

\begin{corollary}[Loss of the pointwise induction]
\label{cor:induction-loss}
\tagNE\ Let $0<\vartheta<1$ and suppose $|M(y)|\le Cy^\vartheta$ for all $y\le\lfloor t/2\rfloor$. Then,
for $u<t<2u$ with $M(u)=0$,
\begin{equation}
\label{eq:induction-loss}
\bigl|M(t)+(t-u)\Upsilon(\sqrt t)\bigr|\ll_\vartheta C\,t^{(1+\vartheta)/2},
\end{equation}
while the direct absolute-value estimate for the full kernel is, in general,
of linear size and yields at best $|M(t)|\ll_\vartheta C\,t$ from this
information alone. For $\vartheta=\frac12+\eps$ the pointwise induction
therefore loses $t^{1/2-\eps}$, and $t^{1/4-\eps/2}$ after the smooth
component $-(t-u)\Upsilon(\sqrt t)$ is extracted.
\end{corollary}

\begin{proof}
By Proposition~\ref{prop:kernel}(iv),
$|\mathrm{Err}|\le2C\sum_{y\le\sqrt t}y^\vartheta+2Ct^\vartheta\sum_{2\le k<\sqrt t}k^{-\vartheta}
\ll_\vartheta Ct^{(1+\vartheta)/2}$. The full kernel satisfies
$\sum_y|c_y|y^\vartheta\ge c_1\ge\frac12(t-u-1)$ by (iii), and
$\sum_y|c_y|y^\vartheta\ll_\vartheta t$ by (ii).
\end{proof}

The smooth component is not small for termwise reasons: by
\eqref{eq:C_tail},
\begin{equation}
\Upsilon(\sqrt t)=-\sum_{n>\sqrt t}\frac{M(n)}{n(n+1)},
\end{equation}
so its size is governed by the global identity
$\sum_{n\ge1}M(n)/(n(n+1))=0$, i.e.\ by the pole of $\zeta$ at $s=1$
(Remark~\ref{rem:why_log2}), and involves values of $M$ beyond $t/2$. Even
under RH the tail gives only $(t-u)\Upsilon(\sqrt t)\ll t^{3/4+\eps/2}$, the
same order as the bound for $\mathrm{Err}(t,u)$; the RH scale for $M(t)$ is a statement
about their sum.

\begin{remark}[Localized intra-excursion system]
\label{rem:sign-cone}
Let $I=(u,x_1]$ be an excursion of length $\ell=x_1-u$ and orientation
$\sigma_\nu$, and put
\begin{equation}
S_r:=\sigma_\nu M(u+r)>0\qquad(1\le r<\ell),\qquad S_0=S_\ell=0.
\end{equation}
For a fixed $\lambda\in(\lambda_*,2)$ set
\begin{equation}
h_\lambda(u):=\min\{\ell-1,(\lambda-1)u\}.
\end{equation}
Then the entire RH-relevant initial segment satisfies simultaneously
\begin{equation}
S_r=-\sigma_\nu\sum_{y\le\lambda u/2}c_y(u+r,u)M(y)>0,
\qquad1\le r\le h_\lambda(u),
\end{equation}
and $|S_r-S_{r-1}|\le1$.  For $\lambda=5/4$ this is the segment
$r\le\min\{\ell-1,u/4\}$ and all dependencies lie below $5u/8$.
\end{remark}

\begin{definition}[Relaxed divisor-sum window class]
\label{def:relaxed-divisor-cone}
Fix $u\ge2$ and $1\le h<u$.  For
$\mathbf a=(a_d)_{d\le u}\in\{-1,0,1\}^u$ with $a_1=1$ and
$\sum_{d\le u}a_d=0$, put
\begin{equation}
D_{\mathbf a}(n):=\sum_{\substack{d\mid n\\ d\le u}}a_d,
\qquad u<n\le u+h.
\end{equation}
We call $\mathbf a$ window-admissible with orientation $\sigma_\nu\in\{-1,1\}$ if
$D_{\mathbf a}(n)\in\{-1,0,1\}$ throughout the window and
\begin{equation}
\sigma_\nu\sum_{u<n\le u+r}D_{\mathbf a}(n)<0
\qquad(1\le r\le h).
\end{equation}
The induced oriented path is
$S_r=-\sigma_\nu\sum_{u<n\le u+r}D_{\mathbf a}(n)>0$.
For the genuine M\"obius prefix, $a_d=\mu(d)$, these are exactly the
increment and excursion conditions coming from $\mu*1=\delta_1$.
\end{definition}

\begin{proposition}[Ballistic counterexample to the relaxed window class]
\label{prop:relaxed-ballistic}
\tagNE\ The relaxed class of Definition~\ref{def:relaxed-divisor-cone} admits
ballistic paths.  Namely, take
\begin{equation}
a_1=1,\qquad a_u=-1,\qquad a_d=0\quad(1<d<u).
\end{equation}
Then for every $u<n<2u$ one has $D_{\mathbf a}(n)=1$, and with $\sigma_\nu=-1$,
\begin{equation}
S_r=r\qquad(1\le r<u).
\end{equation}
If $m_y:=\sum_{d\le y}a_d$, then $m_y=1$ for every $y<u$.  Hence for
any $0<\vartheta<1$, any fixed $1<\lambda<2$, and
$h=\lfloor(\lambda-1)u\rfloor$,
\begin{equation}
\sup_{y\le\lambda u/2}\frac{|m_y|}{y^\vartheta}=1,
\qquad
\frac{S_h}{(u+h)^\vartheta}\asymp_\lambda u^{1-\vartheta}\longrightarrow\infty.
\end{equation}
Thus a maximum principle comparing the normalized localized-window height
to the normalized history is false on this relaxed class.
\end{proposition}

\begin{proof}
For $u<n<2u$, the integer $u$ does not divide $n$.  Hence among the two
nonzero coordinates of $\mathbf a$ only $d=1$ contributes to $D_{\mathbf a}(n)$,
so $D_{\mathbf a}(n)=1$.  The remaining assertions are immediate.
\end{proof}

\begin{remark}[What the relaxation loses]
\label{rem:relaxed-loses}
The counterexample violates the internal divisor identities of the M\"obius
prefix.  Already for every prime $p<u$ the genuine identity requires
$a_1+a_p=0$, whereas the ballistic vector above has
$a_p=0$.  Consequently any successful localized argument of this type
must use some of the relations
\begin{equation}
\sum_{d\mid n}a_d=0\qquad(2\le n\le u),
\end{equation}
not merely the future-window constraints.
\end{remark}

\begin{remark}[The prefix correlation identity]
\label{rem:prefix-correlation}
The complete-excursion bridge identity must not be applied to a truncated
localized window.  For $1\le h<\ell$, put
$x_j=\sigma_\nu\mu(u+j)$,
$S_r=\sum_{j\le r}x_j$, and
$C_g^{\langle h\rangle}=\sum_{j=1}^{h-g}x_jx_{j+g}$.  Lemma~\ref{lem:triangular-corr}
applied to the prefix of length $h$ gives exactly
\begin{equation}
\sum_{r=1}^{h-1}S_r^2+\sum_{g=1}^{h-1}gC_g^{\langle h\rangle}
=S_h\sum_{r=1}^{h-1}S_r.
\end{equation}
The right-hand side vanishes only for a complete bridge.  Thus
$Q_\nu=-\sum_g gC_g^{(\nu)}$ cannot be used as an identity internal to a
truncated horizon unless the excursion actually returns there.  In
particular, passing from a large localized value to the energy of the whole
excursion mixes the retained window with the part deliberately discarded by
the horizon theorem.
\end{remark}

\begin{remark}[The one-step iteration gap]
\label{rem:iteration-gap}
Even a one-step implication sending a large value in the localized window
to a large value at some $y\le\lambda u/2$ would not automatically iterate.
The point $y$ need not lie in the localized window of the excursion that
contains $y$; it may lie beyond $\lambda u'$ relative to that excursion's
start $u'$.  Sparsity of the family of $\lambda$-long excursions does not
repair a single descent chain.  Iteration therefore requires either a
principle valid for all $t$, or a mechanism that places the new point inside
an admissible window at every step.
\end{remark}

\begin{proposition}[Termwise absolute-value contraction gives no nontrivial gain]
\label{prop:cone-obstruction}
\tagNE\ Let $(u,x_1]$ be an excursion, $1\le h<\min\{x_1-u,u\}$, and let
$w_1,\dots,w_h\ge0$, not all zero. Put
$K_y=\sum_{r\le h}w_rc_y(u+r,u)$, so that
$\sum_rw_rM(u+r)=-\sum_yK_yM(y)$. If, for some $0<\vartheta<1$, $\alpha\ge0$ and
$C_\alpha>0$,
\begin{equation}
\sum_y|K_y|y^\vartheta\le C_\alpha\sum_{r\le h}w_r(u+r)^\alpha,
\end{equation}
then
\begin{equation}
\sum_{r\le h}w_rr\le(1+2C_\alpha)\sum_{r\le h}w_r(u+r)^\alpha,
\qquad
\sum_{r\le h}w_r|M(u+r)|
\le(1+2C_\alpha)\sum_{r\le h}w_r(u+r)^\alpha.
\end{equation}
Thus the contraction hypothesis can hold only on windows where the trivial
bound $|M(u+r)|\le r$ already gives the conclusion.
\end{proposition}

\begin{proof}
Proposition~\ref{prop:kernel}(iii) gives
$c_1(u+r,u)\ge(r-1)/2\ge0$, hence
$\frac12\sum_rw_r(r-1)\le K_1\le
C_\alpha\sum_rw_r(u+r)^\alpha$. Since $M(u)=0$ and $M$ has unit steps,
$|M(u+r)|\le r$; since $(u+r)^\alpha\ge1$, $\sum w_r\le\sum w_r(u+r)^\alpha$, and both inequalities follow.
\end{proof}

\begin{remark}[No strict coercivity from the first coefficient]
\label{rem:c1-parity}
Writing $t=u+r$ gives the exact formula
\begin{equation}
c_1(u+r,u)=
\begin{cases}
\lceil r/2\rceil,&u\ \text{even},\\
\lfloor r/2\rfloor,&u\ \text{odd}.
\end{cases}
\end{equation}
In particular $c_1(u+1,u)=0$ when $u$ is odd.  The preceding proposition
therefore rules out only this particular termwise absolute-value mechanism;
it does not rule out a global signed contraction or cancellation among the
coefficients $K_y$.
\end{remark}

\begin{definition}[Multiplicative layers of the prefix]
\label{def:multiplicative-layers}
Let $\Omega(n)$ denote the number of prime factors of $n$, counted with
multiplicity.  For $P\ge2$ and $k\ge1$, a vector
$(a_d)_{d\le P}$ satisfies the first $k$ M\"obius layers if
$a_1=1$ and
\begin{equation}
\sum_{d\mid n}a_d=0
\qquad(2\le n\le P,\ \Omega(n)\le k).
\end{equation}
Thus layer $1$ imposes the identities at primes, layer $2$ adds prime
squares and products of two primes, and so on.
\end{definition}

\begin{proposition}[One-layer propagation into the localized future]
\label{prop:layer-propagation}
\tagNE\ Suppose $(a_d)_{d\le P}$ satisfies the first $k$ M\"obius layers. Then
\begin{equation}
a_d=\mu(d)\qquad(d\le P,\ \Omega(d)\le k).
\end{equation}
Moreover, for every integer $n$ with
\begin{equation}
P<n\le2P,\qquad\Omega(n)\le k+1,
\end{equation}
one has
\begin{equation}
-\sum_{\substack{d\mid n\\d\le P}}a_d=\mu(n).
\end{equation}
Thus the first $k$ multiplicative layers of the prefix determine the
future increments at every integer $n$ in the stated window
$P<n\le2P$ with $\Omega(n)\le k+1$.
\end{proposition}

\begin{proof}
If $d\mid n$ and $d<n$, then $n/d\ge2$, hence $\Omega(d)\le\Omega(n)-1$.
We induct on $\Omega(n)$. For $\Omega(n)=0$, $n=1$ and $a_1=1=\mu(1)$. If
$2\le n\le P$, $\Omega(n)\le k$, and $a_d=\mu(d)$ for every $d\le P$ with
$\Omega(d)<\Omega(n)$, then every proper divisor $d$ of $n$ is such a $d$, and
the layer equation gives
$a_n=-\sum_{d\mid n,\,d<n}a_d=-\sum_{d\mid n,\,d<n}\mu(d)=\mu(n)$.  For the second, every
proper divisor of $n\le2P$ is at most $n/2\le P$ and has
$\Omega(d)\le k$.  Hence the displayed sum equals
$-\sum_{d\mid n,\ d<n}\mu(d)=\mu(n)$.
\end{proof}

\begin{remark}[Multiplicative-layer depth as a quantitative invariant]
\label{q:layer-depth}
Fix $\chi\in(\chi_*,1)$ and the localized window $u<n\le2\chi u$.
Let $k=k(u)$ be a depth up to which the internal M\"obius divisor identities
are imposed below $P=\chi u$. What growth of $k(u)$, if any, suffices to
combine those identities with $M(u)=0$ and the unit-step and excursion-sign
constraints so as to force
$\max_{u<t\le2\chi u}|M(t)|\ll_\eps u^{1/2+\eps}$? A bounded depth would
give a compressed arithmetic mechanism; imposing all layers reconstructs the
M\"obius prefix. Proposition~\ref{prop:layer-propagation} makes the depth a
quantifiable intermediate invariant; the results below treat only fixed
$k$.
\end{remark}
\begin{proposition}[Layer sums in the localized window]
\label{prop:layer-sums}
\tagL\ Fix $1<\lambda<2$ and an integer $j\ge1$ before letting $u\to\infty$. Then
\begin{equation}
\label{eq:layer-sum}
\sum_{\substack{u<n\le\lambda u\\ \Omega(n)=j}}\mu(n)
=(-1)^j\,(1+o(1))\,\frac{(\lambda-1)\,u\,(\log\log u)^{j-1}}{(j-1)!\,\log u}.
\end{equation}
\end{proposition}

\begin{proof}
If $\Omega(n)=j$ and $\mu(n)\ne0$, then $n$ is squarefree with exactly $j$
prime factors and $\mu(n)=(-1)^j$. By Landau's theorem (see
\cite[Ch.~7]{MontgomeryVaughan2007}), the number $\pi_j^\flat(x)$ of such
$n\le x$ satisfies $\pi_j^\flat(x)\sim f_j(x):=x(\log\log x)^{j-1}/((j-1)!\log x)$.
For fixed $j$ and fixed $\lambda>1$,
\begin{equation}
\frac{f_j(\lambda u)}{f_j(u)}
=\lambda\,
\frac{(\log\log(\lambda u))^{j-1}}{(\log\log u)^{j-1}}
\frac{\log u}{\log(\lambda u)}\longrightarrow\lambda.
\end{equation}
Hence $f_j(\lambda u)-f_j(u)=(\lambda-1+o(1))f_j(u)$, as claimed.
\end{proof}

\begin{corollary}[Bounded layer depth]
\label{cor:bounded-depth}
\tagN\ Fix $k\ge1$, $\chi\in(\chi_*,1)$, $P=\chi u$, and the window
$u<n\le2\chi u$.
\begin{enumerate}[label=\textup{(\roman*)}]
\item The increments fixed by the first $k$ layers
(Proposition~\ref{prop:layer-propagation}), namely those with
$\Omega(n)\le k+1$, form a set of $o(u)$ integers of the window.
\item Their summed contribution over the window is
\begin{equation}
\sum_{\substack{u<n\le2\chi u\\ \Omega(n)\le k+1}}\mu(n)
=(-1)^{k+1}(1+o(1))\frac{(2\chi-1)\,u\,(\log\log u)^{k}}{k!\,\log u},
\end{equation}
so, for every fixed $\eps>0$ and all sufficiently large $u$, it exceeds $u^{1-\eps}$ in modulus and has the fixed sign $(-1)^{k+1}$.
\item Let $a_d=\mu(d)$ for $\Omega(d)\le k$ and $a_d=0$ otherwise.
For squarefree $n$ in the window with $\omega(n)=w\ge k+2$,
\begin{equation}
D_{\mathbf a}(n)=\sum_{\substack{d\mid n\\ d\le P}}a_d
=\sum_{j=0}^{k}(-1)^j\binom wj=(-1)^k\binom{w-1}{k},
\end{equation}
which has modulus at least $k+1\ge2$. Thus the zero extension of a
$k$-layer prefix violates the unit-step constraint at every such $n$, i.e.\
at all but $o(u)$ squarefree integers of the window.
\end{enumerate}
\end{corollary}

\begin{proof}
(ii) follows from Proposition~\ref{prop:layer-sums} summed over
$j\le k+1$, the term $j=k+1$ dominating. For (i), Landau's asymptotic
$\#\{n\le x:\Omega(n)=j\}\sim x(\log\log x)^{j-1}/((j-1)!\log x)$ for
fixed $j$ \cite[Ch.~7]{MontgomeryVaughan2007} (all integers, not only
squarefree ones) shows that the window contains $o(u)$ integers with
$\Omega(n)\le k+1$. For (iii), all divisors of $n$
other than $n$ are $\le n/2\le P$, and $n$ itself has $\Omega(n)=w>k$, so
$D_{\mathbf a}(n)$ is the truncated alternating binomial sum, evaluated by
induction on $k$ via $\binom{w-1}{k}=\binom wk-\binom{w-1}{k-1}$. Finally, by
Landau's theorem for squarefree integers with exactly $j$ prime factors
\cite[Ch.~7]{MontgomeryVaughan2007},
$\#\{n\le x:\mu(n)^2=1,\ \omega(n)\le k+1\}=o(x)$, whereas
$\#\{n\le x:\mu(n)^2=1\}\sim6x/\pi^2$; hence all but $o(u)$ squarefree $n$ in
the window have $\omega(n)\ge k+2$.
\end{proof}

\begin{remark}[What fixed layer depth proves]
\label{rem:depth-content}
For every fixed $k$, Corollary~\ref{cor:bounded-depth} shows that the first
$k$ layers determine only a density-zero subset of the localized window,
while the signed sum of the determined increments has modulus
$(1+o(1))(2\chi-1)u(\log\log u)^k/(k!\log u)$, larger than $u^{1-\eps}$ for every fixed $\eps>0$, once $u$ is sufficiently large. Hence, if the full sum on the window is
$O(u^{1/2+\eps})$, substantial cancellation must occur between the
increments determined by any fixed set of initial layers and the remaining
ones. The results proved here do not determine how the required depth grows
with $u$; in particular they do not show that $\asymp\log\log u$ layers
(the typical size of $\Omega$) are necessary or sufficient, which would
require estimates uniform in $k=k(u)$. Part~(iii) shows that the unit-step
constraint is far from vacuous on a $k$-layer relaxation: it forbids the
naive completion at almost every integer of the window, so the free
coordinates $a_d$ with $\Omega(d)>k$ are strongly constrained.
\end{remark}



\subsection{The square range as an exact bilinear form of the prefix}

\begin{proposition}[Bilinear square-range identity]
\label{prop:prefix-bilinear}
\tagL\ Let $u\ge1$ and $1\le t\le u^2$. Then
\begin{equation}
\label{eq:bilinear-general}
M(t)=2M(\min(t,u))-
\sum_{d,j\le u}\mu(d)\mu(j)\left\lfloor\frac{t}{dj}\right\rfloor.
\end{equation}
In particular, if $M(u)=0$ and $u<t\le u^2$,
\begin{equation}
\label{eq:bilinear-excursion}
M(t)=-\mathcal B_u(t),\qquad
\mathcal B_u(t):=\sum_{d,j\le u}\mu(d)\mu(j)
\left\lfloor\frac{t}{dj}\right\rfloor,
\qquad \mathcal B_u(u)=0.
\end{equation}
\end{proposition}

\begin{proof}
Let $\mu_u=\mu\mathbf1_{[1,u]}$. Each factor of $(\mu-\mu_u)*(\mu-\mu_u)$ is
supported on integers $\ge u+1$, so the convolution, and its convolution
with $1$, vanish on $n<(u+1)^2$, in particular on $n\le u^2$. Expanding and using $\mu*1=\delta$ gives
$\mu(n)=2\mu_u(n)-(\mu_u*\mu_u*1)(n)$ for $n\le u^2$; sum over $n\le t$.
\end{proof}

\begin{remark}[Higher-order determinacy]
\label{rem:prefix-k}
For every $k\ge1$ and $n\le u^k$,
\begin{equation}
\label{eq:HB-k}
\mu(n)=\sum_{j=1}^{k}(-1)^{j-1}\binom{k}{j}
\bigl(\mu_u^{*j}*1^{*(j-1)}\bigr)(n).
\end{equation}
Indeed, $(\mu-\mu_u)^{*k}$ is supported on integers at least
$(u+1)^k$, so after convolution with $1^{*(k-1)}$ it still vanishes for
$n<(u+1)^k$, in particular for $n\le u^k$.  Expanding gives
\begin{equation}
0=\sum_{j=0}^{k}(-1)^j\binom{k}{j}
\mu^{*(k-j)}*\mu_u^{*j}*1^{*(k-1)}.
\end{equation}
For $j=0$, $\mu^{*k}*1^{*(k-1)}=\mu$.  For $j\ge1$, pair
$k-j$ factors of $\mu$ with $k-j$ of the factors $1$ to obtain
$\mu^{*(k-j)}*1^{*(k-j)}=\delta$; the remaining $j-1$ factors of $1$
produce the term in \eqref{eq:HB-k}.  Rearranging proves the formula.
Thus determination by the prefix
is not special to $u^2$. What is special to $t\le u^2$ is the order-two form
\eqref{eq:bilinear-excursion}, with rank-one coefficients $\mu(d)\mu(j)$ and
elementary kernel $\lfloor t/(dj)\rfloor$. This is a truncation identity of
Heath-Brown/Vaughan type \cite{HB82}; summed forms of this kind are used in
the fast computation of $M$ \cite{HT23}.
\end{remark}

\begin{remark}[Square-excursion bridges]
\label{q:excursion-constraint}
Let $u=x_{0,\nu}$ and $x_{1,\nu}\le u^2$. Then
$\mathcal B_u(u)=\mathcal B_u(x_{1,\nu})=0$ and
$-\sigma_\nu\mathcal B_u(t)>0$ for $u<t<x_{1,\nu}$. Can the rank-one
coefficients $\mu(d)\mu(j)$, the kernel $\lfloor t/(dj)\rfloor$ and these
bridge and sign constraints be combined to prove the bound
\begin{equation}
\max_{u<t\le x_{1,\nu}}|\mathcal B_u(t)|
\ll_\eps x_{1,\nu}^{1/2+\eps}
\end{equation}
uniformly over square-ordinary excursions? By
Corollary~\ref{cor:power-ordinary-RH} this uniform bound is equivalent to
RH. By Corollary~\ref{cor:linear-window-RH} it suffices to consider
$t\le\min\{x_{1,\nu},\frac54u\}$, where $\mathcal B_u(t)$ reduces to
the localized linear form and uses only $\mu(d)$ with $d\le5u/8$.
\end{remark}


\subsection{Trial-division horizon of an excursion}
\label{sec:prime_horizon}

\begin{lemma}[Square horizon]
\label{lem:square-horizon}
\tagL\ Let $u\ge2$. Trial division by the primes $p\le u$ alone determines $\mu(n)$
for every $n\le u^2$.
\end{lemma}

\begin{proof}
Remove the prime factors $p\le u$ of $n$; if one occurs twice, $\mu(n)=0$.
The cofactor $q\le u^2$ has all prime factors $>u$, so it is $1$ or a single
prime.
\end{proof}

\begin{remark}[Determinacy is not information]
\label{rem:information}
Primes $>u$ are not logically new information, and determination by a
finite prefix is not special to $u^2$: \eqref{eq:HB-k} propagates a prefix
to every power range, and \eqref{eq:hyperbola-M} determines $M$ from $M(1)$.
The ranges $(u,2u]$ and $(u,u^2]$ are special only in the concrete sense
that the exact summatory formula is of order one
(\eqref{eq:linear-window-M}) or two (\eqref{eq:bilinear-excursion}) in the
prefix, and that one-stage trial division by primes $\le u$ suffices on
$(u,u^2]$.
\end{remark}

\begin{definition}[Smooth and rough parts]
\label{def:smooth_rough}
For $y>1$ write $n=s_y(n)\,r_y(n)$, where every prime factor of $s_y(n)$
is $\le y$ and every prime factor of $r_y(n)$ is $>y$. Let $\Pp(n)$ be the
largest prime factor of $n$, with $\Pp(1)=1$.
\end{definition}

\begin{lemma}
\label{lem:rough_part}
\tagL\ $\mu(n)=\mu(s_y(n))\,\mu(r_y(n))$. If $n<y^k$ then $\Omega(r_y(n))\le k-1$,
and if $n\le y^2$ then $r_y(n)$ is $1$ or a prime.
\end{lemma}

\begin{proof}
The two parts are coprime. If $\Omega(r)\ge1$ then $y^{\Omega(r)}<r\le n$.
\end{proof}

\begin{remark}[The horizon $u^2$ is not the multiplicative threshold of an excursion]
\label{rem:horizon_inactive}
For a square-ordinary excursion $\sqrt{v_\nu}\le u_\nu$. Every $n\le v_\nu$
has at most one prime factor $>\sqrt{v_\nu}$, so the threshold that governs
the multiplicative structure inside $I_\nu$ is $\sqrt{v_\nu}$, below
$u_\nu$. The regime $t>u_\nu^2$ of Section~\ref{sec:prime_copy_propagation}
lies beyond $v_\nu$ for a square-ordinary excursion; that section concerns
$M(t)$ globally, and Section~\ref{sec:local_large_prime} gives the local
version inside an excursion.
\end{remark}


\subsection{Large-prime decomposition at level \texorpdfstring{$U^2$}{U squared}}
\label{sec:prime_copy_propagation}

The identities of this section are decompositions by large prime factors of
Buchstab/Legendre type. The asymptotic statements use only the prime number
theorem with classical error term and the Vinogradov--Korobov bound
\cite{MontgomeryVaughan2007,Walfisz1963}.

\begin{definition}[Smooth baseline and large-prime part]
\label{def:baseline_channel}
For $U\ge2$ and $t\ge1$ put
\begin{equation}
\label{eq:baseline_channel_def}
  M^{\rm sm}_U(t):=\sum_{\substack{n\le t\\ \Pp(n)\le U^2}}\mu(n),\qquad
  M^{\rm lp}_U(t):=\sum_{\substack{n\le t\\ \Pp(n)>U^2}}\mu(n),\qquad
  M(t)=M^{\rm sm}_U(t)+M^{\rm lp}_U(t).
\end{equation}
\end{definition}

\begin{proposition}[Single-prime copy]
\label{prop:single_prime_copy}
\tagL\ For a prime $p$ and $t<p^2$,
\begin{equation}
\label{eq:single-prime-copy}
  \mathcal D_p(t):=\sum_{a\le t/p}\mu(ap)=-M(t/p).
\end{equation}
\end{proposition}

\begin{proof}
$a\le t/p<p$, so $p\nmid a$ and $\mu(ap)=-\mu(a)$.
\end{proof}

\begin{proposition}[Sign content of one prime cone]
\label{prop:prime-cone-signs}
\tagNE\ Let $\Sq(y)=\sum_{a\le y}\mu(a)^2$. Among the numbers $ap$ with $a\le y<p$,
the numbers of values $+1,-1,0$ of $\mu$ are
\begin{equation}
\label{eq:cone_counts}
  N_+=\frac{\Sq(y)-M(y)}{2},\qquad N_-=\frac{\Sq(y)+M(y)}{2},\qquad
  N_0=\lfloor y\rfloor-\Sq(y),
\end{equation}
with limiting proportions $3/\pi^2,\ 3/\pi^2,\ 1-6/\pi^2$.
\end{proposition}

\begin{proof}
$\mu(ap)=-\mu(a)$, and $\#\{a\le y:\mu(a)=\mp1\}=(\Sq(y)\mp M(y))/2$.
\end{proof}

\begin{proposition}[First-order window]
\label{prop:linear-new-primes}
\tagL\ If $U^2<t<U^4$, no $n\le t$ has two prime factors $>U^2$ (with
multiplicity), and
\begin{equation}
\label{eq:first-order-total}
  M^{\rm lp}_U(t)=-\sum_{U^2<p\le t}M(t/p).
\end{equation}
\end{proposition}

\begin{proof}
Two such factors would give $n>U^4>t$. Hence $M^{\rm lp}_U(t)=\sum_{U^2<p\le t}
\mathcal D_p(t)$, and $t<U^4<p^2$; apply \eqref{eq:single-prime-copy}.
\end{proof}

\begin{proposition}[Prime-count form]
\label{prop:prime_count_form}
\tagNE\ For $U^2<t<U^4$ and $W:=t/U^2$,
\begin{equation}
\label{eq:prime_count_form}
  M^{\rm lp}_U(t)=-\sum_{a\le W}\mu(a)\bigl(\pi(t/a)-\pi(U^2)\bigr),
  \qquad
  M(t)=M^{\rm sm}_U(t)-\sum_{a\le W}\mu(a)\pi(t/a)+M(W)\pi(U^2).
\end{equation}
\end{proposition}

\begin{proof}
Expand \eqref{eq:first-order-total} as
$-\sum_{U^2<p\le t}\sum_{a\le t/p}\mu(a)$.  This is a finite sum, so
reversing the order is purely combinatorial.  For fixed $a$ the prime runs
over $U^2<p\le t/a$, which is non-empty only if $a<W$.  The displayed sum
may harmlessly be written with $a\le W$: if $W$ is an integer, the endpoint
$a=W$ contributes $\pi(U^2)-\pi(U^2)=0$.
\end{proof}

Thus the only input from primes $>U^2$ is the value of $\pi$ at the points
$t/a$, $a\le W$, weighted by $\mu(a)$ at the small scale $W<U^2$.

\begin{definition}[Copy multiplicities]
\label{def:copy-multiplicity}
For $U^2<t<U^4$, $W=t/U^2$ and $1\le y\le\lfloor W\rfloor$ put
\begin{equation}
\label{eq:Ny-exact}
  N_y(t,U)=\#\Bigl\{p:\ \max\bigl(U^2,\tfrac{t}{y+1}\bigr)<p\le\tfrac ty\Bigr\},
  \qquad\text{so that}\qquad
  M^{\rm lp}_U(t)=-\sum_{1\le y\le\lfloor W\rfloor}N_y(t,U)M(y).
\end{equation}
\end{definition}

\begin{remark}[Asymptotics of the multiplicities]
\label{rem:copy-count-asymptotic}
For fixed $W>1$ and $U\to\infty$,
\begin{equation}
\label{eq:Ny-fixed-W}
  N_y\sim\frac{WU^2}{2\log U}\cdot\frac1{y(y+1)}\quad(1\le y\le\lfloor W\rfloor-1),
  \qquad
  N_{\lfloor W\rfloor}\sim\frac{U^2}{2\log U}\Bigl(\frac W{\lfloor W\rfloor}-1\Bigr)
  \quad(W\notin\mathbb{Z}),
\end{equation}
and $N_W=0$ if $W\in\mathbb{Z}$. The top multiplicity is truncated because
$t/(\lfloor W\rfloor+1)<U^2$.
\end{remark}

\begin{definition}[Coherent coefficient]
\label{def:coherent_coefficient}
\begin{equation}
\label{eq:C_def}
  \Cc(W):=\int_1^W\frac{M(w)}{w^2}\,dw
  =\sum_{a\le W}\mu(a)\Bigl(\frac1a-\frac1W\Bigr)
  =\sum_{a\le W}\frac{\mu(a)}a-\frac{M(W)}W .
\end{equation}
\end{definition}

\begin{proposition}[Fixed $W$]
\label{prop:fixed-scale}
\tagNE\ For fixed $W>1$, as $U\to\infty$,
\begin{equation}
\label{eq:E_fixed_W}
  M^{\rm lp}_U(WU^2)=-\frac{WU^2}{2\log U}\bigl(\Cc(W)+o(1)\bigr).
\end{equation}
\end{proposition}

\begin{proof}
In \eqref{eq:prime_count_form} the sum is finite; for $a<W$ the prime number
theorem gives $\pi(WU^2/a)-\pi(U^2)\sim(W/a-1)U^2/(2\log U)$, and
$\sum_{a\le W}\mu(a)(W/a-1)=W\,\Cc(W)$.
\end{proof}

\begin{lemma}[The multiplicity-weighted coefficient]
\label{lem:copy-weight-coefficient}
\tagNE\ The coefficient
\begin{equation}
\label{eq:Upsilon_def}
  \Upsilon(W):=\sum_{1\le y\le W}\frac{M(y)}{y(y+1)}
  =\sum_{n\le W}\frac{\mu(n)}n-\frac{M(W)}{\lfloor W\rfloor+1}
\end{equation}
satisfies
\begin{equation}
\label{eq:Upsilon_C}
  \Upsilon(W)=\Cc\bigl(\lfloor W\rfloor+1\bigr),\qquad
  \Cc(W)=\Upsilon(W)-M(W)\Bigl(\frac1W-\frac1{\lfloor W\rfloor+1}\Bigr),
\end{equation}
and
\begin{equation}
\label{eq:C_tail}
  \Cc(W)=-\int_W^\infty\frac{M(w)}{w^2}\,dw,\qquad
  \Upsilon(W)=-\sum_{n>W}\frac{M(n)}{n(n+1)}.
\end{equation}
\end{lemma}

\begin{proof}
$M$ is constant on $[y,y+1)$ and $\int_y^{y+1}w^{-2}dw=1/(y(y+1))$, which
gives \eqref{eq:Upsilon_C}. For $s>1$, \eqref{eq:mellin_M} gives
$\int_1^\infty M(w)w^{-s-1}\,dw=1/(s\zeta(s))\to0$ as $s\downarrow1$. For
$s\in[1,1+\delta]$ the bound $M(w)\ll we^{-c\sqrt{\log w}}$ dominates the
integrand by an integrable multiple of $e^{-c\sqrt{\log w}}/w$, so dominated
convergence gives $\int_1^\infty M(w)w^{-2}dw=0$.
\end{proof}

\begin{remark}[$\Cc$ versus $\Upsilon$]
\label{rem:C-vs-Upsilon}
Summing \eqref{eq:Ny-fixed-W} against $M(y)$ suggests the coefficient
$\Upsilon(W)$, but the truncated top multiplicity replaces it by $\Cc(W)$; by
\eqref{eq:Upsilon_C} the two agree exactly when $M(\lfloor W\rfloor)=0$. For
$1<W<2$ one has $M^{\rm lp}_U=-(\pi(t)-\pi(U^2))\sim-(W-1)U^2/(2\log U)$, i.e.\
$\Cc(W)=1-1/W$, whereas $\Upsilon(W)=1/2$.
\end{remark}

\begin{corollary}[Decay of the coefficients]
\label{cor:AR-VK}
\tagC\ By \eqref{eq:C_tail} and \eqref{eq:VK_M}, there is $c_1>0$ such that, for $W>e$,
\begin{equation}
\label{eq:AR-VK}
  \Cc(W),\ \Upsilon(W)\ll\exp\bigl(-c_1(\log W)^{3/5}(\log\log W)^{-1/5}\bigr).
\end{equation}
\end{corollary}

\begin{proposition}[Asymptotic size of the smooth and large-prime components]
\label{prop:both_parts_large}
\tagN\ \begin{enumerate}[label=\textup{(\roman*)},leftmargin=*]
\item Suppose $W>1$ is fixed and $\Cc(W)\ne0$. Then
\begin{align}
\label{eq:E_fixed_W_size}
  M^{\rm lp}_U(WU^2)&\sim-\frac{WU^2}{2\log U}\Cc(W),\\
  M^{\rm sm}_U(WU^2)&\sim+\frac{WU^2}{2\log U}\Cc(W).
\end{align}
For $1<W<2$ one has $\Cc(W)=1-1/W>0$, so the hypothesis holds.
\item There is $C>0$ such that, as $t\to\infty$, uniformly for $U^2<t<U^4$
with $W=t/U^2\ge\exp(C(\log\log t)^2)$,
\begin{equation}
\label{eq:E_log2}
  M^{\rm lp}_U(t)=(1+o(1))\frac{t}{(\log t)^2},\qquad
  M^{\rm sm}_U(t)=-(1+o(1))\frac{t}{(\log t)^2}.
\end{equation}
\end{enumerate}
In particular, in both regimes $|M^{\rm lp}_U(t)|$ and $|M^{\rm sm}_U(t)|$ exceed
$t^{1/2+\eps}$ for every fixed $\eps<1/2$ and all large $t$, and only their
sum $M$ can be compared with the scale $t^{1/2+\eps}$.
\end{proposition}

\begin{proof}
We use the classical form \cite{MontgomeryVaughan2007}: for some $c>0$,
\begin{equation}
\label{eq:pnt_inputs}
  \pi(y)=\li(y)+O\bigl(ye^{-c\sqrt{\log y}}\bigr),\quad
  M(y)\ll ye^{-c\sqrt{\log y}},\quad
  \mt(y):=\sum_{a\le y}\frac{\mu(a)}a\ll e^{-c\sqrt{\log y}},
\end{equation}
and the series
$\sum_{n\ge1}\mu(n)\log n/n=-1$ converges conditionally.  This follows,
for example, from the displayed bound on $\mt$ and partial summation (and
agrees with the limit at $s=1$ of
$-\frac{d}{ds}\zeta(s)^{-1}$).  The integral below converges absolutely by
the bound on $\mt$, and
\begin{equation}
\label{eq:mtilde_integral}
  \int_1^\infty\mt(w)\frac{dw}w
  =\lim_{Y\to\infty}\Bigl(\mt(Y)\log Y-\sum_{a\le Y}\frac{\mu(a)\log a}a\Bigr)=1.
\end{equation}

(i) Proposition~\ref{prop:fixed-scale}, and $M(t)=o(t/\log t)$. For
$1<W<2$, $M=1$ on $[1,2)$.

(ii) Since $t<U^4$, $W<\sqrt t$, so $\log(t/a)\ge\frac12\log t$ for $a\le W$.
By \eqref{eq:prime_count_form},
\begin{equation}
\label{eq:E_three_terms}
  M^{\rm lp}_U(t)=-\sum_{a\le W}\mu(a)\li(t/a)
  -\sum_{a\le W}\mu(a)\bigl(\pi(t/a)-\li(t/a)\bigr)+M(W)\pi(U^2).
\end{equation}
The second term is $\ll t(1+\log W)e^{-c\sqrt{\log t/2}}=o(t/\log^2t)$. For
the third, $\tfrac14\log t<\log U<\tfrac12\log t$, so
\begin{equation}
|M(W)|\pi(U^2)\ll We^{-c\sqrt{\log W}}\frac{U^2}{\log U}
\ll e^{-c\sqrt{\log W}}\frac{t}{\log t}=o(t/\log^2t)
\end{equation}
once $c\sqrt C>1$, where $c$ is the constant in \eqref{eq:pnt_inputs};
we fix such a $C$. The same condition makes the boundary term of the partial
summation below $o(t/\log^2t)$, and the second term of
\eqref{eq:E_three_terms} is $o(t/\log^2t)$ for any $C$. For the first, with $F(w)=w\li(t/w)$,
\begin{align}
\label{eq:partial_summation_li}
  \sum_{a\le W}\mu(a)\li(t/a)
  &=\mt(W)F(W)-\int_1^W\mt(w)F'(w)\,dw,\\
  F'(w)&=\li(t/w)-\frac{t/w}{\log(t/w)}
  =\frac{t/w}{\log^2(t/w)}\Bigl(1+O\Bigl(\frac1{\log t}\Bigr)\Bigr).
\end{align}
The boundary term is $\ll e^{-c\sqrt{\log W}}t/\log t=o(t/\log^2t)$. The
integral equals
\begin{equation}
\frac t{\log^2t}\left(\int_1^W\mt(w)g_t(w)\frac{dw}w+O(1/\log t)\right),
\qquad
g_t(w)=\frac{\log^2t}{\log^2(t/w)}\in[1,4].
\end{equation}
The required uniformity in $(U,W)$ follows by an explicit
$\varepsilon$--$W$ argument.  Given $\varepsilon_1>0$, choose $W_0$ so that
\begin{equation}
 \int_{W_0}^\infty\frac{|\mt(w)|}{w}\,dw<\varepsilon_1.
\end{equation}
For all sufficiently large $t$, the lower bound on $W$ gives $W>W_0$.  On
$[1,W_0]$,
\begin{equation}
 \sup_{1\le w\le W_0}|g_t(w)-1|\ll\frac{\log W_0}{\log t}\to0,
\end{equation}
while on $[W_0,W]$ one has $1\le g_t(w)\le4$.  The omitted interval
$(W,\infty)$ is also contained in the chosen tail.  Therefore
\begin{equation}
 \int_1^W\mt(w)g_t(w)\frac{dw}{w}
 =\int_1^\infty\mt(w)\frac{dw}{w}+o(1)=1+o(1),
\end{equation}
uniformly over all admissible pairs $(U,W)$. Hence
$M^{\rm lp}_U(t)=(1+o(1))t/\log^2t$, and $M^{\rm sm}_U=M-M^{\rm lp}_U$ with $M(t)=o(t/\log^2t)$.
\end{proof}

\begin{remark}[Why $t/\log^2t$]
\label{rem:why_log2}
$\Cc(W)$ captures only the $1/\log t$ term, whose coefficient tends to zero
by Corollary~\ref{cor:AR-VK}. The $1/\log^2t$ term comes from
$\zeta'(s)/\zeta(s)^2\to-1$, that is, from the pole of $\zeta$ at $s=1$,
and does not decay in $W$.
\end{remark}

\begin{remark}[First copy transitions]
\label{rem:first-copy-transitions}
For $1<W<2$ every copy has length $1$ and $M^{\rm lp}_U(t)=-(\pi(t)-\pi(U^2))$. For
$2<W<3$ copies of length $2$ appear but $M(2)=0$. For $3<W<4$ copies of
length $3$ enter with $M(3)=-1$, of opposite net sign.
\end{remark}


\subsection{Local large-prime decomposition of an excursion}
\label{sec:local_large_prime}

\begin{proposition}[Local decomposition]
\label{prop:local_large_prime}
\tagNE\ Let $I_\nu=(u,v]$ be a complete excursion and $y=\sqrt v$. For $u<t\le v$,
\begin{equation}
\label{eq:local_split}
  M(t)=\sum_{\substack{u<n\le t\\ \Pp(n)\le y}}\mu(n)
  -\sum_{a<y}\mu(a)\,\Pi_a(t),\qquad
  \Pi_a(t):=\#\Bigl\{p:\ \max\bigl(y,\tfrac ua\bigr)<p\le\tfrac ta\Bigr\},
\end{equation}
and the bridge condition $M(v)=0$ reads
\begin{equation}
\label{eq:local_bridge}
  \sum_{\substack{u<n\le v\\ \Pp(n)\le\sqrt v}}\mu(n)
  =\sum_{a<\sqrt v}\mu(a)\,\Pi_a(v).
\end{equation}
\end{proposition}

\begin{proof}
$M(u)=0$. An $n\le v=y^2$ has at most one prime factor $p>y$; then
$n=ap$ with $a<v/y=y<p$, so $\mu(n)=-\mu(a)$, and $(a,p)\mapsto ap$ is a
bijection onto these $n$.
\end{proof}

\begin{remark}[Primes in short intervals]
\label{rem:local_meaning}
The decomposition holds for every excursion, with threshold $\sqrt{v_\nu}$.
The large-prime input consists of prime counts in the intervals
$(u/a,t/a]$, of length at most $\ell_\nu/a$ at height about $v_\nu/a$.
Huxley proved the prime number theorem uniformly in intervals of length
$h\ge x^{7/12+\eps}$ \cite{Huxley1972}; the zero-density estimates of
Guth--Maynard give intervals of length $x^{17/30+o(1)}$
\cite{GuthMaynard2026}. For M\"obius sums, Matom\"aki--Ter\"av\"ainen prove
$\sum_{x<n\le x+x^\theta}\mu(n)=o(x^\theta)$ for every $\theta>0.55$
\cite{MatomakiTeravainen2023}. The cited results do not provide, below these ranges, an asymptotic for
$\Pi_a$ uniform in $a$.
\end{remark}

\begin{lemma}[Inverse logarithms]
\label{lem:inverse_log_sum}
\tagC\ Let $C$ be as in Proposition~\ref{prop:both_parts_large}. Uniformly for
\begin{equation}
\exp\bigl(C(\log\log x)^2\bigr)\le z\le\sqrt x,
\end{equation}
\begin{equation}
\label{eq:inverse_log_sum}
  \sum_{a\le z}\frac{\mu(a)}{a\log(x/a)}=-\frac{1+o(1)}{(\log x)^2}.
\end{equation}
\end{lemma}

\begin{proof}
Partial summation gives
\begin{equation}
\frac{\mt(z)}{\log(x/z)}-
\int_1^z\frac{\mt(w)}{w\log^2(x/w)}\,dw.
\end{equation}
Since $z\le\sqrt x$, the boundary term is uniformly
$O(e^{-c\sqrt{\log z}}/\log x)$, hence $o(1/\log^2x)$. Write
$\log^{-2}(x/w)=g_x(w)\log^{-2}x$ with $1\le g_x(w)\le4$, extend
$g_x\mathbf1_{\{w\le z\}}$ by zero, and apply dominated convergence with
\eqref{eq:mtilde_integral}.
\end{proof}

\begin{heuristic}[Opposite drifts]
\label{heu:local_drift}
Let $u\ge\sqrt v\exp(C(\log\log v)^2)$. Near $t\in I_\nu$ the counts
$\Pi_a$ grow for $a<t/\sqrt v$, at rate $1/(a\log(t/a))$ if primes near
$t/a$ have their expected local density uniformly in $a$. Over $(t,t+h]$
the large-prime part then changes by about
$-h\sum_{a<t/\sqrt v}\mu(a)/(a\log(t/a))=(1+o(1))h/\log^2t$ by
Lemma~\ref{lem:inverse_log_sum}. Since $M(v)-M(u)=0$, the $\sqrt v$-smooth
part drifts at rate about $-1/\log^2v$, independently of $\sigma_\nu$. As
\begin{equation}
\frac{\ell_\nu}{\log^2v_\nu}\gg\sqrt{\ell_\nu}
\iff
\ell_\nu\gg\log^4v_\nu,
\end{equation}
for $\ell_\nu\gg(\log v_\nu)^4$ both drift contributions exceed
$\sqrt{\ell_\nu}$, and \eqref{eq:local_bridge} is a balance between two
opposite drifts: the local analogue of
Proposition~\ref{prop:both_parts_large}.
\end{heuristic}


\section{Two-point correlations and the energy of excursions}
\label{sec:chowla}

For integers $0\le u<X$ and $1\le h<X-u$ put
\begin{equation}
C^{(u,X]}_h:=\sum_{u<n\le X-h}\mu(n)\mu(n+h),\qquad
C^{[X]}_h:=C^{(0,X]}_h .
\end{equation}
For an excursion $I_\nu=(x_{0,\nu},x_{1,\nu}]$ one has
$C_h^{(\nu)}=C_h^{(x_{0,\nu},x_{1,\nu}]}$, since $\sigma_\nu^2=1$.
Chowla's conjecture asserts $C^{[X]}_h=o(X)$ for each fixed $h\ge1$. Matom\"aki, Radziwi\l\l\ and Tao \cite{MRT15} prove, for
$K=K(X)\to\infty$, $K\le X$,
\begin{equation}
\sum_{h\le K}\left|\sum_{n\le X}\mu(n)\mu(n+h)\right|=o(KX).
\end{equation}
Since replacing the upper limit $X$ by $X-h$ changes the $h$-th sum by at
most $h$, we shall use the triangular consequence
\begin{equation}
\label{eq:MRT}
\sum_{h\le K}\bigl|C^{[X]}_h\bigr|=o(KX)
\qquad(K\to\infty,\ K=o(X)).
\end{equation}
Their stronger averaged Elliott form permits translated shifts of size
$O(X)$ and covers $\mu$.

\subsection{Signed lag averages and the weighted RH criterion}

\begin{remark}[Unweighted signed averages are degenerate]
\label{rem:chowla-degenerate}
By Lemma~\ref{lem:zeroth-corr}, every excursion satisfies
$\sum_hC_h^{(\nu)}=-q_\nu/2$, with $q_\nu\le\ell_\nu$. Hence
$\ell_\nu^{-2}\sum_hC^{(\nu)}_h=O(\ell_\nu^{-1})$ for purely algebraic
reasons: a signed average over lags of the correlations inside a bridge is
small without any pseudorandomness of $\mu$. Statements of Chowla type
inside an excursion must therefore carry absolute values inside the lag
average, or impose a genuinely weighted condition.
\end{remark}

The weighted sum behaves differently. By \eqref{eq:Q_def},
$Q_\nu=-\sum_hhC_h^{(\nu)}$: the weight $h$ records at which lags the
correlation responsible for the excursion energy sits. The same identity
holds globally at every zero of $M$, and the cross terms between different
excursions cancel exactly.

\begin{proposition}[Weighted correlations at zeros of $M$]
\label{prop:weighted-global}
\tagNE\ Let $v\ge2$ be a zero of $M$. Then
\begin{equation}
\label{eq:weighted-global}
\sum_{n<v}M(n)^2=-\sum_{h<v}h\,C^{[v]}_h,
\qquad
\sum_{h<v}h\,C^{[v]}_h=\sum_{\nu:\,x_{1,\nu}\le v}\ \sum_hh\,C^{(\nu)}_h .
\end{equation}
Thus, in the $h$-weighted sum, the correlations between increments lying in
different completed excursions cancel exactly; internal zero-run increments contribute nothing because $\mu(n)=0$ there.
\end{proposition}

\begin{proof}
Apply Lemma~\ref{lem:triangular-corr} to the block $x_n=\mu(n)$,
$1\le n\le v$: then $S_t=M(t)$ and $S_v=M(v)=0$, which gives the first
identity. Because $M(v)=0$, $\sum_{n\le v}M(n)^2=\sum_{n<v}M(n)^2$.  By
\eqref{eq:quadratic-decomp-zero}, this equals $\sum_\nu Q_\nu$ over the completed excursions ending at or before $v$, and
$Q_\nu=-\sum_hhC^{(\nu)}_h$ by \eqref{eq:Q_def}.
\end{proof}

\begin{theorem}[Weighted-correlation criterion]
\label{thm:weighted-corr-RH}
\tagC\ Using the unconditional return localization (RL) of
Corollary~\ref{cor:KP-local-returns}, RH holds if and only if, for every $\eps>0$,
\begin{equation}
-\sum_{h<v}h\,C^{[v]}_h\ll_\eps v^{2+\eps}
\qquad\text{for every zero $v$ of $M$.}
\end{equation}
\end{theorem}

\begin{proof}
If RH holds, Theorem~\ref{thm:corr-criterion} and
\eqref{eq:weighted-global} give the bound. Conversely, fix a target
$\eps_0>0$ and choose $\eps>0$ so small that
$(1+\eps)(2+\eps)\le2+\eps_0$. Let $x$ be large and put $N=\lfloor x\rfloor$.  If $M(N)=0$, then
$\sum_{n\le x}M(n)^2=B_2(N)$ is already evaluated at a zero.  Otherwise,
let $u=\max\{n\le N:M(n)=0\}$ and let $v> N$ be the next zero.  Then
$(u,v]$ is the completed excursion whose open prefix contains $N$, and
Corollary~\ref{cor:KP-local-returns} gives
$v\le u^{1+\eps}\le x^{1+\eps}$ for all sufficiently large $x$.  In
either case, with $v=N$ in the zero case,
\begin{equation}
\sum_{n\le x}M(n)^2\le\sum_{n<v}M(n)^2
\ll_\eps v^{2+\eps}\le x^{2+\eps_0}.
\end{equation}
Theorem~\ref{thm:corr-criterion} now applies.
\end{proof}

\begin{proposition}[Square-root lag cutoff at zeros]
\label{prop:sqrt-lag-cutoff}
\tagC\ Let $v\to\infty$ through zeros of $M$, and let $K=K(v)\to\infty$ with
$K\le v^{1/2}$. Then
\begin{equation}
\label{eq:small-lag-weighted}
\sum_{h\le K}h\,C_h^{[v]}=o(K^2v).
\end{equation}
In particular,
\begin{equation}
\label{eq:sqrt-lag-tail}
\sum_{n<v}M(n)^2
=-\sum_{\sqrt v<h<v}h\,C_h^{[v]}+o(v^2),
\end{equation}
where the lower limit may be replaced by $\lfloor\sqrt v\rfloor$ without
changing the statement.
\end{proposition}

\begin{proof}
By \eqref{eq:MRT},
$\sum_{h\le K}|C_h^{[v]}|=o(Kv)$. Therefore
\begin{equation}
\left|\sum_{h\le K}hC_h^{[v]}\right|
\le K\sum_{h\le K}|C_h^{[v]}|=o(K^2v).
\end{equation}
Taking $K=\lfloor\sqrt v\rfloor$ and using
Proposition~\ref{prop:weighted-global} gives \eqref{eq:sqrt-lag-tail}.
\end{proof}

\begin{corollary}[RH after deleting all small lags]
\label{cor:RH-large-lags}
\tagC\ RH holds if and only if, for every $\eps>0$,
\begin{equation}
-\sum_{\sqrt v<h<v}h\,C_h^{[v]}\ll_\eps v^{2+\eps}
\qquad\text{for every sufficiently large zero $v$ of $M$.}
\end{equation}
Thus the lags $h\le\sqrt v$ are unconditionally negligible at the critical
$v^2$ scale of the weighted-correlation criterion.
\end{corollary}

\begin{theorem}[Zero-edge lag localization]
\label{thm:theta-lag-localization}
\tagN\ Let $v\to\infty$ through zeros of $M$.  Then
\begin{equation}
\label{eq:weighted-corr-exponent}
 -\sum_{h<v}hC_h^{[v]}
 =\sum_{n<v}M(n)^2
 =v^{1+2\Theta+o(1)}.
\end{equation}
Moreover, for every fixed $0\le\alpha<\Theta$,
\begin{equation}
\label{eq:theta-small-lags-negligible}
 \sum_{h\le v^\alpha}hC_h^{[v]}
 =o\!\left(\sum_{n<v}M(n)^2\right),
\end{equation}
and hence
\begin{equation}
\label{eq:theta-large-lag-mass}
 -\sum_{v^\alpha<h<v}hC_h^{[v]}
 =(1+o(1))\sum_{n<v}M(n)^2.
\end{equation}
Thus the signed weighted correlation total is forced beyond every power
$v^\alpha$ with $\alpha<\Theta$.
\end{theorem}

\begin{proof}
The first equality is Proposition~\ref{prop:weighted-global}, and the
asymptotic exponent is Corollary~\ref{cor:pintz-excursion-exponents}.
For $K=\lfloor v^\alpha\rfloor$ with $\alpha>0$, \eqref{eq:MRT} gives
\begin{equation}
 \left|\sum_{h\le K}hC_h^{[v]}\right|
 \le K\sum_{h\le K}|C_h^{[v]}|
 =o(K^2v)=o(v^{1+2\alpha}).
\end{equation}
Choose a fixed $\delta_0$ with
$0<\delta_0<2(\Theta-\alpha)$.  Along zeros $v\to\infty$,
Corollary~\ref{cor:pintz-excursion-exponents} gives
$B_2(v)\ge v^{1+2\Theta-\delta_0}$ for all sufficiently large $v$, while
the small-lag contribution is $o(v^{1+2\alpha})$.  Their ratio is therefore
$o(v^{-2(\Theta-\alpha)+\delta_0})=o(1)$.  The case
$\alpha=0$ follows trivially from the fixed finite number of lags and
$|C_h^{[v]}|\le v$.
\end{proof}

\begin{remark}[Endpoint versus the universal square-root cutoff]
\label{rem:theta-vs-sqrt-cutoff}
Proposition~\ref{prop:sqrt-lag-cutoff} is uniform without knowing $\Theta$
and reaches the endpoint $K=\sqrt v$ at the absolute $v^2$ scale.
Theorem~\ref{thm:theta-lag-localization} is stronger in relative terms but
requires a strict power gap $\alpha<\Theta$; it does not settle the endpoint
$\alpha=\Theta$.  Under RH it says that every range
$h\le v^{1/2-\delta}$ carries $o(1)$ of the total signed weighted correlation, while if RH
fails the negligible range extends correspondingly beyond the square-root
scale.
\end{remark}

\begin{remark}[Where the cancellation must occur]
\label{rem:cross-lag}
The trivial bound is $\sum_hh|C^{[v]}_h|\le\sum_hh(v-h)\le v^3/6$.
RH requires the signed weighted sum to be $v^{2+o(1)}$. Absolute values
cannot give this at the natural scale: even square-root cancellation
$|C^{[v]}_h|\ll v^{1/2+o(1)}$ in every single lag would give only
$\sum_hh|C^{[v]}_h|\ll v^{5/2+o(1)}$. Corollary~\ref{cor:RH-large-lags}
sharpens the location of the unresolved cancellation: at the $v^2$ scale,
only lags $h>\sqrt v$ remain. By Proposition~\ref{prop:weighted-global} the
signed weighted identity itself decomposes excursion by excursion.
\end{remark}

\subsection{Lag localization of the energy of an excursion}

\begin{proposition}[Energy beyond a lag cutoff]
\label{prop:lag-localization}
\tagNE\ For every excursion $I_\nu$ of length $\ell=\ell_\nu$ and every
$1\le h_0<\ell$,
\begin{equation}
\Bigl|Q_\nu+\sum_{h_0<h<\ell}h\,C^{(\nu)}_h\Bigr|
\le h_0\sum_{h\le h_0}\bigl|C^{(\nu)}_h\bigr| .
\end{equation}
Consequently, if $\sum_{h\le h_0}|C^{(\nu)}_h|\le\delta h_0\ell$, then
\begin{equation}
\sum_{h_0<h<\ell}h\,C^{(\nu)}_h
=-Q_\nu+O(\delta h_0^2\ell).
\end{equation}
\end{proposition}

\begin{proof}
Split $Q_\nu=-\sum_{h<\ell}hC^{(\nu)}_h$ at $h_0$.
\end{proof}

\begin{corollary}[Averaged Chowla on linear-scale intervals]
\label{cor:chowla-linear}
\tagC\ Let $(u_j,u_j+L_j]$ be intervals with $L_j\to\infty$ and $u_j\le BL_j$ for a
fixed $B$, and let $h_j\to\infty$ with $h_j=o(L_j)$. Then
\begin{equation}
\sum_{h\le h_j}\bigl|C^{(u_j,u_j+L_j]}_h\bigr|=o(h_jL_j).
\end{equation}
\end{corollary}

\begin{proof}
Put $X_j=L_j-h_j$. Then $X_j\sim L_j$ and, for all large $j$,
$u_j\le2B X_j$. Apply the shifted form of
\cite[Theorem~1.6]{MRT15}, whose hypotheses explicitly allow translation
parameters $b_r\le A X_j$, to the two-point M\"obius correlation. The same
paper states immediately after Theorem~1.6 that the argument proving its
averaged Chowla theorem also applies with $\lambda$ replaced by $\mu$.
Taking one relative shift fixed and averaging the other over
$1\le h\le h_j$, with the common translation $u_j=O(X_j)$ absorbed into
the allowed $b_r$, gives
\begin{equation}
 \sum_{h\le h_j}\left|
 \sum_{1\le n\le X_j}\mu(u_j+n)\mu(u_j+n+h)
 \right|=o(h_jX_j)=o(h_jL_j).
\end{equation}
For each $h\le h_j$, extending the inner sum from $X_j=L_j-h_j$ to
$L_j-h$ adds at most $h_j-h$ terms. Hence the total triangular boundary
cost is
\begin{equation}
 \sum_{h\le h_j}(h_j-h)=O(h_j^2)=o(h_jL_j),
\end{equation}
which proves the claim.
\end{proof}

\begin{corollary}[Large-lag concentration in linear-scale excursions]
\label{cor:large-lag}
\tagN\ Let $B>0$ be fixed and let $I_j$ be excursions with $\ell_j\to\infty$,
$x_{0,j}\le B\ell_j$ and heights $H_j$ with $H_j^3/\ell_j\to\infty$. For every $h_j\to\infty$ with
$h_j^2\ell_j=o(H_j^3)$,
\begin{equation}
\sum_{h_j<h<\ell_j}h\,C^{(I_j)}_h=-(1+o(1))Q_{I_j}.
\end{equation}
If in addition $\Theta>1/2$, $0<\eta<\Theta-\frac12$ is fixed, there is a
fixed $\lambda'>1$ with $x_{1,j}\le\lambda' x_{0,j}$, and these excursions carry
Kaczorowski--Pintz witnesses
$H_j\ge\frac12x_{0,j}^{\Theta-\eta}$, then for every fixed
\begin{equation}
0<\theta_{\rm lag}<\frac{3(\Theta-\eta)-1}{2}
\end{equation}
one may take $h_j=x_{0,j}^{\theta_{\rm lag}}$, and hence
\begin{equation}
\sum_{x_{0,j}^{\theta_{\rm lag}}<h<\ell_j}h\,C_h^{(I_j)}=-(1+o(1))Q_{I_j}.
\end{equation}
\end{corollary}

\begin{proof}
By Corollary~\ref{cor:chowla-linear},
$\sum_{h\le h_j}|C^{(I_j)}_h|=\delta_jh_j\ell_j$ with $\delta_j\to0$.
Proposition~\ref{prop:lag-localization} gives an error at most
$\delta_jh_j^2\ell_j=o(H_j^3)=o(Q_{I_j})$. In the witness case the hypotheses
$x_{0,j}\le B\ell_j$ and $x_{1,j}\le\lambda' x_{0,j}$ give
$\ell_j\asymp x_{0,j}$.  Hence
$H_j^3/\ell_j\gg x_{0,j}^{3\Theta-3\eta-1}$, so the displayed condition on
$\theta_{\rm lag}$ gives $h_j^2\ell_j=o(H_j^3)$.
\end{proof}

\begin{remark}[Direction of the lag statement]
\label{rem:large-lag-direction}
Corollary~\ref{cor:large-lag} says that the energy lies beyond every cutoff
$h_j$ satisfying $h_j=o(H_j^{3/2}\ell_j^{-1/2})$; it does \emph{not} say
that the energy is concentrated only at lags larger than the endpoint scale
$H_j^{3/2}\ell_j^{-1/2}$ itself.
\end{remark}

\subsection{Uniform two-point cancellation in sublinear intervals}
\label{subsec:sublinear-chowla}

The linear-scale restriction can be removed for mesoscopic lags by a
Fourier-uniform estimate valid in every interval in the range displayed
below. Theorem~1.5 of Matom\"aki--Ter\"av\"ainen
\cite{MatomakiTeravainen2023} gives the twisted M\"obius estimate with the
original logarithmic saving $A<1/3$. Matom\"aki--Shao--Tao--Ter\"av\"ainen
record explicitly, in their discussion following equation~(1.16), that
\cite[Remark~5.2]{MatomakiTeravainen2023} allows this range to be enlarged
to $0<A<1$; see \cite[discussion following (1.16)]{MSTT23}. Thus, for
every fixed $0<A<1$ and $\iota>0$,
\begin{equation}
\label{eq:MT-fourier}
\sup_{\xi\in\mathbb R}\Bigl|\sum_{x<n\le x+L}\mu(n)e(\xi n)\Bigr|
\ll_{A,\iota}\frac{L}{(\log x)^A}
\qquad(x^{3/5+\iota}\le L\le x).
\end{equation}

\begin{theorem}[Uniform two-point Ces\`aro cancellation]
\label{thm:sublinear-two-point}
\tagC\ Fix $0<A<1$ and $\iota>0$. Uniformly for all integers $u,L,R$ with
$u^{3/5+\iota}\le L\le u$ and $1\le R\le L/2$,
\begin{equation}
\label{eq:sublinear-two-point}
\frac1R\sum_{h=1}^{R}\frac{|C^{(u,u+L]}_h|}{L-h}
\ll_{A,\iota}\Bigl(\frac LR\Bigr)^{1/2}(\log u)^{-A}.
\end{equation}
In particular, if $R\ge L(\log u)^{-2A+2\kappa}$ with
$0<\kappa<A$, the left side is $O((\log u)^{-\kappa})$.
\end{theorem}

\begin{proof}
Put $a_i=\mu(u+i)$ for $1\le i\le L$ and
$A_I(\xi)=\sum_ia_ie(i\xi)$. For $h>0$ define
$C_{-h}^{(u,u+L]}:=\sum_{i=1}^{L-h}a_{i+h}a_i=C_h^{(u,u+L]}$ and put
$C_0=q=\sum_i a_i^2$.  These are the Fourier coefficients of
$|A_I(\xi)|^2$.  Parseval gives the precise identity
\begin{equation}
\int_0^1|A_I(\xi)|^4\,d\xi
=q^2+2\sum_{h=1}^{L-1}|C_h^{(u,u+L]}|^2
\le\|A_I\|_\infty^2q
\le L\|A_I\|_\infty^2.
\end{equation}
Hence Cauchy--Schwarz gives
$\sum_{h\le R}|C_h|\le\sqrt{RL}\,\|A_I\|_\infty$, and since
$L-h\ge L/2$ the left side of \eqref{eq:sublinear-two-point} is at most
$2\|A_I\|_\infty/\sqrt{RL}$. Apply \eqref{eq:MT-fourier}.
\end{proof}

\begin{corollary}[Near-log-squared mesoscopic range]
\label{cor:near-log2-lags}
\tagC\ Fix $\iota>0$ and $0<\delta<1$. Uniformly for
\begin{equation}
u^{3/5+\iota}\le L\le u,
\qquad
\frac{L}{(\log u)^{2-\delta}}\le R\le\frac L2,
\end{equation}
one has
\begin{equation}
\frac1R\sum_{h\le R}\frac{|C^{(u,u+L]}_h|}{L-h}
\ll_{\iota,\delta}(\log u)^{-\delta/4}.
\end{equation}
\end{corollary}

\begin{proof}
In Theorem~\ref{thm:sublinear-two-point} take
$A=1-\delta/4$ and $\kappa=\delta/4$.
\end{proof}

\begin{remark}[What is and is not new]
\label{rem:sublinear-status}
Theorem~\ref{thm:sublinear-two-point} is a short Parseval--Cauchy--Schwarz
corollary of the published Fourier estimate \eqref{eq:MT-fourier}; it does
not use the excursion condition. Its role here is to provide an arithmetic
input on every adaptively selected interval, something that ``almost all
intervals'' results cannot supply. The factor $(L/R)^{1/2}$ is inherent in
this direct $L^2$-to-$L^1$ Cauchy--Schwarz step using only the displayed
second-moment control of the $C_h$; additional information on the
distribution of the correlations could in principle improve it.
\end{remark}

\begin{corollary}[Sublinear large-lag estimate]
\label{cor:sublinear-large-lag}
\tagC\ Fix $\iota>0$ and $0<\delta<1$, and let $I=(u,u+L]$ be a complete
excursion with $u^{3/5+\iota}\le L\le u$.  For all sufficiently
large $u$, put
\begin{equation}
R=\left\lceil\frac{L}{(\log u)^{2-\delta}}\right\rceil.
\end{equation}
For sufficiently large $u$, $(\log u)^{2-\delta}\ge3$, and hence
$1\le R\le L/2$; therefore the preceding two-point estimate applies.
Then
\begin{equation}
\label{eq:sublinear-large-lag}
\sum_{R<h<L}hC_h^{(I)}
=-Q_I+O_{\iota,\delta}\!\left(
\frac{L^3}{(\log u)^{4-7\delta/4}}
\right).
\end{equation}
Consequently, if
$Q_I L^{-3}(\log u)^{4-7\delta/4}\to\infty$, then the left side is
$-(1+o(1))Q_I$.
\end{corollary}

\begin{proof}
Corollary~\ref{cor:near-log2-lags} gives
$\sum_{h\le R}|C_h^{(I)}|\ll RL(\log u)^{-\delta/4}$. Hence
\begin{equation}
\left|\sum_{h\le R}hC_h^{(I)}\right|
\le R\sum_{h\le R}|C_h^{(I)}|
\ll R^2L(\log u)^{-\delta/4},
\end{equation}
which is the error in \eqref{eq:sublinear-large-lag} after inserting the
chosen value of $R$ into $Q_I=-\sum_{h<L}hC_h^{(I)}$.
\end{proof}

\begin{remark}[Connection with wide peaks]
\label{rem:sublinear-wide-peaks}
The quantitative form of the short-interval estimate in
\cite[Theorem~1.1]{MatomakiTeravainen2023} gives, after the blocking argument
of Lemma~\ref{lem:meso-lipschitz},
$\omega_u\ll_\varepsilon(\log u)^{-1/3+\varepsilon}+u^{-c}$ for some $c>0$
(depending on the two fixed short-interval exponents used there). Thus, in
the regime $H/h_u\to\infty$ of Proposition~\ref{prop:wide-peaks},
\begin{equation}
Q_I\gg_\varepsilon H^3(\log u)^{1/3-\varepsilon}.
\end{equation}
Combining this with Corollary~\ref{cor:sublinear-large-lag}, a sufficient
condition for large-lag concentration is
\begin{equation}
\left(\frac HL\right)^3
(\log u)^{13/3-7\delta/4-\varepsilon}\longrightarrow\infty.
\end{equation}
For linearly localized Kaczorowski--Pintz excursions satisfying
$v\le\lambda u$ with a fixed $\lambda<2$, this machinery applies when
$\Theta>3/5$ (after choosing the loss parameters so that
$\Theta-\eta>3/5+\iota$); the remaining obstruction is that the excursion
may be much longer than its height.
\end{remark}

\subsection{Fej\'er reduction and the weaker local-variance target}

\begin{remark}[Excursion geometry alone does not give cancellation]
\label{rem:tent-chowla}
For the tent excursion $a_i=+1$ ($i\le k$), $a_i=-1$ ($k<i\le2k$), which
satisfies the bridge and one-sign constraints of a primitive excursion,
$C_h=2k-3h$ for $1\le h<k$. Hence for $R\to\infty$, $R=o(k)$,
\begin{equation}
\frac1R\sum_{h\le R}\frac{|C_h|}{2k-h}\longrightarrow1,
\qquad
\sum_{h<R}(R-h)C_h\asymp kR^2,
\end{equation}
whereas $RL\asymp kR$. Arithmetic input is indispensable.
\end{remark}

\begin{lemma}[Finite Fej\'er identity]
\label{lem:fejer}
\tagL\ For a real sequence $(a_i)_{1\le i\le L}$, extended by $0$ outside
$[1,L]$, put
\begin{equation}
q=\sum_{i=1}^La_i^2,
\qquad
C_h=\sum_{i=1}^{L-h}a_ia_{i+h},
\qquad
F_R:=\sum_{h=1}^{R-1}(R-h)C_h.
\end{equation}
Then for $1\le R\le L$,
\begin{equation}
\label{eq:fejer}
Rq+2F_R
=\sum_{m\in\mathbb Z}\left(\sum_{j=1}^{R}a_{m+j}\right)^2.
\end{equation}
\end{lemma}

\begin{proof}
Expand the square. Each diagonal term occurs $R$ times, and a pair at lag
$h$, $1\le h<R$, occurs $2(R-h)$ times.
\end{proof}

\begin{proposition}[Bridge--Fej\'er identity and interior variance]
\label{prop:bridge-fejer}
\tagNE\ Let $I=(u,u+L]$ be a complete excursion, put
$a_i=\mu(u+i)$ and
$S_t=\sum_{i=1}^ta_i=M(u+t)$ for $0\le t\le L$. Then
$S_0=S_L=0$, and for every $1\le R<L$,
\begin{equation}
\label{eq:bridge-fejer-Q}
Rq+2F_R
=2Q_I-2\sum_{t=1}^{L-1-R}S_tS_{t+R}.
\end{equation}
Moreover
\begin{equation}
\label{eq:fejer-interior-boundary}
Rq+2F_R
=E_{\rm bd}(R)
+\sum_{m=0}^{L-R}\bigl(M(u+m+R)-M(u+m)\bigr)^2,
\end{equation}
where
\begin{equation}
E_{\rm bd}(R)=\sum_{t=1}^{R-1}S_t^2
+\sum_{t=L-R+1}^{L-1}S_t^2
=O(R^3).
\end{equation}
Consequently, if $R=o(\sqrt L)$, then
\begin{equation}
\label{eq:fejer-variance-equivalence}
F_R=o(RL)
\quad\Longleftrightarrow\quad
\sum_{m=0}^{L-R}\bigl(M(u+m+R)-M(u+m)\bigr)^2
=Rq+o(RL).
\end{equation}
If in addition $L\ge u^{1/2+\eps}$ and $L\le u$, then
\begin{equation}
\label{eq:q-squarefree-local}
q=\left(\frac6{\pi^2}+o(1)\right)L,
\end{equation}
and the right side of \eqref{eq:fejer-variance-equivalence} is equivalently
\begin{equation}
\label{eq:local-variance-diagonal}
\sum_{m=0}^{L-R}\bigl(M(u+m+R)-M(u+m)\bigr)^2
=\left(\frac6{\pi^2}+o(1)\right)RL.
\end{equation}
\end{proposition}

\begin{proof}
Extend $S_t$ by $0$ outside $0\le t\le L$. Since
$\sum_{j=1}^Ra_{m+j}=S_{m+R}-S_m$, \eqref{eq:fejer} gives
\begin{equation}
Rq+2F_R=\sum_m(S_{m+R}-S_m)^2
=2\sum_mS_m^2-2\sum_mS_mS_{m+R},
\end{equation}
The natural range of the last sum is $0\le m\le L-R$, but the terms
$m=0$ and $m=L-R$ vanish because $S_0=S_L=0$, leaving $1\le t\le L-1-R$;
this is \eqref{eq:bridge-fejer-Q}. Splitting the same sum into $m<0$, $0\le m\le L-R$, and $m>L-R$
gives \eqref{eq:fejer-interior-boundary}: for $-R<m<0$ the contribution is
$S_{m+R}^2$, for $L-R<m<L$ it is $S_m^2$, and the endpoint terms vanish
because $S_0=S_L=0$.  Since $|a_i|\le1$ and the path returns
to zero, $|S_t|\le t$ and $|S_{L-t}|\le t$, whence
$E_{\rm bd}(R)\le2\sum_{t<R}t^2=O(R^3)$. Thus
$E_{\rm bd}(R)=o(RL)$ when $R=o(\sqrt L)$, proving the equivalence.
Finally the classical estimate
$\sum_{n\le x}\mu(n)^2=(6/\pi^2)x+O(\sqrt x)$ gives
$q=(6/\pi^2)L+O(\sqrt{u+L})$.  Under
$u^{1/2+\eps}\le L\le u$,
\begin{equation}
 \frac{\sqrt{u+L}}{L}\le
 \frac{\sqrt{2u}}{u^{1/2+\eps}}\ll u^{-\eps},
\end{equation}
so the error is uniform and \eqref{eq:q-squarefree-local} follows.
\end{proof}

\begin{remark}[Fej\'er density at frequency zero]
\label{rem:fejer-spectrum}
With $A_I(\xi)=\sum_{i=1}^La_ie(i\xi)$ and
\begin{equation}
K_R(\xi)=\frac1R\left|\sum_{j=1}^Re(j\xi)\right|^2,
\end{equation}
one has
\begin{equation}
q+\frac{2F_R}{R}=\int_0^1|A_I(\xi)|^2K_R(\xi)\,d\xi.
\end{equation}
Thus $F_R=o(RL)$ says that the Fej\'er average of the local spectral density
at frequency $0$, at resolution $1/R$, equals $q+o(L)$. It is not a
statement that $|A_I(\xi)|^2$ is spectrally flat on the whole circle.
\end{remark}

\begin{remark}[Three distinct targets]
\label{rem:three-chowla-targets}
The following conditions must not be conflated.

First, absolute averaged Chowla,
$\sum_{h<R}|C_h|=o(RL)$, implies only
$|F_R|\le R\sum_{h<R}|C_h|=o(R^2L)$, not the sharper
$F_R=o(RL)$ when $R\to\infty$. Conversely, $F_R=o(RL)$ allows cancellation
between the $C_h$ and does not control their absolute values.

Second, $F_R$ does not by itself control the small-lag contribution to the
excursion energy. If
\begin{equation}
V_R:=\sum_{h<R}hC_h,
\end{equation}
then exactly
\begin{equation}
\label{eq:F-V-relation}
F_R+V_R=R\sum_{h<R}C_h.
\end{equation}
Thus a Fej\'er estimate alone does not imply $V_R=o(Q_I)$.

Third, Proposition~\ref{prop:bridge-fejer} requires both endpoints of the
block to be zeros of $M$. For a block beginning at a zero but ending before
the return, additional boundary terms remain.
\end{remark}

\begin{remark}[Why the variance target may go below $3/5$]
\label{rem:variance-target}
For the Fej\'er problem one does not need the stronger absolute estimate
$\sum|C_h|$. By Proposition~\ref{prop:bridge-fejer}, on a complete excursion
and for $R=o(\sqrt L)$ it is enough to prove the diagonal local-variance law
\begin{equation}
\sum_{m=0}^{L-R}\bigl(M(u+m+R)-M(u+m)\bigr)^2
=Rq+o(RL).
\end{equation}
This is strictly weaker than absolute averaged Chowla and is therefore a
plausible target below the Fourier-uniform threshold $L\ge u^{3/5+\iota}$.
The all-interval $U^2$/Fourier estimates currently used above do not, by
themselves, supply this asymptotic for every adaptively selected block in
the range $L\ge u^{1/2+\eps}$. Results holding for almost all starting
points cannot be substituted, because excursion endpoints are selected by
$M$ itself.

No fixed-prefix formula for the increments in
\eqref{eq:local-variance-diagonal} is asserted here. In particular, the
identity $\sum_{d\le x}\mu(d)\lfloor x/d\rfloor=1$ introduces lower-scale
terms such as $M(x/2)$ when one truncates at $x/2$; omitting those terms
would give an invalid reduction.
\end{remark}



\section{Extreme values in stochastic models of excursions}
\label{sec:extreme}
\label{subsec:extreme}

The counts of Section~\ref{subsec:geometric} bound the number of excursions
in a family, not their heights. They acquire a second function when combined
with a tail hypothesis on a normalized excursion statistic: they fix the
scale of the family maximum, with no independence assumption.

Let $(\mathfrak X,\mathcal G,\Pr)$ be a probability space carrying a random
sequence of disjoint integer intervals $I_k=(x_{0,k},x_{1,k}]$ with
$x_{1,k}\le x_{0,k+1}$, and real random variables $Z_k$ (for instance a model
of the M\"obius walk with $Z_k=H_k/\sqrt{\ell_k}$). A family $\F$ is given by
a rule selecting indices from $(x_{0,k},x_{1,k})_k$, such as $\F_{\mathfrak h}$; let
$k_\F(1)<k_\F(2)<\cdots$ enumerate its members, and put
\begin{equation}
Z_{\max,\F}(X):=\max\{Z_k:k\in\F,\ x_{0,k}\le X\}\qquad(\max\emptyset=-\infty).
\end{equation}
The lemmas of Section~\ref{subsec:geometric} hold for every realization.

\begin{hypothesis}[$\mathrm{SG}_\varpi$]
\label{hyp:SG}
There are $C,\kappa_0,\varpi>0$ such that for every $j\ge1$ and $z\ge0$
\begin{equation}
\Pr\bigl(k_\F(j)\text{ exists and }Z_{k_\F(j)}>z\bigr)\le Ce^{-\kappa_0z^\varpi}.
\end{equation}
For $\varpi=2$ this is a uniform sub-Gaussian tail, denoted (SG).
\end{hypothesis}

\begin{theorem}[Extreme scale from a count]
\label{thm:ev-union}
\tagNE\ Assume $(\mathrm{SG}_\varpi)$ for $\F$ and suppose that $E_\F(X)\le \mathcal N(X)$ for
every realization, with $\mathcal N(X)\ge1$ deterministic. Then
\begin{equation}
\label{eq:ev-union}
\Pr\bigl(Z_{\max,\F}(X)>z\bigr)\le \mathcal N(X)\,Ce^{-\kappa_0z^\varpi}.
\end{equation}
Consequently, for every $\delta>0$,
\begin{equation}
\label{eq:ev-scale}
\Pr\Bigl(Z_{\max,\F}(X)>\Bigl(\frac{1+\delta}{\kappa_0}\log \mathcal N(X)\Bigr)^{1/\varpi}\Bigr)
\le C\,\mathcal N(X)^{-\delta},
\end{equation}
and $Z_{\max,\F}(X)=O_{\Pr}\bigl((\log \mathcal N(X))^{1/\varpi}\bigr)$ as $\mathcal N(X)\to\infty$.
\end{theorem}

\begin{proof}
The members of $\F$ starting at most $X$ are the first $E_\F(X)\le \mathcal N(X)$
members of the enumeration, so
$\{Z_{\max,\F}(X)>z\}\subset\bigcup_{1\le j\le\lfloor\mathcal N(X)\rfloor}\{k_\F(j)\text{ exists},\
Z_{k_\F(j)}>z\}$, and \eqref{eq:ev-union} is the union bound. Substituting
$\kappa_0z^\varpi=(1+\delta)\log \mathcal N(X)$ gives \eqref{eq:ev-scale}.
\end{proof}

\begin{corollary}[Triple-logarithmic scale on the square family]
\label{cor:ev-square}
\tagNE\ Assume (SG) for the square-exceptional family $\F_{x^2}$ of the model, and
write $Z_{\max,\square}:=Z_{\max,\F_{x^2}}$. Then, for $X\ge e^e$ and every
$\delta>0$,
\begin{equation}
\label{eq:ev-square}
\Pr\bigl(Z_{\max,\square}(X)>z\bigr)\le4C(\log\log X)e^{-\kappa_0z^2},
\end{equation}
\begin{equation}
\label{eq:ev-square-b}
\Pr\Bigl(Z_{\max,\square}(X)>\sqrt{\tfrac{1+\delta}{\kappa_0}\log\log\log X}\Bigr)
\le4C(\log\log X)^{-\delta}.
\end{equation}
Thus $Z_{\max,\square}(X)=O_{\Pr}(\sqrt{\log\log\log X})$. The same holds for
the family $\F_{cx^2}$ of Lemma~\ref{lem:partition} with a constant depending
on $c$.
\end{corollary}

\begin{proof}
Use \eqref{eq:square-exception-count} in \eqref{eq:ev-union} and note that
$e^{-\kappa_0z^2}=(\log\log X)^{-(1+\delta)}$ at the stated $z$.
\end{proof}

The union bound gives convergence in probability; a Borel--Cantelli argument
along a sequence adapted to $\mathcal N$ gives an almost sure bound with the same
constant.

\begin{theorem}[Almost sure extreme scale]
\label{thm:ev-as}
\tagNE\ Assume $(\mathrm{SG}_\varpi)$ for $\F$ and $E_\F(X)\le \mathcal N(X)$ for every realization,
where $\mathcal N$ is continuous, nondecreasing and unbounded. Then almost surely
\begin{equation}
\limsup_{X\to\infty}\frac{Z_{\max,\F}(X)}{(\log \mathcal N(X))^{1/\varpi}}\le \kappa_0^{-1/\varpi}.
\end{equation}
\end{theorem}

\begin{proof}
Fix $\delta>0$. For $j\ge j_0$ let $X_j:=\inf\{X:\mathcal N(X)\ge2^j\}$.
By continuity, monotonicity and unboundedness, after increasing $j_0$ if
necessary one has $\mathcal N(X_j)=2^j$. Put $z_j=\bigl((1+\delta)(j+1)\log2/\kappa_0\bigr)^{1/\varpi}$ and
$A_j=\{Z_{\max,\F}(X_{j+1})>z_j\}$. By \eqref{eq:ev-union},
$\Pr(A_j)\le2^{j+1}Ce^{-(1+\delta)(j+1)\log2}=C2^{-(j+1)\delta}$, which is
summable, so almost surely $A_j$ occurs only finitely often. For
$X_j\le X\le X_{j+1}$, monotonicity of $Z_{\max,\F}$ in $X$ gives
$Z_{\max,\F}(X)\le z_j$ outside $A_j$, while $\log \mathcal N(X)\ge j\log2$. Hence
the $\limsup$ is at most
$\lim_j\bigl((1+\delta)(j+1)/(\kappa_0j)\bigr)^{1/\varpi}=((1+\delta)/\kappa_0)^{1/\varpi}$. Let
$\delta\downarrow0$ along a countable sequence.
\end{proof}

With $\mathcal N_\square(X)=2+\log_2(\log X/\log2)$, which is continuous and
increasing, Theorem~\ref{thm:ev-as} gives under (SG)
\begin{equation}
\limsup_{X\to\infty}\frac{Z_{\max,\square}(X)}{\sqrt{\log\log\log X}}\le \kappa_0^{-1/2}
\quad\text{almost surely.}
\end{equation}
Applied to the three counts available for any sequence of disjoint
intervals (all excursions: at most $X+1$; Lemma~\ref{lem:linear-lacunary};
Lemma~\ref{lem:power-lacunary}), Theorems~\ref{thm:ev-union}
and~\ref{thm:ev-as} give three levels:
\begin{equation}
\begin{array}{c|c|c}
\text{family} & \text{members starting}\le X & \text{extreme scale under (SG)}\\
\hline
\text{all excursions} & \le X+1 & \sqrt{\log X}\\
x_1>\lambda x_0 & O(\log X) & \sqrt{\log\log X}\\
x_1>cx_0^2\ \ (\text{or }x_0^A) & O(\log\log X) & \sqrt{\log\log\log X}
\end{array}
\end{equation}
Each level of lacunarity removes one logarithm from the extreme scale.
Under $(\mathrm{SG}_\varpi)$ the last column reads $(\log X)^{1/\varpi}$,
$(\log\log X)^{1/\varpi}$, $(\log\log\log X)^{1/\varpi}$.

\begin{remark}[Scope]
\label{rem:ev-scope}
For $M$ itself the square-exceptional family is finite by
Corollary~\ref{cor:KP-local-returns}, so
Theorems~\ref{thm:ev-union}--\ref{thm:ev-as} concern stochastic models of
the M\"obius walk, in which the pathwise counts of
Section~\ref{subsec:geometric} persist while no analogue of the
Kaczorowski--Pintz localization is available. Hypothesis (SG) is an
assumption on the model; the theorems state exactly what the counts
contribute to its extremes.
\end{remark}


\section{A one-frequency model for the net area curve}
\label{sec:model}

\begin{heuristic}[Scale and period]
\label{heu:model}
Motivated by \eqref{eq:explicitM}, model $M(x)\approx\sqrt x\,F(\log x)$
with $F$ almost periodic, and retain only the lowest zero $\gamma_1$ (RH
does not imply that its coefficient dominates). Within an excursion of
length $\asymp x$ this gives $H_\nu\asymp\sqrt x$ and
$\sqrt{H_\nu}\ell_\nu\asymp x^{5/4}$, so $\Anet$ oscillates at scale
\begin{equation}
\label{eq:model_scale}
  |\Anet(x)|\asymp x^{5/4},
\end{equation}
the scale of Corollary~\ref{cor:rh}, with period $2\pi/\gamma_1$ in
$\log x$, i.e.\ a factor $e^{2\pi/\gamma_1}\approx1.5597$ in $x$. In this
model consecutive zeros of $M$ are separated by the factor
$e^{\pi/\gamma_1}=\lambda_*$, the sharp threshold of this counting mechanism
in Corollary~\ref{cor:descending}(iii): a model excursion is exactly
$\lambda_*$-long.
\end{heuristic}


\section{Numerical data}
\label{sec:data}

The following census covers all $n\le N$, where $N=455\,052\,511$; it
contains $87\,385$ excursions and records $(u_\nu,\tau_\nu,v_\nu,H_\nu,\sigma_\nu)$ and the
curves $\Ap,\Am,\Anet$:
\begin{equation}
\label{eq:census_values}
  \Ap(N)=6.3449\times10^9,\quad \Am(N)=-7.0877\times10^9,\quad
  \Anet(N)=-7.4279\times10^8,
\end{equation}
\begin{equation}
\label{eq:census_max}
  \max_\nu\frac{H_\nu}{\sqrt{\ell_\nu}}=2.0743,\qquad
  \max_\nu\frac{\mathfrak b_\nu}{\ell_\nu^{3/2}}=1.0214 .
\end{equation}
Since $|M(x)|\le0.571\sqrt x$ for $33\le x\le10^{16}$ \cite{Hurst18}, and
$|M(n)|\le\sqrt n$ for $n<33$ with equality only at $n=1$, the
constant $C_0=1$ is admissible in Lemma~\ref{lem:transfer}, which gives
$\Ap(N)+|\Am(N)|\le\frac45N^{5/4}=5.317\times10^{10}$ against the observed
$1.3433\times10^{10}$, a ratio of $3.96$. Heuristic~\ref{heu:renewal} with
$L=\sqrt N$ predicts $\#\{\ell_\nu>\sqrt N\}\asymp N^{1/4}\approx146$. At
this range $N^{0.05}\approx2.7$, so factors $x^\eps$ are not separable from
constants. Reproducing these values requires the excursion file, the
area-integration convention of Proposition~\ref{prop:explicit}, and the
terminal convention of Definition~\ref{def:Ms-terminal}.

\section{Conclusions}
\label{sec:conclusions}

\subsection*{What is new and what is reformulation}
Results and formulations proved in this paper and believed to be new: the excursion skeleton and its exact geometry; the counting lemmas for
disjoint intervals; the horizon-localization theorem with threshold
$\lambda_*=e^{\pi/\gamma_1}$ and the descending horizons; the localization
of returns $x_1\le x_0^{1+o(1)}$; the exact area exponent $1+\Theta/2$, which
settles the equivalence of the $x^{5/4+\eps}$ area bound with RH; the
two-sided localization of large values and the mean logarithmic length
$\le\pi/\gamma_1$ of witness excursions; wide peaks; the prefix-dependency
cutoff $\chi_*=\lambda_*/2$; channel descent with quantified persistence;
the weighted correlation identity at zeros with $\Theta$-edge lag
localization; the large-prime decomposition $M=M^{\rm sm}_U+M^{\rm lp}_U$
with both parts of size $t/\log^2t$. Reformulations: the $L^1$, $L^2$,
energy and correlation forms of Littlewood's criterion, the exact exponents
of $B_1,B_2,\En$ from Pintz, the excursion-block series of $1/\zeta$, and the
transport of averaged Chowla, Fourier uniformity and short-interval
cancellation to excursion blocks.

\subsection*{Lines this paper can help to pursue}
\begin{enumerate}[label=\textup{(\arabic*)}]
\item \emph{Lowering $\lambda_*$.} By Corollary~\ref{cor:KP-mean-log-length}
the witness excursions have mean logarithmic length $\le\pi/\gamma_1$; any
positional information giving a fixed surplus for the witnesses beyond
$\lambda x_0$ lowers the horizon below $\lambda_*$
(Remark~\ref{rem:surplus}).
\item \emph{Local variance below $3/5$.} On complete excursions with
$L\ge u^{1/2+\eps}$ the diagonal law
$\sum_m(M(u+m+R)-M(u+m))^2=(6/\pi^2+o(1))RL$ for some $R\to\infty$,
$R=o(\sqrt L)$, is equivalent to $F_R=o(RL)$
(Proposition~\ref{prop:bridge-fejer}) and is strictly weaker than absolute
averaged Chowla; it is the natural target for adaptively selected blocks
below the Fourier-uniform range.
\item \emph{Absolute two-point cancellation on adaptive blocks} of length
$x_0^{1/2+\eps}\le\ell<x_0^{3/5}$, extending
Theorem~\ref{thm:sublinear-two-point}.
\item \emph{The lag endpoint $\alpha=\Theta$} in
Theorem~\ref{thm:theta-lag-localization}: under RH, relative negligibility
of $h\le v^{1/2}$ would locate the whole weighted correlation mass at lags
$h>v^{1/2+o(1)}$.
\item \emph{Collective use of persistence.} Channel descent iterates with a
constant loss; the multi-index persistence of
Lemma~\ref{lem:channel-persistence} and cancellation inside the deep
remainder $\mathcal R_K$ are the two places where the loss can be recovered.
\item \emph{Layer depth.} The depth $k(u)$ of Möbius layers needed to force
the RH scale on the localized window is a quantitative invariant; fixed $k$
is excluded (Corollary~\ref{cor:bounded-depth}), and estimates uniform in
$k\asymp\log\log u$ are the next step.
\item \emph{Exponent meters.} $B_1$, $B_2$, $\En$, $\mathcal A_{\rm abs}$
and the block series all have exponents that are explicit affine functions
of $\Theta$; they give several independent observables of $\Theta$ with
different weights on long and short excursions.
\item \emph{Boundary of the block series} at $\Re s=\Theta$, and a third
bridge relation for complete excursions involving the defect
$\ell_\nu/2-H_\nu$.
\item \emph{Stochastic models.} Tail hypotheses of type (SG) for
$H/\sqrt\ell$ in random-walk models, where the extreme-value theorems of
Section~\ref{sec:extreme} apply.
\end{enumerate}

\subsection*{Lines shown here not to work}
\begin{enumerate}[label=\textup{(\arabic*)}]
\item Unweighted signed lag averages inside a bridge are $O(\ell^{-1})$ for
algebraic reasons (Remark~\ref{rem:chowla-degenerate}); only absolute or
weighted averages carry arithmetic content.
\item Square-root cancellation in each individual lag gives only
$v^{5/2+o(1)}$ for the weighted sum, not $v^{2+o(1)}$
(Remark~\ref{rem:cross-lag}).
\item Pointwise channel descent loses a factor $k^\vartheta/(2K)<1/2$ per
step and cannot be iterated to a bound (Remark~\ref{rem:channel-loss}).
\item Termwise absolute-value contraction of the kernel gives nothing beyond
$|M(u+r)|\le r$ (Proposition~\ref{prop:cone-obstruction}).
\item The relaxed divisor-sum class that keeps only future-window
constraints admits ballistic paths (Proposition~\ref{prop:relaxed-ballistic}).
\item Any fixed number of multiplicative layers determines only $o(u)$
increments, whose sum is $\asymp u(\log\log u)^k/\log u$
(Corollary~\ref{cor:bounded-depth}).
\item Dependency reduction alone: the window is determined by the prefix up
to $\chi_*u$, but determinacy is not cancellation
(Remark~\ref{rem:information}).
\item The global Catalan maximum tail is the wrong statistic at square
horizons (Remark~\ref{rem:max-tail-square-limit}), and Catalan selector
bounds are one-sided (Remark~\ref{rem:catalan-selector-gap}).
\item Short-interval cancellation $|M(x+h)-M(x)|=o(h)$ and energy bounds do
not force returns and cannot lower $\lambda_*$ (Remark~\ref{rem:surplus}).
\item Power and fractional-square horizons are inactive for $M$
(Corollary~\ref{cor:KP-local-returns}); only the horizons $x_0\psi(x_0)$
and $\lambda x_0$ carry information for the actual Mertens excursions.
\item The truncated bridge identity cannot be applied inside a horizon
window (Remark~\ref{rem:prefix-correlation}).
\end{enumerate}

\begin{thebibliography}{99}

\bibitem{Mirsky1949}
L.~Mirsky, \emph{Arithmetical pattern problems relating to divisibility by $r$th powers}, Proc. London Math. Soc. (2) \textbf{50} (1949), 497--508.

\bibitem{Bjorner2011}
A.~Bj\"orner, \emph{A cell complex in number theory}, Adv. in Appl. Math.
\textbf{46} (2011), 71--85.

\bibitem{Chowla1965}
S.~Chowla, \emph{The Riemann Hypothesis and Hilbert's Tenth Problem}, Gordon
and Breach, New York, 1965.

\bibitem{dlVP1896}
C.-J. de la Vall\'ee Poussin, \emph{Recherches analytiques sur la th\'eorie
des nombres premiers}, Ann. Soc. Sci. Bruxelles \textbf{20} (1896), 183--256.

\bibitem{DelegliseRivat1996}
M.~Del\'eglise and J.~Rivat, \emph{Computing the summation of the M\"obius
function}, Experiment. Math. \textbf{5} (1996), 291--295.

\bibitem{Denjoy1931}
A.~Denjoy, \emph{L'hypoth\`ese de Riemann sur la distribution des z\'eros de
$\zeta(s)$, reli\'ee \`a la th\'eorie des probabilit\'es}, C. R. Acad. Sci.
Paris \textbf{192} (1931), 656--658.

\bibitem{Ford2002}
K.~Ford, \emph{Vinogradov's integral and bounds for the Riemann zeta
function}, Proc. London Math. Soc. (3) \textbf{85} (2002), 565--633.

\bibitem{Hadamard1896}
J.~Hadamard, \emph{Sur la distribution des z\'eros de la fonction $\zeta(s)$
et ses cons\'equences arithm\'etiques}, Bull. Soc. Math. France \textbf{24}
(1896), 199--220.

\bibitem{Hardy1914}
G.~H. Hardy, \emph{Sur les z\'eros de la fonction $\zeta(s)$ de Riemann},
C. R. Acad. Sci. Paris \textbf{158} (1914), 1012--1014.

\bibitem{Harper2020}
A.~J. Harper, \emph{Moments of random multiplicative functions, I: low
moments, better than squareroot cancellation, and critical multiplicative
chaos}, Forum Math. Pi \textbf{8} (2020), e1.

\bibitem{Haselgrove1958}
C.~B. Haselgrove, \emph{A disproof of a conjecture of P\'olya}, Mathematika
\textbf{5} (1958), 141--145.

\bibitem{Ingham1942}
A.~E. Ingham, \emph{On two conjectures in the theory of numbers}, Amer. J.
Math. \textbf{64} (1942), 313--319.

\bibitem{Landau1911}
E.~Landau, \emph{\"Uber die \"Aquivalenz zweier Haupts\"atze der analytischen
Zahlentheorie}, Sitzungsber. Akad. Wiss. Wien \textbf{120} (1911), 973--988.

\bibitem{Littlewood1912}
J.~E. Littlewood, \emph{Quelques cons\'equences de l'hypoth\`ese que la
fonction $\zeta(s)$ de Riemann n'a pas de z\'eros dans le demi-plan
$\Re(s)>\frac12$}, C. R. Acad. Sci. Paris \textbf{154} (1912), 263--266.

\bibitem{MaierMontgomery2009}
H.~Maier and H.~L. Montgomery, \emph{The sum of the M\"obius function},
Bull. London Math. Soc. \textbf{41} (2009), 213--226.

\bibitem{MR16}
K.~Matom\"aki and M.~Radziwi\l\l, \emph{Multiplicative functions in short
intervals}, Ann. of Math. (2) \textbf{183} (2016), 1015--1056.

\bibitem{Mertens1897}
F.~Mertens, \emph{\"Uber eine zahlentheoretische Function}, Sitzungsber.
Akad. Wiss. Wien, Math.-Naturwiss. Kl. \textbf{106} (1897), Abt.~2a,
761--830.

\bibitem{Mobius1832}
A.~F. M\"obius, \emph{\"Uber eine besondere Art von Umkehrung der Reihen},
J. Reine Angew. Math. \textbf{9} (1832), 105--123.

\bibitem{Pintz1983}
J.~Pintz, \emph{Oscillatory properties of $M(x)=\sum_{n\le x}\mu(n)$, I},
Acta Arith. \textbf{42} (1982/83), 49--55.

\bibitem{Pintz1987}
J.~Pintz, \emph{An effective disproof of the Mertens conjecture}, Ast\'erisque
\textbf{147--148} (1987), 325--333.

\bibitem{PlattTrudgian2021}
D.~Platt and T.~Trudgian, \emph{The Riemann hypothesis is true up to
$3\cdot10^{12}$}, Bull. London Math. Soc. \textbf{53} (2021), 792--797.

\bibitem{Redheffer1977}
R.~M. Redheffer, \emph{Eine explizit l\"osbare Optimierungsaufgabe}, Internat.
Schriftenreihe Numer. Math. \textbf{36} (1977).

\bibitem{Riemann1859}
B.~Riemann, \emph{Ueber die Anzahl der Primzahlen unter einer gegebenen
Gr\"osse}, Monatsber. K\"onigl. Preuss. Akad. Wiss. Berlin (1859), 671--680.

\bibitem{Sarnak2010}
P.~Sarnak, \emph{Three lectures on the M\"obius function, randomness and
dynamics}, lecture notes, Institute for Advanced Study, 2010.

\bibitem{Soundararajan2009}
K.~Soundararajan, \emph{Partial sums of the M\"obius function}, J. Reine
Angew. Math. \textbf{631} (2009), 141--152.

\bibitem{Stieltjes1905}
C.~Hermite and T.~J. Stieltjes, \emph{Correspondance d'Hermite et de
Stieltjes}, vol.~1, Gauthier-Villars, Paris, 1905.

\bibitem{Tao2016}
T.~Tao, \emph{The logarithmically averaged Chowla and Elliott conjectures for
two-point correlations}, Forum Math. Pi \textbf{4} (2016), e8.

\bibitem{Tao2017}
T.~Tao, \emph{Equivalence of the logarithmically averaged Chowla and Sarnak
conjectures}, in: Number Theory -- Diophantine Problems, Uniform
Distribution and Applications, Springer, 2017, 391--421.

\bibitem{TaoTeravainen2018}
T.~Tao and J.~Ter\"av\"ainen, \emph{Odd order cases of the logarithmically
averaged Chowla conjecture}, J. Th\'eor. Nombres Bordeaux \textbf{30} (2018),
997--1015.

\bibitem{TaoTeravainen2019}
T.~Tao and J.~Ter\"av\"ainen, \emph{The structure of logarithmically averaged
correlations of multiplicative functions, with applications to the Chowla
and Elliott conjectures}, Duke Math. J. \textbf{168} (2019), 1977--2027.

\bibitem{vonMangoldt1897}
H.~von Mangoldt, \emph{Beweis der Gleichung $\sum_{k=1}^\infty\mu(k)/k=0$},
Sitzungsber. K\"onigl. Preuss. Akad. Wiss. Berlin (1897), 835--852.

\bibitem{Wintner1944}
A.~Wintner, \emph{Random factorizations and Riemann's hypothesis}, Duke
Math. J. \textbf{11} (1944), 267--275.


\bibitem{CP17}
P.~Carmona and N.~P\'etr\'elis, \emph{Limit theorems for random walk excursion
conditioned to enclose a typical area},
\href{https://arxiv.org/abs/1709.06448}{arXiv:1709.06448}.

\bibitem{DKW15}
D.~Denisov, M.~Kolb and V.~Wachtel, \emph{Local asymptotics for the area of
random walk excursions}, J. London Math. Soc. \textbf{91} (2015), 495--513.

\bibitem{GuthMaynard2026}
L.~Guth and J.~Maynard, \emph{New large value estimates for Dirichlet
polynomials}, Ann. of Math. (2) \textbf{203} (2026), no.~2, 623--675,
\href{https://doi.org/10.4007/annals.2026.203.2.6}{doi:10.4007/annals.2026.203.2.6}.

\bibitem{HB82}
D.~R. Heath-Brown, \emph{Prime numbers in short intervals and a generalized
Vaughan identity}, Canad. J. Math. \textbf{34} (1982), 1365--1377.

\bibitem{HT23}
H.~A. Helfgott and L.~Thompson, \emph{Summing $\mu(n)$: a faster elementary
algorithm}, Res. Number Theory \textbf{9} (2023), Paper No.~6.

\bibitem{Hurst18}
G.~Hurst, \emph{Computations of the Mertens function and improved bounds on
the Mertens conjecture}, Math. Comp. \textbf{87} (2018), no.~310, 1013--1028.

\bibitem{Huxley1972}
M.~N. Huxley, \emph{On the difference between consecutive primes}, Invent.
Math. \textbf{15} (1972), 164--170.

\bibitem{Janson07}
S.~Janson, \emph{Brownian excursion area, Wright's constants in graph
enumeration, and other Brownian areas}, Probab. Surveys \textbf{4} (2007),
80--145.

\bibitem{KP86}
J.~Kaczorowski and J.~Pintz, \emph{Oscillatory properties of arithmetical
functions. I}, Acta Math. Hungar. \textbf{48} (1986), no.~1--2, 173--185.
\href{https://doi.org/10.1007/BF01949062}{doi:10.1007/BF01949062};
\href{https://mathscinet.ams.org/mathscinet-getitem?mr=858395}{MR858395}.

\bibitem{KotnikTeRiele2006}
T.~Kotnik and H.~te Riele, \emph{The Mertens conjecture revisited}, in
Algorithmic Number Theory (ANTS VII), Lecture Notes in Comput. Sci.
\textbf{4076}, Springer, Berlin, 2006.

\bibitem{MRT15}
K.~Matom\"aki, M.~Radziwi\l\l\ and T.~Tao, \emph{An averaged form of Chowla's
conjecture}, Algebra Number Theory \textbf{9} (2015), no.~9, 2167--2196.

\bibitem{MSTT23}
K.~Matom\"aki, X.~Shao, T.~Tao and J.~Ter\"av\"ainen, \emph{Higher uniformity
of arithmetic functions in short intervals I. All intervals}, Forum Math.
Pi \textbf{11} (2023), Paper No.~e29.

\bibitem{MatomakiTeravainen2023}
K.~Matom\"aki and J.~Ter\"av\"ainen, \emph{On the M\"obius function in all
short intervals}, J. Eur. Math. Soc. \textbf{25} (2023), no.~4, 1207--1225.

\bibitem{MontgomeryVaughan2007}
H.~L. Montgomery and R.~C. Vaughan, \emph{Multiplicative Number Theory I:
Classical Theory}, Cambridge Stud. Adv. Math. \textbf{97}, Cambridge
University Press, 2007.

\bibitem{Ng2004}
N.~Ng, \emph{The distribution of the summatory function of the M\"obius
function}, Proc. London Math. Soc. (3) \textbf{89} (2004), no.~2, 361--389.

\bibitem{NguyenThe04}
M.~Nguy\^en Th\^e, \emph{Area and inertial moment of Dyck paths}, Combin.
Probab. Comput. \textbf{13} (2004), 697--716.

\bibitem{OdlyzkoTeRiele1985}
A.~M. Odlyzko and H.~J.~J. te Riele, \emph{Disproof of the Mertens
conjecture}, J. Reine Angew. Math. \textbf{357} (1985), 138--160.

\bibitem{Pintz2026}
J.~Pintz, \emph{Oscillation of partial sums of the M\"obius function and
zeros of Riemann's zeta function}, \href{https://arxiv.org/abs/2608.24878v2}{arXiv:2608.24878v2} [math.NT] (1 September 2026).

\bibitem{Richard05}
C.~Richard, \emph{On $q$-functional equations and excursion moments},
\href{https://arxiv.org/abs/math/0503198}{arXiv:math/0503198}.

\bibitem{Takacs91}
L.~Tak\'acs, \emph{A Bernoulli excursion and its various applications},
Adv. Appl. Probab. \textbf{23} (1991), 557--585.

\bibitem{Titchmarsh1986}
E.~C. Titchmarsh, \emph{The Theory of the Riemann Zeta-Function}, 2nd ed.,
revised by D.~R. Heath-Brown, Clarendon Press, Oxford, 1986.

\bibitem{Walfisz1963}
A.~Walfisz, \emph{Weylsche Exponentialsummen in der neueren Zahlentheorie},
VEB Deutscher Verlag der Wissenschaften, Berlin, 1963.

\end{thebibliography}

\end{document}